\documentclass[11pt]{article}

\usepackage[margin=1.1in]{geometry}

\usepackage{amssymb,mathtools,natbib}

% Anchored formal environments for causalsmith papers.
% First mandatory argument = obj_id (machine anchor joining the paper to the
% Lean crosswalk; invisible in the rendered PDF). Optional argument = name.
% Each environment also defines \label{obj:<obj_id>} so prose can use
% \cref{obj:T-1}; cleveref derives the printed kind and number from the target.
\usepackage{amsmath}
\usepackage{mathrsfs}
\usepackage{amsthm}
\usepackage{booktabs}
% Long displays may break across pages; let TeX stretch a little harder before an
% overfull \hbox so wide formulas/tables are less likely to run off the page.
\allowdisplaybreaks
\setlength{\emergencystretch}{3em}
\newtheorem{theorem}{Theorem}
\newtheorem{lemma}{Lemma}
\newtheorem{assumption}{Assumption}
\newtheorem{definition}{Definition}
\newtheorem{proposition}{Proposition}
\newtheorem{remark}{Remark}
\newtheorem{citedresult}{Cited result}
% The amsthm counter is `algorithmbox` (not `algorithm`) so a later `\usepackage{algorithm}` cannot clash.
\newcounter{algorithmbox}
\renewcommand{\thealgorithmbox}{\arabic{algorithmbox}}
\NewDocumentEnvironment{theoremv}{m o s}{\begin{theorem}\label{obj:#1}{\normalfont[\texttt{\detokenize{#1}}]\IfValueT{#2}{\ (#2)}\IfBooleanT{#3}{\verificationfootnotemark}\ }}{\end{theorem}}
\NewDocumentEnvironment{lemmav}{m o s}{\begin{lemma}\label{obj:#1}{\normalfont[\texttt{\detokenize{#1}}]\IfValueT{#2}{\ (#2)}\IfBooleanT{#3}{\verificationfootnotemark}\ }}{\end{lemma}}
\NewDocumentEnvironment{assumptionv}{m o s}{\begin{assumption}\label{obj:#1}{\normalfont[\texttt{\detokenize{#1}}]\IfValueT{#2}{\ (#2)}\IfBooleanT{#3}{\verificationfootnotemark}\ }}{\end{assumption}}
\NewDocumentEnvironment{definitionv}{m o s}{\begin{definition}\label{obj:#1}{\normalfont[\texttt{\detokenize{#1}}]\IfValueT{#2}{\ (#2)}\IfBooleanT{#3}{\verificationfootnotemark}\ }}{\end{definition}}
% A `propositionv` is a proved result tiered as a Proposition (theorem-grade, machine-verified).
\NewDocumentEnvironment{propositionv}{m o s}{\begin{proposition}\label{obj:#1}{\normalfont[\texttt{\detokenize{#1}}]\IfValueT{#2}{\ (#2)}\IfBooleanT{#3}{\verificationfootnotemark}\ }}{\end{proposition}}
% A `remarkv` holds NON-load-bearing interpretive / open-question content (e.g. a causalsmith `oeq:`
% open question). It is neither proved nor assumed; no theorem may depend on it.
\NewDocumentEnvironment{remarkv}{m o s}{\begin{remark}\label{obj:#1}{\normalfont[\texttt{\detokenize{#1}}]\IfValueT{#2}{\ (#2)}\IfBooleanT{#3}{\verificationfootnotemark}\ }}{\end{remark}}
% An `algorithmv` is a DEFINITION-kind object presented as a boxed, numbered-step procedure (an
% estimator / selection rule built in stages). Same anchor contract as `definitionv`; its body is
% ordinary LaTeX whose steps are `\begin{enumerate}` items, so tex2html needs no extra machinery.
% Presented in the CONVENTIONAL algorithm style readers expect from `algorithm`+`algorithmic`:
% a rule above, an "Algorithm N (Title)" caption line, a rule under the caption, the numbered
% steps, and a closing rule. It is NOT a float — the anchor contract requires the object to stay
% where the outline places it, and floats drift. The obj-id tag is suppressed here (it is
% bookkeeping, and a caption line carrying it reads as debug output in a finished paper).
\NewDocumentEnvironment{algorithmv}{m o s}%
  {\par\addvspace{\medskipamount}\noindent\rule{\linewidth}{0.8pt}\par\nobreak\vspace{2pt}%
   \refstepcounter{algorithmbox}\label{obj:#1}%
   \noindent\textbf{Algorithm \thealgorithmbox}\IfValueT{#2}{ \textnormal{(#2)}}%
   \IfBooleanT{#3}{\verificationfootnotemark}\par\nobreak\vspace{2pt}%
   \noindent\rule{\linewidth}{0.4pt}\par\nobreak\vspace{2pt}\normalfont}%
  {\par\nobreak\vspace{2pt}\noindent\rule{\linewidth}{0.8pt}\par\addvspace{\medskipamount}}
% A `citedv` is an IMPORTED external result the paper relies on but does NOT prove
% (a causalsmith `gate_class:"cited"` node). It is machine-checked against the cited
% source at F2.5, never proved here; the fixed provenance note makes that explicit
% so the environment can never be read as a contribution of this paper. The \citep
% attribution is supplied by the surrounding prose (P2), keyed to references.bib.
\NewDocumentEnvironment{citedv}{m o s}{\begin{citedresult}\label{obj:#1}{\normalfont[\texttt{\detokenize{#1}}]\IfValueT{#2}{\ (#2)}\IfBooleanT{#3}{\verificationfootnotemark}\ \textit{(imported result; machine-checked against the cited source, not proved here.)}\ }}{\end{citedresult}}
% \leanref{obj_id}{display text}: an inline first-mention link to a from-note object's
% verified Lean code. In the PDF it renders the display text plainly (the Lean link is a
% web-only affordance); tex2html turns it into a drawer-opening span.
\newcommand{\leanref}[2]{#2}
% For a theorem/lemma whose Lean declaration accepts a source-matched published proposition as an
% input, the starred anchored environment puts the mark in its heading; the generated text remains
% outside the frozen mathematical environment. Keep the combined macro for older paper bundles.
\newcommand{\verificationfootnotemark}{\footnotemark}
\newcommand{\verificationfootnotetext}[1]{\footnotetext{#1}}
\newcommand{\verificationfootnote}[1]{\verificationfootnotemark\verificationfootnotetext{#1}}
% xcolor before hyperref (hyperref would otherwise pull in plain `color` for colorlinks, and a
% later xcolor load clashes). The page-1 identity stamp at the end of this file needs a grey.
\usepackage{xcolor}
\usepackage{graphicx}
\usepackage{eso-pic}
% hyperref follows natbib and the environment macros; cleveref must follow hyperref.
\usepackage[colorlinks=true,linkcolor=blue,citecolor=blue,urlcolor=blue]{hyperref}
\usepackage[capitalise,nameinlink,noabbrev]{cleveref}
% Pin the names explicitly: labels are placed inside the underlying amsthm environments, and these
% declarations make the PDF contract independent of cleveref's theorem-name inference. House style:
% reference names always print with a capital first letter ("Theorem 1", "Definition 2"), in any
% sentence position — hence `capitalise` above and capitalized \crefname forms below.
\crefname{theorem}{Theorem}{Theorems}
\Crefname{theorem}{Theorem}{Theorems}
\crefname{lemma}{Lemma}{Lemmas}
\Crefname{lemma}{Lemma}{Lemmas}
\crefname{assumption}{Assumption}{Assumptions}
\Crefname{assumption}{Assumption}{Assumptions}
\crefname{definition}{Definition}{Definitions}
\Crefname{definition}{Definition}{Definitions}
\crefname{proposition}{Proposition}{Propositions}
\Crefname{proposition}{Proposition}{Propositions}
\crefname{remark}{Remark}{Remarks}
\Crefname{remark}{Remark}{Remarks}
\crefname{citedresult}{Cited result}{Cited results}
\Crefname{citedresult}{Cited result}{Cited results}
\crefname{algorithmbox}{Algorithm}{Algorithms}
\Crefname{algorithmbox}{Algorithm}{Algorithms}
% ---------------------------------------------------------------------------
% Page-1 identity stamp (arXiv style) and the pinned title-block date.
%
% Split of responsibility: THIS file owns the FORMAT (placement, font, colour, punctuation),
% the generated `paper_stamp.tex` owns only the VALUES (number+version, area, revised date,
% DOI, link target). P4 refreshes this template on every emit, so a layout fix reaches every
% bundle from a recompile; and because `paper_stamp.tex` is not one of the P2 assembly
% digest's authored inputs, a DOI that arrives after publication needs only that one small
% file rewritten plus a recompile — never a redraft, never a reassembly.
%
% Defaults first, so a bundle WITHOUT a stamp file compiles exactly as it always did:
% `\cspaperdate` falls back to `\today` and no stamp is drawn.
\providecommand*{\cspaperdate}{\today}  % the \date{} target
\providecommand*{\csstampid}{}          % e.g. CSWP-2026-008v3 (empty ⇒ no stamp at all)
\providecommand*{\csstamparea}{}        % e.g. Stat
\providecommand*{\csstampdate}{}        % e.g. 27 Aug 2026
\providecommand*{\csstampdoi}{}         % e.g. 10.5281/zenodo.123 (empty ⇒ DOI part omitted)
\providecommand*{\csstamplink}{}        % href target (DOI url, else the paper's page; may be empty)
\IfFileExists{paper_stamp.tex}{% GENERATED by CausalSmith P4 — paper identity stamp values. Do not hand-edit.
%
% Field order on the stamp (owned by paper_macros.tex):
%   <number>v<version> · doi:<doi> · [<Area>] <date>   — the DOI is never the last token, so a
%   text extractor cannot run it into whatever follows. Without a DOI that part is omitted.
% Format lives in paper_macros.tex; only the values are here, and this file is NOT one of
% the P2 assembly digest's authored inputs. A DOI arriving after publication therefore needs
% this file rewritten plus a `latexmk` recompile — never a redraft or a reassembly.
\renewcommand*{\csstampid}{CSWP-2026-030v3}
\renewcommand*{\csstamparea}{Stat}
\renewcommand*{\csstampdate}{1 Oct 2026}
\renewcommand*{\csstampdoi}{}
\renewcommand*{\csstamplink}{https://causalsmith.org/papers/stat_transport_cace_roughnuisance_length_v1}
\renewcommand*{\cspaperdate}{1 October 2026}
}{}
\makeatletter
% Field order is deliberate: the DOI is NEVER the last token on the line. A trailing identifier
% runs into whatever a text extractor emits next, which makes it unsafe to match mechanically
% (the Zenodo stamp matcher reads exactly this line); the date is a safe terminator.
%   <number>v<version> · doi:<doi> · [<Area>] <date>      — with a DOI
%   <number>v<version> · [<Area>] <date>                  — without
\newcommand{\cs@stampsep}{\,\textperiodcentered\,}
\newcommand{\cs@stamptext}{%
  \csstampid
  \ifx\csstampdoi\@empty\else\cs@stampsep doi:\csstampdoi\fi
  \ifx\csstamparea\@empty\else\cs@stampsep[\csstamparea]\fi
  \ifx\csstampdate\@empty\else\quad\csstampdate\fi
}
\ifx\csstampid\@empty\else
  \definecolor{csstampgrey}{gray}{0.45}
  % Background shipout material: it occupies no space, so the text block and the \thanks
  % footnote are byte-identical to an unstamped compile. The starred form applies to the
  % NEXT page shipped only — i.e. page 1. Placed in the left margin (the text block starts
  % at 1.1in), reading bottom-to-top, in a Times-like face at 8pt.
  \AddToShipoutPictureBG*{%
    \AtPageLowerLeft{%
      \put(\LenToUnit{0.42in},\LenToUnit{1.1in}){%
        \rotatebox{90}{%
          \fontfamily{ptm}\fontsize{8}{10}\selectfont
          \ifx\csstamplink\@empty
            \textcolor{csstampgrey}{\cs@stamptext}%
          \else
            \expandafter\href\expandafter{\csstamplink}{\textcolor{csstampgrey}{\cs@stamptext}}%
          \fi
        }%
      }%
    }%
  }
\fi
\makeatother


\title{Honest confidence intervals for transported complier effects with rough nuisance functions}

\author{CausalSmith\thanks{Machine-generated by the CausalSmith research pipeline.}}

\date{\cspaperdate}

\begin{document}

\maketitle

\begin{abstract}
We characterize the optimal expected length of uniformly honest confidence intervals for a target-population complier average treatment effect in an equal-size two-sample experiment. The data comprise \(n\) independent source records containing a scalar covariate, encouragement, treatment receipt and outcome, together with \(n\) independent target covariate records. Source and target densities, the encouragement propensity and conditional arm means are unknown and belong to fixed Hölder classes with exponent \(1/8\). For \(n\ge256\), under fixed density envelopes, encouragement overlap, bounded potential outcomes, instrument validity, monotonicity, and transport of conditional complier shares and complier-weighted outcome effects, the minimax expected length is of order \(\min\{1,n^{-1/3}/a\}\), uniformly over \(0<a\le1/4\), where \(a\) is the transported first-stage evaluation floor. Coverage is required over the full model with positive transported compliance. An observable split-sample cubic estimator and calibrated affine-score inversion attain this order with a single interval sequence across all floors. A matching lower bound uses a continuum of admissible instrumental-variable laws with a no-defiers potential-outcome realization and applies to arbitrary measurable confidence sets. The frontier has a constant-order branch when the strength floor is at or below the rough estimation scale and an inverse-strength branch above it.
\end{abstract}

\section{Introduction}
\label{sec:introduction}

Transporting an instrumental-variable effect requires learning both the population weights and the encouragement contrasts that identify the effect among target-population compliers. When transported compliance is small, uncertainty in those contrasts is amplified by the compliance denominator. This paper characterizes that interaction through the optimal expected length of confidence intervals with finite-sample coverage uniformly over a specified continuous-covariate model.

Our experiment combines two independent samples of equal size: source records containing covariates, encouragement, receipt and outcomes, and target records containing covariates. The covariate is scalar and lies in the unit interval. Source and target densities, the source encouragement propensity, and outcome and receipt arm means belong to fixed Hölder classes with exponent \(1/8\). Density envelopes are fixed and positive, encouragement probabilities lie between \(1/4\) and \(3/4\), and potential outcomes lie in \([0,1]\). Conditional instrument randomization, consistency, exclusion and monotonicity connect observed source contrasts to complier quantities. Two transport equalities, holding for target-almost every covariate value, carry the conditional complier share and the conditional complier-weighted outcome effect across populations. Under these restrictions, \cref{obj:lem:transported-cace-identification} identifies the target complier average treatment effect as the ratio of target-averaged source outcome and receipt contrasts.

The main result concerns intervals whose coverage is at least \(1-\alpha\) at every law in this model, where \(\alpha\in(0,1)\) is the fixed noncoverage probability. Precision is evaluated on laws whose transported first stage lies between \(a\), the evaluation floor, and \(1/4\). Thus global coverage and strength-specific expected length are distinct requirements. For fixed density envelopes and Hölder radius, \cref{obj:thm:sharp-length-frontier-full-elbow} establishes
\[
L_n(a)\asymp \min\{1,n^{-1/3}/a\},
\qquad n\ge256,\quad 0<a\le1/4,
\]
where \(L_n(a)\) is the minimax worst-case expected length among globally honest connected intervals. The comparison constants depend only on the fixed regularity parameters and noncoverage probability. The bounds are uniform across the stated strength range and contain no leading logarithmic factor.

A single observable interval sequence attains this frontier. The construction expresses the transported contrasts as smooth rational functionals of seven observable marked densities. A clipped histogram pilot, followed by linear, quadratic and cubic corrections using separate evaluation blocks, yields the uniform mean-square bound in \cref{obj:lem:marked-cubic-mse}. Inverting an affine score with a deterministic radius calibrated from that bound gives the interval in \cref{obj:def:score-interval}. Its coverage holds for every positive sample size, and its expected length is bounded by a constant times the rough estimation scale divided by the actual transported first stage, capped at constant order. The evaluation floor enters the performance criterion; the same procedure serves every strength slice.

The converse preserves the causal structure of the problem. The construction in \cref{obj:def:legal-iv-mixture,obj:lem:legal-iv-mixture-full-elbow} supplies admissible observed probability laws together with compatible potential receipts and outcomes satisfying monotonicity and the transport restrictions. Its actual compliance share is the larger of the evaluation floor and the rough estimation scale. A continuum of common effect values across sign mixtures, combined with observational proximity to a center experiment, produces the matching expected-length lower bound. \Cref{obj:thm:honest-lower-full-elbow} applies to arbitrary measurable confidence sets contained in the bounded effect domain. The constant-order branch therefore reflects the information available under global coverage when the strength floor is at or below \(n^{-1/3}\).

This characterization joins three established lines of work. Transported complier-effect identification and regular inference are developed by \citet[Theorems 3.1, 4.2 and 4.3]{ChenHuang2025Generalizing}, while \citet[Theorem 6]{Clark2024} establish regular inference for transported principal effects under their identification and nuisance-rate conditions. Rough transport-functional estimation has a direct antecedent in \citet[Theorems 6 and 8]{ZengKennedyBodnarNaimi2025Efficient}; their lower-bound exponent specializes to one third under the stated one-dimensional smoothness and design conditions. Higher-order correction follows the functional-estimation approach of \citet[Section 1]{Robins2008}. Identification-robust score inversion is developed for LATE by \citet[Theorem 3.2]{Ma2026IdentificationRobust}, and its confidence-set geometry is studied by \citet[Section 3 and Theorems 1 and 2]{SmuclerLanniMasip2025Score}. The immediate expected-length benchmark is \citet{CausalSmith2026TransportedLATEStrengthFrontier}, which treats supplied geometry and specified cell-based weight-learning settings. The present result connects rough functional estimation and score inversion to a sharp expected-length criterion with fully learned continuous geometry, equal sample sizes and uniform coverage under the declared causal model.

The analysis covers scalar covariates, equal source and target sample sizes, fixed Hölder exponent (1/8), and the causal and transport restrictions stated in \cref{sec:model}. \Cref{sec:related-work} gives the closest methodological comparisons; \cref{sec:model} specifies the observed experiment and inference criterion; and \cref{sec:main-results} states the frontier and its mechanism. \Cref{sec:discussion} interprets the calibration and formulates the unequal-sample problem. The technical construction and supporting bounds appear in \cref{sec:upper-appendix}, the legal lower family in \cref{sec:lower-appendix}, and the auxiliary comparators and verification scope in \cref{sec:verification-appendix}. Full proofs of the anchored results are collected in \cref{sec:deferred-proofs}.


\section{Related work}
\label{sec:related-work}

The closest comparisons concern transported complier effects, estimation of rough transport functionals, and inference under weak identification. Our immediate benchmark is \citet{CausalSmith2026TransportedLATEStrengthFrontier}, which characterizes asymptotically honest expected length through first-stage strength and transport-weight dispersion when the covariate and assignment geometry is supplied. Its feasible cell-based procedures learn target weights under uniform source cells and balanced assignment, or under regular nonuniform cells with supplied source probabilities and assignment propensities. Those results assume proportional source and target sample sizes and cell counts growing more slowly than the square root of the source sample size. The present comparison concerns the continuous-covariate model in \cref{obj:def:model-class}: the densities, encouragement propensity and arm means have Hölder exponent \leanref{sym:\beta}{\ensuremath{\beta}}, fixed at \(1/8\), and the two samples each contain \leanref{sym:n}{\ensuremath{n}} records.
% Source: [Honest Expected Length for Transported Complier Effects with Weak First Stages](https://causalsmith.org/papers/stat_transported_late_strength_frontier_v1/).

For identification and regular inference, \citet[Theorems 3.1, 4.2 and 4.3]{ChenHuang2025Generalizing} provide a particularly close comparison. Their observation scheme combines study records containing covariates, encouragement, receipt and outcomes with target covariates. Theorem 3.1 identifies the target complier average causal effect as a ratio of target-averaged source encouragement contrasts under their sampling overlap, target-population IV validity, and mean exchangeability restrictions for the outcome contrast and first stage. Their regular inference results distinguish supplied and estimated selection probabilities. Theorem 4.2 assumes known selection and assignment models, a fixed source-to-target sample-size ratio, and finite weighted second moments. Theorem 4.3 treats a correctly specified logistic selection model estimated by maximum likelihood, with known assignment probabilities and the stated moment and full-rank conditions. These theorems support delta-method inference at a fixed positive first stage. Our identification uses the transport restrictions on the complier-weighted outcome contrast and complier share in \cref{obj:ass:transport-complier-outcome,obj:ass:transport-complier-share}, while the inference criterion in \cref{obj:def:honest-intervals} requires coverage uniformly over the model class.
% Source: [Generalizing causal effects with noncompliance](https://arxiv.org/html/2506.00149).

A distinct observation scheme underlies \citet[Sections 5 and 7]{RudolphVanDerLaan2017Robust}. They transport outcomes to a site where covariates, encouragement and receipt are observed, so that the target first stage can be estimated directly. Their identification combines a common conditional outcome model across sites with target encouragement exchangeability, exclusion, monotonicity and positivity. Section 5 derives the transported-CACE influence curve and constructs a ratio of targeted estimators, with inference based on the influence curve or the multivariate delta method. Section 7 applies this construction to the Moving to Opportunity experiment. In the covariate-only target scheme studied here, both the outcome and receipt contrasts are obtained by averaging source contrasts over the target covariate distribution, as stated in \cref{obj:lem:transported-cace-identification}. This difference determines which transport restrictions identify the denominator as well as the numerator.
% Source: [Robust estimation of encouragement-design intervention effects transported across sites](https://pmc.ncbi.nlm.nih.gov/articles/PMC5784860/).

Transported principal-effect inference supplies another close comparison. Within their principal-stratification identification framework, \citet[Theorem 6]{Clark2024} establish asymptotic linearity, normality and efficiency for an influence-function estimator using trial assignment, receipt and outcome data together with target covariates. Their theorem assumes a positive limiting source-to-target sample-size ratio, nuisance estimation on a separate independent sample, uniform positivity of the relevant true and estimated probabilities, bounded outcome regressions, and consistency of the component influence functions. Efficient regular inference further requires each of the three specified products of nuisance estimation errors to be negligible relative to the inverse square root of sample size. These conditions describe a regular estimation regime for transported principal effects. The rough-class analysis here instead calibrates the observable estimator through the uniform mean-square bound in \cref{obj:lem:marked-cubic-mse}, which supplies the error control used by the score interval.
% Source: [Transportability of Principal Causal Effects](https://arxiv.org/html/2405.04419).

The reduced-form estimation scale has a direct methodological antecedent in \citet[Theorems 6 and 8]{ZengKennedyBodnarNaimi2025Efficient}. Their Theorem 6 gives a minimax lower bound for generalization when source and target covariates coincide, the joint covariate density for source membership and a specified treatment arm is constant, and the compound selection--treatment probability and outcome regression satisfy the stated positivity bounds. Its exponent specializes to one third for one-dimensional covariates and average nuisance smoothness one eighth. Their Theorem 8 gives conditional bias and variance bounds for a quadratic estimator of a transportation functional. That result assumes independent training data, uniform positivity, and uniformly well-conditioned true and estimated projection Gram matrices. Its bias bound includes both a product of projection approximation errors and a remainder multiplying two nuisance errors by the inverse-Gram estimation error. Thus, the exponent and the use of higher-order corrections are established antecedents, with the design, approximation and remainder conditions forming part of their results.
% Source: [Efficient Generalization and Transportation](https://arxiv.org/html/2302.00092).

The cubic construction in \cref{obj:def:cubic-estimator} belongs to the higher-order functional-estimation approach developed by \citet[Section 1]{Robins2008}. The marked densities in \cref{obj:def:marked-density-vector} express each transported contrast as the integral of a smooth rational function, as specified in \cref{obj:def:smooth-functional}. This representation permits a clipped pilot followed by observable linear, quadratic and cubic corrections computed from separate sample blocks. The construction specializes higher-order bias correction to the joint source--target observation scheme and provides the uniform mean-square control needed for finite-sample calibration.
% Source: [Higher order influence functions and minimax estimation of nonlinear functionals](https://arxiv.org/abs/0805.3040).

Score inversion connects that estimation problem to identification-robust inference. \citet[Theorem 3.2]{Ma2026IdentificationRobust} establish uniformly correct asymptotic size for an orthogonalized test of a LATE moment under their high-dimensional nuisance, sparsity, overlap, moment and empirical regularity conditions. The result accommodates weak identification within that class. Our score interval in \cref{obj:def:score-interval} uses a finite-sample calibration derived from a uniform mean-square bound over the specified rough class. The shared principle is to test an affine moment in the candidate effect; the calibration and the model over which coverage is required determine the resulting guarantee.
% Source: [Identification-robust inference for the LATE with high-dimensional covariates](https://arxiv.org/html/2302.09756).

The geometry of this inversion is clarified by \citet[Section 3 and Theorems 1 and 2]{SmuclerLanniMasip2025Score}. Their studentized score inversion reduces to a quadratic inequality and can produce bounded or unbounded intervals, disconnected sets, the whole line, an empty set or a singleton. Under their independent nuisance estimation, boundedness, overlap, consistency and product-rate conditions, Theorem 1 shows agreement with the doubly robust Wald interval up to stochastic endpoint errors of order \(O_P(n^{-1})\), with probability tending to one, at a fixed law with a nonzero first stage. Theorem 2 gives a nonnormal ratio limit under their local-to-zero moment and nuisance-negligibility conditions. In our construction, a deterministic calibration radius and the bounded causal domain yield the connected acceptance interval and fallback rules specified in \cref{obj:def:score-interval}. These comparisons locate the inversion principle within established weak-identification methodology.
% Source: [A note on the properties of the confidence set for the local average treatment effect obtained by inverting the score test](https://arxiv.org/html/2506.10449).

Two older principles organize the criterion used here. Anderson--Rubin inversion constructs confidence sets by testing an affine reduced-form moment without dividing by a possibly weak first stage \citep[pp. 46--63]{Anderson1949}; Fieller's ratio analysis likewise allows denominator uncertainty to determine confidence-set geometry \citep[pp. 175--185]{Fieller1954}. The present interval specializes that logic to the bounded causal domain and replaces an asymptotic critical value with the finite-sample rough-functional radius in \cref{obj:def:score-interval}. On the decision-theoretic side, weak identification can force unbounded confidence sets on an unrestricted parameter space \citep[pp. 1365--1387]{Dufour1997}, while honest expected length over bounded functional classes is governed by between-class modulus arguments \citep[pp. 1805--1840]{CaiLow2004Adaptation}. Here bounded outcomes cap length at constant order, and the legal transported-IV family supplies the testing comparison that yields the second branch of the frontier.

The contribution is the joint equal-size, rough-geometry, uniform-coverage and expected-length characterization in \cref{obj:thm:sharp-length-frontier-full-elbow}. Its upper bound uses an observable cubic estimator and a calibrated score interval; its lower bound is realized by the primitive IV family in \cref{obj:def:legal-iv-mixture,obj:lem:legal-iv-mixture-full-elbow}. That family retains the causal and transport restrictions of the model while producing the observed experiments needed for the length comparison. Together, these results connect the established rough-functional exponent and score-inversion principle to a sharp expected-length criterion for transported complier effects with unknown continuous covariate and assignment geometry.


\section{Model, transport restrictions and inference criterion}
\label{sec:model}

Throughout, sample sizes are positive integers, and all random variables and conditional expectations are measurable. We use \(\lVert\cdot\rVert_\infty\) for the supremum norm and \(|\cdot|\) for absolute value. Generic positive constants may depend on the fixed density envelopes, smoothness radius and coverage level; their values may change between occurrences. The notation \(u_n\lesssim v_n\) denotes an inequality up to such a constant, uniformly over the stated model class and strength range, and \(u_n\asymp v_n\) denotes inequalities in both directions.

\subsection{Primitive variables and the observed experiment}

The scalar covariate domain is \leanref{sym:\mathcal X}{\ensuremath{\mathcal X}}=[0,1]. Population membership is coded by \(S=1\) for the source and \(S=0\) for the target. A primitive full-data record under the full-data law \leanref{sym:P_n^F}{\ensuremath{P_n^F}} is
\[
O^F=(S,X,D(0),D(1),Y(0),Y(1)),
\]
where \(D(0),D(1)\in\{0,1\}\) are potential receipts under encouragement, \(Y(0),Y(1)\in[0,1]\) are potential outcomes under receipt, and \(C=\mathbf 1\{D(1)>D(0)\}\) is the complier indicator. The target conditional effect among compliers is represented by
\[
\theta=
\frac{\mathbb E_{P_n^F}[(Y(1)-Y(0))C\mid S=0]}
     {\mathbb E_{P_n^F}[C\mid S=0]},
\]
whenever the denominator is positive.

The source assignment law augments a source full-data draw with encouragement \(Z\), using propensity \leanref{sym:e}{\ensuremath{e}}, and consistency induces the observed receipt \(D=D(Z)\) and outcome \(Y=Y(D)\). Its observed record is \(O^{\mathrm S}=(X,Z,D,Y)\) with one-record law \leanref{sym:P_{\mathrm S}}{\ensuremath{P_{\mathrm S}}}. The target observation is only \(X^{\mathrm T}=X\), with one-record law \leanref{sym:P_{\mathrm T}}{\ensuremath{P_{\mathrm T}}}. Thus \(P_n^F\) determines the causal variables, the source assignment mechanism determines the observed source law, and the target-population covariate marginal determines \(P_{\mathrm T}\). The samples \leanref{sym:\mathcal S_n}{\ensuremath{\mathcal S_n}} and \leanref{sym:\mathcal T_n}{\ensuremath{\mathcal T_n}} contain \(n\) records from these two observed laws. Their densities and conditional arm means receive the regularity restrictions below.

\begin{assumptionv}{ass:source-iid}[Independent source sampling]
The observed source records \(\mathcal S_n\) are independent and identically distributed with common law \(P_{\mathrm S}\), the law of one observed source record:
\[
\mathcal S_n\sim P_{\mathrm S}^{\otimes n}.
\]
\end{assumptionv}

We first impose independent sampling within each population and independence between the two samples.

\begin{assumptionv}{ass:target-iid}[Independent target sampling]
The observed target covariate records \(\mathcal T_n\) are independent and identically distributed with common law \(P_{\mathrm T}\), the law of one observed target covariate:
\[
\mathcal T_n\sim P_{\mathrm T}^{\otimes n}.
\]
\end{assumptionv}

Independent source sampling makes the source records repeated observations from a common population law, as in standard two-sample transport designs \citep{ZengKennedyBodnarNaimi2025Efficient}. The corresponding target condition fixes the population represented by the covariate sample.

\begin{assumptionv}{ass:sample-independence}[Source–target sample independence]
The observed source records \(\mathcal S_n\) and observed target covariate records \(\mathcal T_n\) are independent:
\[
\mathcal S_n\perp\!\!\!\perp\mathcal T_n.
\]
\end{assumptionv}

Independent target sampling allows the target covariate distribution to be learned from repeated draws under the same two-sample framework \citep{ZengKennedyBodnarNaimi2025Efficient}. We complete the sampling design by requiring the two collections of records to be independent.

\begin{definitionv}{synth_1}[The scalar covariate space]
The covariate space is
\[
\mathcal X := [0,1] \subset \mathbb R.
\]
\end{definitionv}

Source–target sample independence describes separately drawn samples and gives the joint observed law a product structure \citep{ZengKennedyBodnarNaimi2025Efficient}.

\subsection{Density, overlap and smoothness restrictions}

The regularity conditions use fixed envelopes and a common Hölder convention. We define these quantities together so that density, propensity and regression restrictions share the same scale.



The lower density envelope \leanref{sym:c_f}{\ensuremath{c_f}} and upper density envelope \leanref{sym:C_f}{\ensuremath{C_f}} control covariate support across populations. The common radius \leanref{sym:L}{\ensuremath{L}} bounds both components of the Hölder ball \leanref{sym:\mathcal H^{\beta}([0,1],L)}{\ensuremath{\mathcal H^\beta([0,1],L)}}. At exponent \(\beta=1/8\), the class permits rough continuous nuisance functions. We first impose density bounds for the source and target designs.

\begin{assumptionv}{ass:source-density-bounds}[Source density and support bounds]
The source covariate density \(f_{\mathrm S}\), with \(c_f\) and \(C_f\) the fixed lower and upper density envelopes, satisfies the following conditions:
\begin{itemize}
\item \textbf{(Envelope constants.)} \(0<c_f<1<C_f\).
\item \textbf{(Full-data law and support.)} \(P_n^F\), the full-data law, is a probability law. Under this law, the covariate \(X\) and both potential outcomes \(Y(0)\) and \(Y(1)\) belong to \([0,1]\) almost surely.
\item \textbf{(Density identity.)} The covariate marginal of \(P_{\mathrm S}\), the law of one observed source record, has density \(f_{\mathrm S}\) with respect to Lebesgue measure restricted to \(\mathcal X=[0,1]\). Thus, for every Borel set \(B\subseteq\mathbb R\),
\[
P_{\mathrm S}(X\in B)=\int_{B\cap[0,1]} f_{\mathrm S}(x)\,dx.
\]
\item \textbf{(Pointwise bounds.)} For every \(x\in\mathcal X\),
\[
c_f\le f_{\mathrm S}(x)\le C_f.
\]
\end{itemize}
\end{assumptionv}

A uniformly positive bounded source design density supplies source observations throughout the covariate domain, a standard support condition for transport analysis \citep{ZengKennedyBodnarNaimi2025Efficient}. The full-data support clause also fixes the scale of the potential outcomes. The target density obeys the same envelopes.

\begin{assumptionv}{ass:target-density-bounds}[Target covariate density bounds]
The target covariate law \(P_{\mathrm T}\), with target density \(f_{\mathrm T}\), satisfies the following conditions on \(\mathcal X=[0,1]\), the covariate domain, with real constants \(c_f\) and \(C_f\), the lower and upper density envelopes:
\begin{itemize}
\item \textbf{(Density identity.)} For every Borel set \(B\subseteq\mathbb R\),
\[
P_{\mathrm T}(B)=\int_{B\cap\mathcal X}\max\{f_{\mathrm T}(x),0\}\,dx.
\]
\item \textbf{(Density envelope.)} For every \(x\in\mathcal X\),
\[
c_f\le f_{\mathrm T}(x)\le C_f.
\]
\item \textbf{(Positive lower envelope.)}
\[
0<c_f.
\]
\end{itemize}
\end{assumptionv}

Uniformly positive bounded target density gives both populations the common support used in transport comparisons \citep{ZengKennedyBodnarNaimi2025Efficient}. Together, the two density envelopes bound the target-to-source density ratio. We next control variation within that support.

\begin{assumptionv}{ass:source-density-holder}[Source density Hölder regularity]
The source covariate density \(f_{\mathrm S}\) belongs to \(\mathcal H^{\beta}([0,1],L)\), the class of continuous functions on \([0,1]\) with supremum norm and \(\beta\)-Hölder seminorm bounded separately by \(L\), where \(\beta\) is the common Hölder exponent and \(L\) is the common Hölder radius:
\[
\beta=\frac18,\qquad L>1,\qquad
f_{\mathrm S}\in\mathcal H^{\beta}([0,1],L).
\]
Thus the bounds are
\[
\lVert f_{\mathrm S}\rVert_{\infty}\le L,
\qquad
[f_{\mathrm S}]_{\beta}\le L,
\]
where \([f_{\mathrm S}]_{\beta}\) denotes the \(\beta\)-Hölder seminorm of the source density.
\end{assumptionv}

The one-dimensional isotropic Hölder ball controls the source density's local variation while allowing low regularity, the smoothness framework used in higher-order functional estimation \citep{RobinsLiMukherjeeTchetgenVanDerVaart2017HigherOrder}. The target density receives the same exponent and radius.

\begin{assumptionv}{ass:target-density-holder}[Target density Hölder regularity]
The target covariate density \(f_{\mathrm T}\) satisfies the following conditions:
\begin{itemize}
\item \textbf{(Radius.)} The common Hölder radius \(L\) satisfies \(L>1\).
\item \textbf{(Exponent.)} The common Hölder exponent is \(\beta=1/8\).
\item \textbf{(Regularity.)} \(f_{\mathrm T}\in\mathcal H^\beta([0,1],L)\): it is continuous on \([0,1]\) and satisfies
\[
\sup_{x\in[0,1]}|f_{\mathrm T}(x)|\le L,
\qquad
[f_{\mathrm T}]_\beta
=
\sup_{\substack{x,y\in[0,1]\\x\ne y}}
\frac{|f_{\mathrm T}(x)-f_{\mathrm T}(y)|}{|x-y|^\beta}
\le L.
\]
\end{itemize}
\end{assumptionv}

Target density Hölder regularity places the target design in the same one-dimensional smoothness class as the source density, giving common control over the covariate distributions entering our transported functional. The related transport rate comparison specifies smoothness of outcome regressions and inverse selection--treatment weights \citep{ZengKennedyBodnarNaimi2025Efficient}. For source encouragement, we impose positivity before its smoothness restriction.

\begin{assumptionv}{ass:propensity-overlap}[Source encouragement propensity overlap]
The source encouragement propensity \(e\) satisfies
\[
\frac14\le e(x)\le\frac34
\qquad\text{for every }x\in\mathcal X.
\]
\end{assumptionv}

Source instrument positivity ensures that both encouragement arms are represented at every covariate value \citep{ChenHuang2025Generalizing}. Its fixed bounds concern assignment probabilities, whereas compliance strength will be governed by a separate condition. The propensity also satisfies the common Hölder restriction.

\begin{assumptionv}{ass:propensity-holder}[Hölder regularity of the propensity]
The source encouragement propensity \(e\) belongs to \(\mathcal H^\beta([0,1],L)\), with Hölder exponent \(\beta=1/8\) and common Hölder radius \(L>1\); explicitly,
\[
\lVert e\rVert_\infty\le L,
\qquad
[e]_\beta\le L.
\]
\end{assumptionv}

A Hölder-smooth propensity controls how assignment probabilities vary with covariates in our model. The related transport rate comparison imposes smoothness on an inverse selection--treatment weight and an outcome regression \citep{ZengKennedyBodnarNaimi2025Efficient}. The remaining observed nuisance functions are the receipt and outcome regressions within each encouragement arm.

\begin{assumptionv}{ass:arm-means-holder}[Hölder regularity of arm means]
The source conditional arm means \(m_{Az}\), for responses \(A\in\{Y,D\}\) and encouragement arms \(z\in\{0,1\}\), satisfy the following conditions, with common Hölder radius \(L\) and Hölder exponent \(\beta\):
\begin{itemize}
\item \textbf{(Smoothness.)} \(L>1\), \(\beta=1/8\), and each \(m_{Az}\) is continuous on \([0,1]\), with
\[
\lVert m_{Az}\rVert_\infty\le L,
\qquad
[m_{Az}]_\beta\le L.
\]
\item \textbf{(Range.)} For every \(A\), \(z\), and \(x\in[0,1]\),
\[
0\le m_{Az}(x)\le 1.
\]
\item \textbf{(Source arm moments.)} Under the usual measurability conditions, for every \(A\), \(z\), and measurable set \(B\subseteq[0,1]\),
\[
\int_{\{X\in B,\;Z=z\}} A\,dP_{\mathrm S}
=
\int_B f_{\mathrm S}(x)\,
e(x)^z\bigl(1-e(x)\bigr)^{1-z}\,
m_{Az}(x)\,dx,
\]
where \(P_{\mathrm S}\) is the law of an observed source record, \(f_{\mathrm S}\) is the source covariate density, and \(e\) is the source encouragement propensity.
\end{itemize}
\end{assumptionv}

Hölder-smooth observed arm regressions control the covariate variation of the conditional receipt and outcome means \citep{ZengKennedyBodnarNaimi2025Efficient}. Their range reflects the bounded responses, and the arm-moment identities tie the chosen regression versions to the observed source law.

\subsection{Instrument validity and complier transport}

The causal restrictions connect observed encouragement-arm contrasts to effects among compliers. We begin with consistency of treatment receipt and outcomes in the source population.

\begin{assumptionv}{ass:receipt-consistency}[Source treatment receipt consistency]
In the source population, observed receipt equals potential receipt under the realized encouragement:
\[
D=D(Z)\qquad\text{almost surely}.
\]
\end{assumptionv}

Receipt consistency connects realized treatment receipt to its potential value under the assigned encouragement \citep{ChenHuang2025Generalizing}. The outcome restriction completes this link between observed records and potential variables.

\begin{assumptionv}{ass:outcome-consistency}[Outcome consistency and boundedness]
In the source population, the observed outcome equals the potential outcome under the realized receipt and lies in the unit interval, both almost surely:
\[
Y=Y(D), \qquad 0\le Y\le 1.
\]
\end{assumptionv}

Outcome consistency identifies the observed response with the potential outcome under realized receipt \citep{ChenHuang2025Generalizing}; boundedness fixes its scale for the inference problem. We then require conditional randomization of encouragement.

\begin{assumptionv}{ass:instrument-randomization}[Conditional instrument randomization]
Let \(O\) denote the primitive full-data record and let \(P_n^{F,\mathrm S}=P_n^F(\,\cdot\mid S=1)\) be its source-population law. The joint source assignment law \(\mathbb P_{\mathrm{assign}}\) of \((O,Z)\) is a probability law, and its \(O\)-marginal is \(P_n^{F,\mathrm S}\). The zero extension of the source propensity,
\[
\widetilde e(x)=
\begin{cases}
e(x),&x\in[0,1],\\
0,&x\notin[0,1],
\end{cases}
\]
is measurable. Put \(\pi_1(x)=e(x)\) and \(\pi_0(x)=1-e(x)\). For every measurable full-data event \(B\) and each \(z\in\{0,1\}\),
\[
\mathbb P_{\mathrm{assign}}(O\in B, Z=z)
=\int_B \pi_z(X(O))\,dP_n^{F,\mathrm S}(O).
\]
Thus encouragement \(Z\) is conditionally independent of \((D(0),D(1),Y(0),Y(1))\) given \(X\) and source membership \(S=1\):
\[
Z\perp\!\!\!\perp (D(0),D(1),Y(0),Y(1))\mid X,S=1.
\]
\end{assumptionv}

Conditional instrument exchangeability makes source encouragement comparable across potential receipt and outcome types within covariate strata \citep{ChenHuang2025Generalizing}. This permits a covariate-dependent randomized encouragement design. The next restriction specifies the causal channel through which encouragement affects outcomes.

\begin{assumptionv}{ass:exclusion}[Exclusion through treatment receipt]
For each \(z\in\{0,1\}\), the source potential outcome generated by setting encouragement to \(z\) equals \(Y(D(z))\) almost surely under \(P_n^F\), the primitive full-data law at sample-size index \(n\). Here \(D(z)\) is potential receipt under encouragement \(z\), and \(Y(d)\) is the potential outcome under receipt \(d\).
\end{assumptionv}

Instrument exclusion assigns encouragement's outcome effect to the receipt it induces \citep{ChenHuang2025Generalizing}. Combined with conditional randomization and consistency, it supplies the outcome interpretation of the source arm contrast. Monotonicity supplies its principal-stratum interpretation.

\begin{assumptionv}{ass:monotonicity}[Monotonicity of treatment receipt]
The potential receipts \(D(1)\) and \(D(0)\), under unit and zero encouragement respectively, satisfy
\[
D(1)\ge D(0)
\]
almost surely under \(P_n^F\), the primitive full-data law at sample-size index \(n\).
\end{assumptionv}

The no-defiers monotonicity condition makes individual receipt weakly increasing in encouragement \citep{ChenHuang2025Generalizing}. Under these instrument restrictions, the source receipt contrast represents a conditional complier share, and the source outcome contrast represents a conditional complier-weighted effect; their transported interpretation is given in \cref{obj:lem:transported-cace-identification}.

Transport is imposed through two equalities at target-almost-every covariate value. The first carries the complier-weighted outcome contrast across populations.

\begin{assumptionv}{ass:transport-complier-outcome}[Transport of complier outcomes]
Let \(C\) be the indicator of the complier principal stratum, let \(Y(d)\) be the potential outcome under receipt \(d\), and let \(P_n^F\) be the primitive full-data law at sample-size index \(n\). For \(P_{\mathrm T}\)-almost every covariate value \(x\), where \(P_{\mathrm T}\) is the law of one observed target covariate,
\[
\mathbb E_{P_n^F}[(Y(1)-Y(0))C\mid X=x,S=0]
=
\mathbb E_{P_n^F}[(Y(1)-Y(0))C\mid X=x,S=1],
\]
where \(S=0\) denotes target membership and \(S=1\) denotes source membership.
\end{assumptionv}

After aligning encouragement potential outcomes with receipt-mediated outcomes under exclusion and monotonicity, this complier-weighted outcome restriction corresponds to Assumption 2 of \citet{ChenHuang2025Generalizing}. It transports the conditional contribution to the effect numerator, including local complier prevalence. The second equality transports that prevalence itself.

\begin{assumptionv}{ass:transport-complier-share}[Transport of complier shares]
Let \(C\) be the indicator of the complier principal stratum, and let \(P_n^F\) be the primitive full-data law at sample-size index \(n\). For \(P_{\mathrm T}\)-almost every covariate value \(x\), where \(P_{\mathrm T}\) is the law of one observed target covariate,
\[
\mathbb E_{P_n^F}[C\mid X=x,S=0]
=
\mathbb E_{P_n^F}[C\mid X=x,S=1],
\]
where \(S=0\) denotes target membership and \(S=1\) denotes source membership.
\end{assumptionv}

The pointwise complier-share equality is the transport condition chosen for the present model and implies the averaged equality in Assumption 4 of \citet{ChenHuang2025Generalizing}. It is substantively plausible when source participation changes covariate composition but, within the measured covariate strata, does not change how encouragement shifts receipt. Together, the two equalities carry the numerator and denominator of the target complier effect through the target covariate distribution. The contribution developed below concerns estimation over the rough nuisance class and the honest expected-length frontier under these identification conditions.



The source contrast \leanref{sym:\Delta_A}{\ensuremath{\Delta_A}} and target-averaged contrast \leanref{sym:T_A}{\ensuremath{T_A}} separate the source conditional response from target covariate weighting. The target complier share \leanref{sym:\mu}{\ensuremath{\mu}} determines the normalization of the target complier effect \leanref{sym:\theta}{\ensuremath{\theta}}. Because both potential outcomes lie in \([0,1]\), the effect domain \leanref{sym:\Theta}{\ensuremath{\Theta}} is \([-1,1]\). Positive compliance gives the conditional target effect a well-defined denominator.

\begin{assumptionv}{ass:positive-strength}[Positive transported first stage]
The transported compliance share and first stage, \(\mu\), satisfies
\[
\mu>0.
\]
\end{assumptionv}

Target instrument relevance requires a positive transported complier share \citep{ChenHuang2025Generalizing}. The model permits this share to approach zero as the sample-size index varies. Under the preceding causal and transport restrictions, \cref{obj:lem:transported-cace-identification} gives
\[
\mu=T_D,
\qquad
\theta=\frac{T_Y}{T_D}.
\]
Transported compliance therefore determines both the population represented by the causal effect and the first stage that scales the observed outcome contrast. The bounded domain gives confidence intervals a common length scale even when that first stage is small.

\subsection{Global honesty and expected length}

We collect the sampling, regularity and causal restrictions into a single law-level class before defining coverage.

\begin{definitionv}{def:model-class}[Model class $\mathcal M_n$]
The model class \(\mathcal M_n\) consists of law triples
\[
P=(P_n^F,P_{\mathrm S},P_{\mathrm T}),
\]
where \(P_n^F\) is the primitive full-data law, \(P_{\mathrm S}\) is the law of one observed source record, and \(P_{\mathrm T}\) is the law of one observed target covariate, satisfying all of the following restrictions at sample-size index \(n\):
\begin{itemize}
\item \textbf{(Sampling.)} \cref{obj:ass:source-iid,obj:ass:target-iid,obj:ass:sample-independence}.
\item \textbf{(Covariate densities.)} \cref{obj:ass:source-density-bounds,obj:ass:target-density-bounds,obj:ass:source-density-holder,obj:ass:target-density-holder}.
\item \textbf{(Propensity and arm means.)} \cref{obj:ass:propensity-overlap,obj:ass:propensity-holder,obj:ass:arm-means-holder}.
\item \textbf{(Causal restrictions.)} \cref{obj:ass:receipt-consistency,obj:ass:outcome-consistency,obj:ass:instrument-randomization,obj:ass:exclusion,obj:ass:monotonicity}.
\item \textbf{(Transport and strength.)} \cref{obj:ass:transport-complier-outcome,obj:ass:transport-complier-share,obj:ass:positive-strength}.
\end{itemize}
\end{definitionv}

A law triple \leanref{sym:P}{\ensuremath{P}} in the model class \leanref{sym:\mathcal M_n}{\ensuremath{\mathcal M_n}} specifies both the causal target and the observed two-sample experiment. For probabilities and expectations involving a data-dependent interval, the observed law is \(P_{\mathrm S}^{\otimes n}\otimes P_{\mathrm T}^{\otimes n}\). The following conventions fix the coverage level, the estimator's sample-size threshold and the measure used to evaluate interval length.



The noncoverage probability \leanref{sym:\alpha}{\ensuremath{\alpha}} is fixed across the model class. The sample-size threshold \leanref{sym:n_0}{\ensuremath{n_0}} governs use of the cubic estimator, and restricted Lebesgue length \leanref{sym:\lambda_{\Theta}}{\ensuremath{\lambda_\Theta}} measures uncertainty on the causal domain. A generic random interval \leanref{sym:C_n}{\ensuremath{C_n}} is a measurable function of the two observed samples: formally, its graph \(\{(\omega,t):t\in C_n(\omega)\}\) is measurable in the product of the observed sample space and the Borel \(\sigma\)-field on \(\Theta\). We require its coverage to hold uniformly over the full model class.

\begin{definitionv}{def:honest-intervals}[Honest interval class \(\mathfrak H_n\)]
The class of uniformly honest connected intervals is
\[
\mathfrak H_n=
\left\{
C_n:
C_n\subseteq\Theta\ \text{is a measurable connected interval},
\quad
\inf_{P\in\mathcal M_n}P(\theta\in C_n)\ge 1-\alpha
\right\}.
\]
Here \(\Theta\) is the bounded domain of the target-population complier effect \(\theta\), \(\mathcal M_n\) is the model class, and \(\alpha\) is the prescribed noncoverage probability.
\end{definitionv}

The honest interval class \leanref{sym:\mathfrak H_n}{\ensuremath{\mathfrak H_n}} imposes finite-sample coverage at every law in \(\mathcal M_n\), including laws with arbitrarily small positive transported compliance. Connectedness makes the reported set an interval, and containment in \(\Theta\) bounds its length by two. To evaluate how its expected length changes with first-stage strength, we next select a strength slice.

\begin{definitionv}{def:strength-slice}[First-stage strength slice]
The first-stage strength slice, with \(\mu(P)\) denoting the transported first stage under \(P\), is defined for real parameters \(c_f,C_f,L,a\) and every \(n\in\mathbb N\) by
\[
\mathcal M_n(a)
=
\left\{
P\in\mathcal M_n:
0<a\le\frac14,\quad
a\le\mu(P)\le\frac14
\right\}.
\]
Here \(\mathcal M_n\) is the model class of \cref{obj:def:model-class} with parameters \(c_f,C_f,L\).
\end{definitionv}

The evaluation floor \leanref{sym:a}{\ensuremath{a}} indexes the strength slice \leanref{sym:\mathcal M_n(a)}{\ensuremath{\mathcal M_n(a)}}, whose transported first stages lie between \(a\) and \(1/4\). It enters the laws over which expected length is evaluated. Coverage continues to be required over all of \(\mathcal M_n\), and the attaining interval in \cref{obj:def:score-interval} uses a single construction across these slices. This ordering separates global validity from the strength-specific evaluation of precision.

We express the resulting optimal expected-length criterion as follows.

\begin{definitionv}{def:length-frontier}[Minimax length frontier \(L_n(a)\)]
The minimax expected-length frontier is
\[
L_n(a)
=
\inf_{C_n\in\mathfrak H_n}
\sup_{P\in\mathcal M_n(a)}
\mathbb E_P\!\left[\lambda_{\Theta}(C_n)\right],
\qquad n\ge n_0,\quad 0<a\le\frac14,
\]
where \(\mathfrak H_n\) is the class of uniformly honest connected intervals, \(\mathcal M_n(a)\) is the model slice with transported first stage between \(a\) and \(1/4\), and \(\lambda_{\Theta}\) denotes Lebesgue length restricted to the effect domain \(\Theta\).
\end{definitionv}

The minimax expected length \leanref{sym:L_n(a)}{\ensuremath{L_n(a)}} compares globally honest intervals by their worst expected length on a specified strength slice. Its infimum ranges over procedures satisfying the full coverage requirement in \cref{obj:def:honest-intervals}, while its supremum evaluates precision on \cref{obj:def:strength-slice}. This criterion links the roughness of the observed nuisance functions to the uncertainty remaining about a transported complier effect at a given first-stage scale.


\section{Main results and an attaining interval}
\label{sec:main-results}

The minimax expected length has order
\[
L_n(a)\asymp \min\{1,n^{-1/3}/a\},
\qquad n\ge n_0,quad 0<a\le 1/4.
\]
A single observable interval attains the upper bound while maintaining coverage over the full model class, and the lower bound applies to every measurable confidence set, including disconnected ones. This section states the exact interval and frontier. The technical construction and supporting bounds appear in \cref{sec:upper-appendix,sec:lower-appendix}; full proofs are collected in \cref{sec:deferred-proofs}.

\subsection{Observable representation and cubic correction}

The construction collects the observable density components in a seven-coordinate vector. Its source arm densities and response-marked arm densities encode the conditional contrasts without estimating conditional means separately.

\begin{definitionv}{def:marked-density-vector}[Marked-density vector \(F\)]
The marked-density vector for the model in \cref{obj:def:model-class} is
\[
F(x)=\bigl(f_{\mathrm T}(x),q_0(x),q_1(x),
r_{Y0}(x),r_{Y1}(x),r_{D0}(x),r_{D1}(x)\bigr),
\]
where the source arm densities and response-marked arm densities are
\[
q_z(x)=f_{\mathrm S}(x)P_{\mathrm S}(Z=z\mid X=x),
\qquad
r_{Az}(x)=q_z(x)m_{Az}(x),
\qquad A\in\{Y,D\},\quad z\in\{0,1\}.
\]
Here \(m_{Az}(x)\) is the source conditional mean of response \(A\) in encouragement arm \(z\).
\end{definitionv}

Weighting the arm contrast by the target density gives the transported functional.

\begin{definitionv}{def:smooth-functional}[Transported contrast functional \(T_A\)]
The target-averaged source contrast for response \(A\in\{Y,D\}\) is
\[
T_A=\int_0^1\Phi_A(F(x))\,dx,
\qquad
\Phi_A(F)=f_{\mathrm T}
\left(\frac{r_{A1}}{q_1}-\frac{r_{A0}}{q_0}\right).
\]
\end{definitionv}

Under the causal and transport assumptions, \(T_D=\mu(P)\) is the transported first stage and \(T_Y=\mu(P)\theta(P)\); see \cref{obj:lem:transported-cace-identification}. Thus inference can invert the affine moment \(T_Y-vT_D\) without dividing by an estimated first stage.

For \(n\ge n_0=256\), each source and target sample is split into one pilot block and three evaluation blocks. Blockwise marked histograms estimate all seven coordinates. The pilot resolution satisfies \(K_0\asymp n^{4/5}\), while the quadratic and cubic correction resolutions satisfy \(K_2\asymp n^{4/3}\) and \(K_3\asymp n\), up to dyadic rounding. The transported-contrast estimator \(\widehat T_A\) expands \(\Phi_A\) around the clipped pilot and adds linear, quadratic, and cubic terms formed from separate evaluation blocks.

The three correction orders cancel the corresponding Taylor terms while retaining independence among repeated residual factors. At smoothness \(\beta=1/8\), the fourth-order remainder and correction variances are controlled by the uniform mean-square envelope constant \(C_{\mathrm{mse}}=C_{\mathrm{mse}}(c_f,C_f,L)\). Specifically,
\[
\sup_{P\in\mathcal M_n}
\mathbb E_P\!\left[(\widehat T_A-T_A)^2\right]
\le C_{\mathrm{mse}}n^{-2/3},
\qquad A\in\{Y,D\}.
\]
The exact histogram conventions, estimator, and computable constant are given in \cref{obj:def:cubic-estimator,obj:lem:marked-cubic-mse}.

\subsection{Score interval and sharp frontier}

The mean-square scale yields the following calibrated affine-score acceptance set and interval.

\begin{definitionv}{synth_10}[Calibrated score acceptance set]
The calibrated affine-score acceptance set \(\mathcal A_n\) is defined, for real inputs \(\alpha,c_f,C_f,L\), a sample size \(n\in\mathbb N\), and a pair of source and target samples each containing \(n\) records, by
\[
\mathcal A_n
=
\left\{
\theta\in[-1,1]:
\left|\widehat T_Y-\theta\widehat T_D\right|
\le
2\sqrt{\frac{2C_{\mathrm{mse}}}{\alpha}}\,
n^{-1/3}
\right\}.
\]
Here \(\widehat T_Y\) and \(\widehat T_D\) are the cubic estimators computed from the supplied samples with inputs \(c_f,C_f\), and \(C_{\mathrm{mse}}\) is the mean-square envelope evaluated at \(c_f,C_f,L\). Each cubic estimator equals zero when \(n<n_0\), the estimator threshold, and otherwise equals its pilot integral plus its linear, quadratic, and cubic corrections. The square root and negative power use their total real-valued conventions, including at zero.
\end{definitionv}

\begin{definitionv}{def:score-interval}[Inverted score interval \(I_n\)]
The inverted score interval is
\[
I_n=
\begin{cases}
\Theta,& n<n_0,\\
\mathcal A_n,& n\ge n_0\ \text{and}\ \mathcal A_n\ne\varnothing,\\
\{0\},& n\ge n_0\ \text{and}\ \mathcal A_n=\varnothing.
\end{cases}
\]
Here \(\mathcal A_n\) is the calibrated affine-score acceptance set, \(\Theta\) is the bounded effect domain, and \(n_0\) is the threshold for using the cubic estimator.
\end{definitionv}

The full-domain fallback supplies coverage below the estimator threshold. For larger samples, the procedure retains a nonempty affine-score acceptance interval and uses the stated singleton convention when the acceptance set is empty. The exact result combines all-positive-sample-size coverage with the two-sided frontier for \(n\ge256\).

\begin{theoremv}{thm:sharp-length-frontier-full-elbow}[Sharp length frontier with elbow]
Fix real constants \(\alpha,c_f,C_f,L\) satisfying:
\begin{itemize}
\item \textbf{(Noncoverage probability.)} \(0<\alpha<1\).
\item \textbf{(Lower density envelope.)} \(0<c_f<1\).
\item \textbf{(Upper density envelope.)} \(1<C_f\).
\item \textbf{(Hölder radius.)} \(1<L\).
\end{itemize}
Let \(I_n\) be the score interval of \cref{obj:def:score-interval}, and let \(L_n(a)\) be the minimax expected restricted length defined in \cref{obj:def:length-frontier}. Set
\[
n_0=256.
\]
There exist finite constants \(0<c_0<C_0\), depending only on \((\alpha,c_f,C_f,L)\), such that \(I_n\in\mathfrak H_n\), the class of honest intervals in \cref{obj:def:honest-intervals}, for every integer \(n>0\), and, for every integer \(n\ge n_0\) and every real \(a\) with \(0<a\le 1/4\),
\[
c_0\min\{1,n^{-1/3}/a\}
\le L_n(a)
\le C_0\min\{1,n^{-1/3}/a\}.
\]
Moreover, for every integer \(n\ge n_0\) and every law \(P\in\mathcal M_n\) as defined in \cref{obj:def:model-class}, writing
\[
\mu(P)=T_D(P)
\]
for the transported first stage, the same interval sequence satisfies
\[
\mathbb E_P\!\left[\lambda_{\Theta}(I_n)\right]
\le C_0\min\{1,n^{-1/3}/\mu(P)\},
\]
where \(\lambda_{\Theta}\) denotes Lebesgue length restricted to the effect domain \(\Theta\).
\end{theoremv}

The two regimes follow by comparing the score radius with the transported first stage. Above \(n^{-1/3}\), affine inversion turns contrast error into effect uncertainty of order \(n^{-1/3}/\mu(P)\). At or below that scale, the bounded effect domain caps length at constant order. The legal instrumental-variable mixture in \cref{obj:lem:legal-iv-mixture-full-elbow} produces matching separations in both regimes. The supporting upper and lower statements are \cref{obj:thm:honest-upper-full-elbow,obj:thm:honest-lower-full-elbow}.


\section{Discussion and open questions}
\label{sec:discussion}

The expected-length bounds in \cref{obj:thm:sharp-length-frontier-full-elbow} provide precision benchmarks for the declared source--target experiment: independent samples with \(n\) source records and \(n\) target covariate records, under the restrictions in \cref{obj:def:model-class}. For \(n\ge n_0\), fixed calibration and regularity parameters, and an evaluation floor \(0<a\le1/4\), the two regimes describe the interaction between transported compliance and the rough estimation scale. When \(a>n^{-1/3}\), the optimal worst-case expected length has order \(n^{-1/3}/a\). When \(a\le n^{-1/3}\), it has constant order. Thus increasing the common sample size improves the precision benchmark once the strength floor exceeds the estimation scale, while global honesty requires positive worst-case expected length across the weaker-strength slices.

The floor \(a\) indexes the performance criterion in \cref{obj:def:strength-slice,obj:def:length-frontier}; the interval in \cref{obj:def:score-interval} uses the observed transported contrasts. Its law-specific expected-length guarantee in \cref{obj:thm:honest-upper-full-elbow} relates precision to the actual transported first stage \(\mu(P)\). A single interval therefore supplies global coverage and responds to the compliance strength of the underlying law, with inverse-strength control of expected length above the rough estimation scale and bounded length throughout the effect domain. This guarantee concerns expected length under repeated sampling; it is compatible with variation in the reported interval across realizations.

The fixed value \(\beta=1/8\) specifies the rough-nuisance regularity class analyzed here. At this regularity, the cubic correction is used to control the nonlinear remainder at the required scale. The split-sample estimator includes linear, quadratic, and cubic Taylor corrections across all seven learned marked-density coordinates and has mean-square error bounded by \(C_{\mathrm{mse}}n^{-2/3}\) under the model assumptions for \(n\ge256\). The resulting contrast-estimation radius is therefore of order \(n^{-1/3}\), which becomes the elbow in the length frontier.

The displayed finite-sample radius, \(2\sqrt{2C_{\mathrm{mse}}/\alpha}\,n^{-1/3}\), makes the hidden rate constant explicit: it depends on the density envelopes, the common H\"older radius, and the desired coverage level through \(C_{\mathrm{mse}}\) and \(\alpha\). The threshold \(n_0=256\) ensures that the deterministic four-block construction and its resolution choices are defined. Numerical precision is governed by the calibration constant and the transported first stage. For example, take the admissible inputs \(c_f=0.9\), \(C_f=1.1\), \(L=1.01\), and \(\alpha=0.05\). Substitution into the displayed formula in \cref{obj:lem:marked-cubic-mse} gives \(C_{\mathrm{mse}}\approx1.16\times10^{27}\), so the certified score tolerance is approximately \(4.31\times10^{14}n^{-1/3}\): about \(6.79\times10^{13}\) at \(n=256\) and \(4.31\times10^{12}\) at \(n=10^6\). These values exceed the effect-domain diameter and show that the explicit worst-case envelope is highly conservative; a realized affine-score interval can be shorter, but that realized length is not the uniform certificate. An analyst supplies fixed, scientifically defensible bounds on density overlap and Hölder size before applying this finite-sample calibration. The expected-length bound then uses the corresponding \(C_0\) and the law-specific first stage.

An application of this benchmark requires the source encouragement experiment to satisfy conditional randomization, exclusion, and monotonicity, together with the sampling and regularity conditions of the model. Transport to the target population rests on the conditional equalities for the complier-weighted outcome contrast and the complier share in \cref{obj:ass:transport-complier-outcome,obj:ass:transport-complier-share}. Under these conditions, the precision benchmark concerns the target-population average treatment effect among compliers, using source receipt and outcome observations together with target covariate observations.

Separate recruitment costs motivate the unequal-sample extension of the same decision problem in \cref{obj:def:unequal-sample-handle}. Let \(N_{\mathrm S}\) denote the source sample size and \(N_{\mathrm T}\) the target sample size. Retaining the law-level restrictions while allowing these counts to differ gives the following criterion.

\begin{definitionv}{def:unequal-sample-handle}[Unequal-sample frontier \(R_{N_{\mathrm S},N_{\mathrm T}}(a)\)]
The unequal-sample minimax expected-length frontier is
\[
R_{N_{\mathrm S},N_{\mathrm T}}(a)
=
\inf_{C_{N_{\mathrm S},N_{\mathrm T}}\in
\mathfrak H^{\mathrm u}_{N_{\mathrm S},N_{\mathrm T}}}
\sup_{P\in\mathcal M^{\mathrm u}_{N_{\mathrm S},N_{\mathrm T}}(a)}
\mathbb E_P\!\left[
\lambda_{\Theta}\!\left(C_{N_{\mathrm S},N_{\mathrm T}}\right)
\right],
\qquad
0<a\le\frac14,
\]
for positive integers \(N_{\mathrm S}\) and \(N_{\mathrm T}\), where \(\lambda_{\Theta}\) denotes Lebesgue length restricted to the effect domain \(\Theta\).

Here \(\mathcal M^{\mathrm u}_{N_{\mathrm S},N_{\mathrm T}}\) is the class of full-data and one-record observed-law triples satisfying every law-level density, smoothness, overlap, causal, transport, and positive-strength condition of the equal-sample model class \(\mathcal M_n\), observed under
\[
P_{\mathrm S}^{\otimes N_{\mathrm S}}
\otimes
P_{\mathrm T}^{\otimes N_{\mathrm T}},
\]
where \(P_{\mathrm S}\) is the law of one observed source record and \(P_{\mathrm T}\) is the law of one observed target covariate. The independent source and target samples contain \(N_{\mathrm S}\) records \((X,Z,D,Y)\) and \(N_{\mathrm T}\) records \(X\), respectively. With \(\mu(P)\) denoting the transported first stage, the strength slice is
\[
\mathcal M^{\mathrm u}_{N_{\mathrm S},N_{\mathrm T}}(a)
=
\left\{
P\in\mathcal M^{\mathrm u}_{N_{\mathrm S},N_{\mathrm T}}:
a\le\mu(P)\le\frac14
\right\}.
\]

Let \(C_{N_{\mathrm S},N_{\mathrm T}}\) range over measurable connected random intervals contained in \(\Theta\) and based on this product experiment. With \(\theta(P)\) denoting the target-population average treatment effect among compliers and \(\alpha\) the prescribed noncoverage probability, define
\[
\mathfrak H^{\mathrm u}_{N_{\mathrm S},N_{\mathrm T}}
=
\left\{
C_{N_{\mathrm S},N_{\mathrm T}}:
\inf_{P\in\mathcal M^{\mathrm u}_{N_{\mathrm S},N_{\mathrm T}}}
P\!\left(
\theta(P)\in C_{N_{\mathrm S},N_{\mathrm T}}
\right)
\ge1-\alpha
\right\}.
\]
\end{definitionv}

The unequal-sample minimax expected-length criterion \(R_{N_{\mathrm S},N_{\mathrm T}}(a)\) preserves the distinction between global coverage and performance on a strength slice. Coverage is required over the entire unequal-sample model, while the supremum defining expected length evaluates laws whose transported first stage lies between the specified floor and \(1/4\), inclusive. The two sample counts consequently enter through the observation experiment, with the population restrictions and the estimand held fixed.

\subsection{Open questions}

The unequal-sample criterion makes allocation between source outcomes and target covariates a statistical design question. Characterizing its order would provide a basis for comparing the precision gains from additional records in either sample, subject to their respective recruitment costs. The central question is the joint behavior of the criterion across sample-size and strength regimes.

\begin{remarkv}{oeq:unequal-sample-frontier}[Unequal-sample length phases]
What is the exact piecewise order of \(R_{N_{\mathrm S},N_{\mathrm T}}(a)\), the unequal-sample minimax expected-length frontier, when the source sample size \(N_{\mathrm S}\) and target sample size \(N_{\mathrm T}\) diverge at unrelated rates? Can a single finite-sample honest interval attain this order in every phase, together with a matching legal instrumental-variable lower family satisfying the no-defiers restriction?
\end{remarkv}

Resolving \cref{obj:oeq:unequal-sample-frontier} would connect an attaining procedure and an information bound for each allocation regime. Requiring the lower family to satisfy the no-defiers restriction keeps that comparison within the same causal model as the proposed intervals.

\subsection{Limitations and future work}

The established benchmark uses a scalar covariate, bounded outcomes, fixed Hölder smoothness with exponent \(\beta=1/8\), and independent records within and across the two samples. Its matching expected-length bounds concern equal source and target sample sizes. The unequal-allocation problem in \cref{obj:def:unequal-sample-handle} defines an extension of the criterion; its general rate characterization and an interval attaining that characterization remain the questions in \cref{obj:oeq:unequal-sample-frontier}.

Further work could examine higher-dimensional covariates, inference under unknown smoothness, and sampling structures with dependent records. These extensions require corresponding coverage and expected-length analyses. Applications with alternative population transport restrictions likewise require an identification analysis suited to those restrictions before the precision benchmark can be extended.


\appendix

\section*{Appendices}

\section{Identification and upper-bound proofs}
\label{sec:upper-appendix}

\subsection{Identification from primitive variables}

The causal interpretation of the transported contrasts rests on consistency, conditional instrument randomization, exclusion, and monotonicity. The two transport restrictions then carry the source conditional contrasts to target complier quantities.

\begin{lemmav}{lem:transported-cace-identification}[Transported complier effect identification]
Let \(n\in\mathbb N\) and let \(P\) be a transport law satisfying:
\begin{itemize}
\item \textbf{(Lower envelope.)} \(0<c_f<1\), where \(c_f\) is the lower density envelope.
\item \textbf{(Upper envelope.)} \(1<C_f\), where \(C_f\) is the upper density envelope.
\item \textbf{(Smoothness radius.)} \(1<L\), where \(L\) is the common Hölder radius.
\item \textbf{(Model membership.)} \(P\in\mathcal M_n\), the model class of \cref{obj:def:model-class} with constants \(c_f,C_f,L\).
\end{itemize}
Write \(P_n^F\) for the full-data law, \(D(z)\) for potential receipt under encouragement \(z\), and \(Y(d)\) for potential outcome under receipt \(d\). For each \(z\in\{0,1\}\), there are versions of the source conditional means such that, almost everywhere under the source covariate law,
\[
m_{Dz}(x)
 =\mathbb E_{P_n^F}[D(z)\mid X=x,S=1],
\qquad
m_{Yz}(x)
 =\mathbb E_{P_n^F}[Y(D(z))\mid X=x,S=1].
\]
Here \(m_{Dz}\) and \(m_{Yz}\) are the source receipt and outcome arm means.

Define the source arm contrasts and the complier indicator by
\[
\Delta_D(x)=m_{D1}(x)-m_{D0}(x),
\qquad
\Delta_Y(x)=m_{Y1}(x)-m_{Y0}(x),
\qquad
C=\mathbf 1\{D(1)=1,\ D(0)=0\}.
\]
There are common versions of the conditional means in each of the following chains for which, for \(P_{\mathrm T}\)-almost every \(x\), where \(P_{\mathrm T}\) is the target covariate law,
\[
\begin{aligned}
\Delta_D(x)
&=\mathbb E_{P_n^F}[D(1)-D(0)\mid X=x,S=1]\\
&=\mathbb E_{P_n^F}[C\mid X=x,S=1]
 =\mathbb E_{P_n^F}[C\mid X=x,S=0],\\
\Delta_Y(x)
&=\mathbb E_{P_n^F}[Y(D(1))-Y(D(0))\mid X=x,S=1]\\
&=\mathbb E_{P_n^F}[(Y(1)-Y(0))C\mid X=x,S=1]\\
&=\mathbb E_{P_n^F}[(Y(1)-Y(0))C\mid X=x,S=0].
\end{aligned}
\]
Moreover,
\[
D(1)-D(0)=C
\qquad P_n^F\text{-almost surely}.
\]
For \(A\in\{D,Y\}\), define the transported contrast using the target covariate density \(f_{\mathrm T}\) by
\[
T_A=\int \Delta_A(x)f_{\mathrm T}(x)\,dx.
\]
The transported first stage equals the target complier share, and the transported outcome contrast equals the target complier-weighted outcome effect:
\[
T_D=\mu=P_n^F(C=1\mid S=0),
\qquad
T_Y=\mathbb E_{P_n^F}[(Y(1)-Y(0))C\mid S=0].
\]
The target complier average treatment effect \(\theta\) satisfies
\[
\theta
=\frac{\mathbb E_{P_n^F}[(Y(1)-Y(0))C\mid S=0]}
       {P_n^F(C=1\mid S=0)}
=\frac{T_Y}{T_D}
\in\Theta,
\]
where \(\Theta\) is the effect domain specified by the model.
\end{lemmav}

\begin{proof}[Proof of \cref{obj:lem:transported-cace-identification}]
Fix \(P\in\mathcal M_n\) as in the statement. Throughout, conditional means are understood through their defining integral identities over measurable covariate sets.

\begin{enumerate}
\item Choose common conditional-mean versions for the transported quantities. By \cref{obj:ass:transport-complier-share,obj:ass:transport-complier-outcome}, there are measurable functions \(q_D,q_Y:\mathbb R\to\mathbb R\) satisfying
\[
\begin{aligned}
q_D(x)
&=\mathbb E_{P_n^F}[C\mid X=x,S=0]
 =\mathbb E_{P_n^F}[C\mid X=x,S=1],\\
q_Y(x)
&=\mathbb E_{P_n^F}[(Y(1)-Y(0))C\mid X=x,S=0]\\
&=\mathbb E_{P_n^F}[(Y(1)-Y(0))C\mid X=x,S=1],
\end{aligned}
\]
as versions in the respective populations.

By \cref{obj:ass:monotonicity}, the binary receipt pair
\((D(0),D(1))\) belongs almost surely to
\(\{(0,0),(0,1),(1,1)\}\). For these three pairs, respectively, both the receipt difference and the complier indicator equal \(0,1,0\). Consequently,
\[
D(1)-D(0)=C
\qquad P_n^F\text{-almost surely}.
\]
Checking the four binary receipt pairs also gives the pointwise identity
\[
Y(D(1))-Y(D(0))
=(Y(1)-Y(0))(D(1)-D(0)).
\]
Substitution therefore yields
\[
Y(D(1))-Y(D(0))=(Y(1)-Y(0))C
\qquad P_n^F\text{-almost surely}.
\]
Each population law obtained by conditioning on \(S\) is absolutely continuous with respect to \(P_n^F\). Thus these identities persist in the source population and in every restricted integral defining a conditional mean. In particular,
\[
\begin{aligned}
q_D(x)&=\mathbb E_{P_n^F}[D(1)-D(0)\mid X=x,S=1],\\
q_Y(x)&=\mathbb E_{P_n^F}[Y(D(1))-Y(D(0))\mid X=x,S=1]
\end{aligned}
\]
are source conditional-mean versions.
% lean: CausalSmith.Stat.TransportCaceRoughnuisanceLength.conditionalMean_congr_full_ae

\item Construct source arm-mean versions. Define the covariate domain, potential arm responses, zero extensions, and assignment weights by
\[
\mathcal X=[0,1],\qquad
A(z)=
\begin{cases}
D(z),&A=D,\\
Y(D(z)),&A=Y,
\end{cases}
\quad A\in\{D,Y\},\ z\in\{0,1\},
\]
\[
\widetilde m_{Az}(x)=
\begin{cases}
m_{Az}(x),&x\in\mathcal X,\\
0,&x\notin\mathcal X,
\end{cases}
\qquad
\pi_z(x)=
\begin{cases}
e(x),&z=1,\\
1-e(x),&z=0,
\end{cases}
\quad x\in\mathbb R.
\]
Also write \(P_{\mathrm S}^{X}\) for the source covariate marginal:
\[
P_{\mathrm S}^{X}(B)=P_{\mathrm S}(X\in B),
\qquad B\subseteq\mathbb R\text{ Borel}.
\]
The source marginal identification supplied by conditional randomization gives
\[
P_{\mathrm S}^{X}(B)=P_n^F(X\in B\mid S=1).
\]
By \cref{obj:ass:source-density-bounds}, this marginal is supported on \(\mathcal X\), so
\[
\widetilde m_{Az}=m_{Az}
\qquad P_{\mathrm S}^{X}\text{-almost everywhere}.
\]
The regularity and range conditions in \cref{obj:ass:arm-means-holder} make the zero extensions measurable and give
\[
0\le \widetilde m_{Az}(x)\le1
\quad(x\in\mathbb R).
\]
The binary receipts and potential-outcome bounds in
\cref{obj:ass:source-density-bounds} give, in the source population,
\[
0\le A(z)\le1
\qquad\text{almost surely}.
\]
Moreover, \cref{obj:ass:propensity-overlap} gives
\[
\frac14\le\pi_z(x)\le\frac34
\qquad(x\in\mathcal X).
\]
These bounds ensure the integrability of the variables used below.

For a Borel set \(B\subseteq\mathbb R\), put
\[
B_{\mathcal X}=B\cap\mathcal X.
\]
The source arm-moment identity in \cref{obj:ass:arm-means-holder} gives
\[
\int_{B_{\mathcal X}}
 f_{\mathrm S}(x)\pi_z(x)m_{Az}(x)\,dx
=
\mathbb E_{P_{\mathrm S}}
 \bigl[\mathbf1_{\{X\in B_{\mathcal X},\,Z=z\}}A\bigr].
\]
Receipt and outcome consistency, as supplied by
\cref{obj:ass:receipt-consistency,obj:ass:outcome-consistency},
replace the observed response on the arm \(Z=z\) by \(A(z)\).
Conditional randomization in \cref{obj:ass:instrument-randomization}
makes the full-data law restricted to that arm have density
\(\pi_z(X)\) relative to the source full-data law. Hence
\[
\begin{aligned}
\mathbb E_{P_{\mathrm S}}
 \bigl[\mathbf1_{\{X\in B_{\mathcal X},\,Z=z\}}A\bigr]
&=\mathbb E
 \bigl[\mathbf1_{\{X\in B_{\mathcal X},\,Z=z\}}A(z)\bigr]\\
&=\mathbb E_{P_n^F}
 \bigl[\mathbf1_{\{X\in B_{\mathcal X}\}}\pi_z(X)A(z)
       \mid S=1\bigr],
\end{aligned}
\]
where the expectation in the middle is under the source assignment law. The last equality follows first for indicators of measurable full-data events from the assignment probability identity, and then for the displayed bounded measurable integrand by integration.

Using the source density identity and source support, the preceding equalities become
\[
\begin{aligned}
&\mathbb E_{P_n^F}
 \bigl[\mathbf1_{\{X\in B\}}\pi_z(X)\widetilde m_{Az}(X)
       \mid S=1\bigr]\\
&\qquad=
\int_{B_{\mathcal X}}
 f_{\mathrm S}(x)\pi_z(x)m_{Az}(x)\,dx\\
&\qquad=
\mathbb E_{P_n^F}
 \bigl[\mathbf1_{\{X\in B\}}\pi_z(X)A(z)\mid S=1\bigr].
\end{aligned}
\]
Thus \(\pi_z(X)\widetilde m_{Az}(X)\) is a version of the conditional mean of \(\pi_z(X)A(z)\) given \(X\) in the source population. Since \(\pi_z(X)\) is measurable with respect to \(X\), the conditional-expectation multiplication rule gives
\[
\pi_z(X)\widetilde m_{Az}(X)
=
\mathbb E_{P_n^F}[\pi_z(X)A(z)\mid X,S=1]
=
\pi_z(X)\mathbb E_{P_n^F}[A(z)\mid X,S=1].
\]
The positive lower bound on \(\pi_z(X)\) permits cancellation:
\[
\widetilde m_{Az}(X)
=\mathbb E_{P_n^F}[A(z)\mid X,S=1]
\qquad\text{almost surely in the source population}.
\]
Equivalently, for every Borel \(B\subseteq\mathbb R\),
\[
\mathbb E_{P_n^F}
 \bigl[\mathbf1_{\{X\in B\}}A(z)\mid S=1\bigr]
=
\int_B\widetilde m_{Az}(x)\,dP_{\mathrm S}^{X}(x).
\]
% lean: CausalSmith.Stat.TransportCaceRoughnuisanceLength.armMeanVersion_conditionalMean

\item Identify the contrasts with the common versions. Define, for \(A\in\{D,Y\}\),
\[
\widetilde\Delta_A(x)
=\widetilde m_{A1}(x)-\widetilde m_{A0}(x),
\qquad
q_A=
\begin{cases}
q_D,&A=D,\\
q_Y,&A=Y.
\end{cases}
\]
The source almost-everywhere equalities for the arm means imply
\[
\Delta_A=\widetilde\Delta_A
\qquad P_{\mathrm S}^{X}\text{-almost everywhere}.
\]
Subtracting the two arm integral identities from the preceding step gives
\[
\int_B\widetilde\Delta_A(x)\,dP_{\mathrm S}^{X}(x)
=
\mathbb E_{P_n^F}
 \bigl[\mathbf1_{\{X\in B\}}(A(1)-A(0))\mid S=1\bigr]
=
\int_Bq_A(x)\,dP_{\mathrm S}^{X}(x)
\]
for every Borel \(B\subseteq\mathbb R\).

To obtain almost-everywhere equality from these identities, define
\[
\mathcal E_{k,A}
=
\bigl\{x\in\mathbb R:
 |\widetilde\Delta_A(x)|\le k,\ |q_A(x)|\le k\bigr\},
\qquad k\in\mathbb N.
\]
Both functions are measurable and real-valued, and
\[
\bigcup_{k\in\mathbb N}\mathcal E_{k,A}=\mathbb R.
\]
On each \(\mathcal E_{k,A}\), they are integrable under the finite source covariate law. Applying the preceding integral equality to
\(B\cap\mathcal E_{k,A}\), uniqueness for integrable functions with equal integrals over every measurable set yields
\[
\widetilde\Delta_A=q_A
\qquad
P_{\mathrm S}^{X}\text{-almost everywhere on }\mathcal E_{k,A}.
\]
Taking the countable union, and then using the equality between the raw and extended contrasts, gives
\[
\Delta_D=q_D,\qquad \Delta_Y=q_Y
\qquad P_{\mathrm S}^{X}\text{-almost everywhere}.
\]

The density identities in
\cref{obj:ass:source-density-bounds,obj:ass:target-density-bounds}
transfer these equalities to the target covariate law. Indeed, if a Borel set \(B\) is source-null, then
\[
0=P_{\mathrm S}^{X}(B)
=\int_{B\cap\mathcal X}f_{\mathrm S}(x)\,dx
\ge c_f\int_{B\cap\mathcal X}1\,dx.
\]
Since \(c_f>0\), the set \(B\cap\mathcal X\) is Lebesgue-null, and consequently
\[
P_{\mathrm T}(B)
=\int_{B\cap\mathcal X}f_{\mathrm T}(x)\,dx=0.
\]
Thus
\[
P_{\mathrm T}\ll P_{\mathrm S}^{X},
\qquad
\Delta_D=q_D,\quad\Delta_Y=q_Y
\quad P_{\mathrm T}\text{-almost everywhere}.
\]
Together with the conditional-mean identities for \(q_D,q_Y\), this supplies common versions for both chains in the statement.
% lean: CausalSmith.Stat.TransportCaceRoughnuisanceLength.target_absolutelyContinuous_source

\item Integrate the common versions in the target population. The target density identity and its positive envelope in
\cref{obj:ass:target-density-bounds} convert the defining transported integrals into integrals under \(P_{\mathrm T}\). Applying the target conditional-mean identities to the whole covariate space gives
\[
\begin{aligned}
T_D
&=\int_{\mathcal X}f_{\mathrm T}(x)\Delta_D(x)\,dx\\
&=\int_{\mathbb R}\Delta_D(x)\,dP_{\mathrm T}(x)
 =\int_{\mathbb R}q_D(x)\,dP_{\mathrm T}(x)\\
&=\mathbb E_{P_n^F}[C\mid S=0]
 =P_n^F(C=1\mid S=0)=\mu.
\end{aligned}
\]
The penultimate equality uses
\(C=\mathbf1_{\{C=1\}}\), so its integral is the probability of the complier event. Similarly,
\[
\begin{aligned}
T_Y
&=\int_{\mathcal X}f_{\mathrm T}(x)\Delta_Y(x)\,dx\\
&=\int_{\mathbb R}\Delta_Y(x)\,dP_{\mathrm T}(x)
 =\int_{\mathbb R}q_Y(x)\,dP_{\mathrm T}(x)\\
&=\mathbb E_{P_n^F}[(Y(1)-Y(0))C\mid S=0].
\end{aligned}
\]
% lean: CausalSmith.Stat.TransportCaceRoughnuisanceLength.transported_cace_identification

\item Complete the arm identities and the effect-domain assertion. For each \(z\in\{0,1\}\), the constructed versions satisfy
\[
\begin{aligned}
m_{Dz}(x)
&=\widetilde m_{Dz}(x)
 =\mathbb E_{P_n^F}[D(z)\mid X=x,S=1],\\
m_{Yz}(x)
&=\widetilde m_{Yz}(x)
 =\mathbb E_{P_n^F}[Y(D(z))\mid X=x,S=1]
\end{aligned}
\]
for \(P_{\mathrm S}^{X}\)-almost every \(x\). These are the required source arm identities.

By the definition of the target complier effect and the two transported integral identities,
\[
\theta
=
\frac{\mathbb E_{P_n^F}[(Y(1)-Y(0))C\mid S=0]}
     {P_n^F(C=1\mid S=0)}
=
\frac{T_Y}{T_D}.
\]
By \cref{obj:ass:positive-strength}, \(T_D=\mu>0\).

Finally, the potential-outcome support restriction in
\cref{obj:ass:source-density-bounds} persists under the target population law. On the event \(C=1\),
\[
-1\le Y(1)-Y(0)\le1,
\]
whereas on \(C=0\) the complier-weighted effect is zero. Hence, target-almost surely,
\[
-C\le (Y(1)-Y(0))C\le C.
\]
All three variables are integrable by their bounds. Integration therefore gives
\[
-\mu
\le\mathbb E_{P_n^F}[(Y(1)-Y(0))C\mid S=0]
=T_Y
\le\mu.
\]
Dividing by \(\mu=T_D>0\) yields
\[
-1\le\frac{T_Y}{T_D}=\theta\le1,
\qquad\text{so}\qquad
\theta\in[-1,1]=\Theta.
\]
% lean: CausalSmith.Stat.TransportCaceRoughnuisanceLength.targetCACE_mem_parameterSpace
\end{enumerate}
\end{proof}

The common-version qualification makes the conditional identities compatible with target integration. Positive density envelopes provide common covariate support, monotonicity identifies the receipt contrast with complier prevalence, and exclusion identifies the outcome contrast with the complier-weighted receipt effect. Positive transported strength makes their ratio the target complier average effect.

\subsection{Coordinate and sampling conventions}

The estimator uses one target-density coordinate and six marked source coordinates. The next two definitions fix the coordinate names, observable marks, channels, and index conversion once for the remaining construction.

\begin{definitionv}{synth_3}[Exact block size]
The block size \(m_b\), for a nonnegative integer \(n\) and an index \(b\in\{0,1,2,3\}\), is defined by
\[
m_b := \max\left\{
\left\lfloor \frac{(b+1)n}{4} \right\rfloor
-
\left\lfloor \frac{bn}{4} \right\rfloor,
0
\right\}.
\]
\end{definitionv}

\begin{definitionv}{synth_2}[Deterministic sample blocks]
The deterministic block \(\mathcal B_b^G\) in sample channel \(G\), for sample size \(n\in\{0,1,2,\ldots\}\) and block index \(b\in\{0,1,2,3\}\), uses zero-based record indices and is defined by
\[
\mathcal B_b^G
=
\left\{
i\in\mathbb Z:
0\le i<n,\quad
\left\lfloor\frac{bn}{4}\right\rfloor
\le i<
\left\lfloor\frac{(b+1)n}{4}\right\rfloor
\right\}.
\]
\end{definitionv}

\begin{definitionv}{synth_19}[Target-density coordinate]
In the marked-density vector, \(t\) denotes the target covariate density:
\[
t:=f_{\mathrm T},\qquad t(x)=f_{\mathrm T}(x),\quad x\in[0,1].
\]
Thus \(F_0=t=f_{\mathrm T}\), and the named-coordinate subscript \(t\) denotes coordinate \(0\): \(V_t:=V_0\) for \(V=F\), \(V=H_K^{(b)}\), \(V=R_K^{(b)}\), or \(V=W_r\). In particular, \(F_t=f_{\mathrm T}\). In \(\operatorname{bump}(t)\) and \(\int_{\Theta}P(t\in C_n)\,dt\), \(t\) is instead a local real argument or integration variable, respectively.
\end{definitionv}

\begin{definitionv}{synth_14}[Coordinate marks \(W_{i,r}\) and sample blocks \(\mathcal B_b^G\)]
The coordinate index set and coordinate ordering are
\[
\mathcal I=\{0,\ldots,6\},
\qquad
(F_0,F_1,F_2,F_3,F_4,F_5,F_6)
=(t,q_0,q_1,r_{Y0},r_{Y1},r_{D0},r_{D1}).
\]
The sample channel of coordinate \(i\in\mathcal I\) is
\[
g(i)=
\begin{cases}
\mathrm T,&i=0,\\
\mathrm S,&i\in\{1,\ldots,6\}.
\end{cases}
\]
Vector, histogram, residual, and mark subscripts follow the correspondence
\[
\begin{array}{c|ccccccc}
i&0&1&2&3&4&5&6\\ \hline
\text{coordinate}&t&q_0&q_1&r_{Y0}&r_{Y1}&r_{D0}&r_{D1}.
\end{array}
\]
For \(V=F\), \(V=H_K^{(b)}\), \(V=R_K^{(b)}\), or \(V=W_r\), the named-coordinate subscripts satisfy
\[
(V_t,V_{q_0},V_{q_1},V_{r_{Y0}},V_{r_{Y1}},V_{r_{D0}},V_{r_{D1}})
=(V_0,V_1,V_2,V_3,V_4,V_5,V_6).
\]

Index the source records and target covariate records from zero:
\[
O_r^{\mathrm S}=(X_r^{\mathrm S},Z_r,D_r,Y_r),
\qquad
X_r^{\mathrm T},
\qquad r=0,\ldots,n-1.
\]
Let \(Y_r^{\mathrm{clip}}\) denote the clipped observable outcome used in
\cref{obj:synth_7,obj:synth_8,obj:def:cubic-estimator}.
The observable coordinate marks, with this same clipping convention throughout, are
\[
\begin{aligned}
W_{0,r}&=1,\\
W_{1,r}&=\mathbf 1\{Z_r=0\},
&
W_{2,r}&=\mathbf 1\{Z_r=1\},\\
W_{3,r}&=\mathbf 1\{Z_r=0\}Y_r^{\mathrm{clip}},
&
W_{4,r}&=\mathbf 1\{Z_r=1\}Y_r^{\mathrm{clip}},\\
W_{5,r}&=\mathbf 1\{Z_r=0\}D_r,
&
W_{6,r}&=\mathbf 1\{Z_r=1\}D_r.
\end{aligned}
\]
The mark \(W_{0,r}\) accompanies target record \(r\); the remaining marks accompany source record \(r\).

For \(G\in\{\mathrm S,\mathrm T\}\) and \(b\in\{0,1,2,3\}\), the blocks and their sizes are
\[
\mathcal B_b^G
=
\left\{
r\in\{0,\ldots,n-1\}:
\left\lfloor\frac{bn}{4}\right\rfloor
\le r<
\left\lfloor\frac{(b+1)n}{4}\right\rfloor
\right\},
\qquad
m_b=\left\lfloor\frac{(b+1)n}{4}\right\rfloor
-\left\lfloor\frac{bn}{4}\right\rfloor.
\]
Their one-based counterparts are
\[
\widetilde{\mathcal B}_b^G
=\{r+1:r\in\mathcal B_b^G\}
=\{\lfloor bn/4\rfloor+1,\ldots,\lfloor(b+1)n/4\rfloor\}.
\]
One-based covariates and marks are defined together by
\[
\widetilde X_s^G=X_{s-1}^G,
\qquad
\widetilde W_{i,s}=W_{i,s-1},
\qquad s=1,\ldots,n.
\]
Thus, for a covariate set \(E\), the corresponding block sums satisfy
\[
\sum_{s\in\widetilde{\mathcal B}_b^{g(i)}}
\widetilde W_{i,s}\,
\mathbf 1\{\widetilde X_s^{g(i)}\in E\}
=
\sum_{r\in\mathcal B_b^{g(i)}}
W_{i,r}\,\mathbf 1\{X_r^{g(i)}\in E\}.
\]
Block \(0\) supplies the training records, and blocks \(1,2,3\) supply the evaluation records. Each coordinate-\(i\) sum uses covariates and marks with the same record index in channel \(g(i)\).
\end{definitionv}

\subsection{Histograms and cubic estimator}

A normalized marked histogram estimates each coordinate on a common dyadic partition. Block 0 trains the clipped pilot; blocks 1--3 provide independent factors for the correction terms.

\begin{definitionv}{synth_7}[Blockwise marked histogram definition]
The marked histogram \(H_{i,K}^{(b)}(x)\), for coordinate \(i\), resolution \(K\), block \(b\), and evaluation point \(x\), is defined by
\[
H_{i,K}^{(b)}(x)=
\begin{cases}
0,&n<256,\\[2mm]
\displaystyle
\frac{K}{m_b}
\sum_{\ell=0}^{K-1}
\mathbf 1\{x\in E_{\ell,K}\}
\sum_{r\in\mathcal B_b^{g(i)}}
W_{i,r}\,
\mathbf 1\{X_r^{g(i)}\in E_{\ell,K}\},
&n\ge256.
\end{cases}
\]
Here \(\mathcal S_n\), the sample of source records, and \(\mathcal T_n\), the sample of target covariates, each contain \(n\) records, and the parameters range over
\[
n,K\in\{0,1,\ldots\},\qquad
i\in\mathcal I=\{0,\ldots,6\},\qquad
b\in\{0,\ldots,3\},\qquad
x\in\mathbb R.
\]
For \(K\ge1\), the cells \(E_{\ell,K}\) are
\[
E_{\ell,K}=
\begin{cases}
[\ell/K,(\ell+1)/K),&0\le\ell<K-1,\\
[\ell/K,1],&\ell=K-1.
\end{cases}
\]
For \(K=0\), the cell sum is empty and \(H_{i,0}^{(b)}(x)=0\).

Both samples use zero-based record indices \(r\in\{0,\ldots,n-1\}\), with an empty index set when \(n=0\). For sample channel \(G\in\{\mathrm S,\mathrm T\}\), the deterministic block \(\mathcal B_b^G\) and its size \(m_b\) are
\[
\mathcal B_b^G
=
\left\{
r\in\{0,\ldots,n-1\}:
\left\lfloor\frac{bn}{4}\right\rfloor
\le r<
\left\lfloor\frac{(b+1)n}{4}\right\rfloor
\right\},
\qquad
m_b=
\left\lfloor\frac{(b+1)n}{4}\right\rfloor
-
\left\lfloor\frac{bn}{4}\right\rfloor.
\]
The sample channel \(g(i)\) assigned to coordinate \(i\) is
\[
g(i)=
\begin{cases}
\mathrm T,&i=0,\\
\mathrm S,&i\in\{1,\ldots,6\}.
\end{cases}
\]
The covariate \(X_r^{g(i)}\) is the observed covariate of zero-based record \(r\) in that channel.

For zero-based source record \(r\), let \(Z_r\in\{0,1\}\) denote encouragement, \(D_r\in\{0,1\}\) treatment receipt, and \(Y_r\in\mathbb R\) the observed outcome. The observable marks \(W_{i,r}\) are
\[
\begin{aligned}
W_{0,r}&=1,\\
W_{1,r}&=\mathbf 1\{Z_r=0\},
&
W_{2,r}&=\mathbf 1\{Z_r=1\},\\
W_{3,r}&=\mathbf 1\{Z_r=0\}\max\{0,\min\{1,Y_r\}\},
&
W_{4,r}&=\mathbf 1\{Z_r=1\}\max\{0,\min\{1,Y_r\}\},\\
W_{5,r}&=\mathbf 1\{Z_r=0\}\mathbf 1\{D_r=1\},
&
W_{6,r}&=\mathbf 1\{Z_r=1\}\mathbf 1\{D_r=1\}.
\end{aligned}
\]
The constant mark \(W_{0,r}\) accompanies zero-based target record \(r\); each remaining mark accompanies zero-based source record \(r\). In every block sum, the mark and covariate use the same record index.
\end{definitionv}

\begin{definitionv}{synth_4}[Dyadic pilot resolution rule]
The pilot resolution \(K_0(n)\), for sample size \(n\in\mathbb N\), is defined by
\[
K_0(n)=
\begin{cases}
1, & n<n_0,\\[2pt]
2^{\left\lfloor \frac45\,\frac{\log m_0}{\log 2}\right\rfloor},
& n\ge n_0,
\end{cases}
\]
where \(n_0\) is the threshold for using the cubic estimator and \(m_0\) is the size of block \(0\) at sample size \(n\).
\end{definitionv}

\begin{definitionv}{synth_8}[Clipped marked-density pilot vector]
The pilot vector \(\widehat F(x)\), indexed by \(i\in\{0,\ldots,6\}\), is defined for real envelope parameters \(c_f,C_f\), a sample-size index \(n\in\mathbb N\), source and target samples \((\mathcal S_n,\mathcal T_n)\) each of size \(n\), and \(x\in\mathbb R\). For \(n<256\), set
\[
\widehat F(x)=\left(1,\frac12,\frac12,0,0,0,0\right).
\]
For \(n\ge256\), use \(K_0(n)\), the dyadic pilot resolution with exponent \(4/5\), and the block-\(0\) marked histogram \(H_{i,K_0(n)}^{(0)}(x)\) of \cref{obj:synth_7}, with the observable marks and zero-based record indexing of \cref{obj:synth_14}. Define \(\widehat F_i(x)\) by clipping this histogram value with the following lower and upper cutoffs, respectively:
\[
\begin{cases}
(c_f,C_f), & i=0,\\
(c_f/4,3C_f/4), & i\in\{1,2\},\\
(0,3C_f/4), & i\in\{3,4,5,6\}.
\end{cases}
\]
\end{definitionv}

\begin{definitionv}{synth_9}[Marked histogram residual]
The residual \(R_{i,K}^{(b)}(x)\) is defined by
\[
R_{i,K}^{(b)}(x)
=
H_{i,K}^{(b)}(x)-\widehat F_i(x),
\]
where the marked histogram \(H_{i,K}^{(b)}\) and the clipped pilot vector \(\widehat F\) are defined in \cref{obj:synth_7,obj:synth_8}, respectively. The parameters are arbitrary nonnegative integers \(n\) and \(K\), data consisting of \(n\) source records and \(n\) target covariates, a coordinate index \(i\) among the seven coordinates, a block index \(b\in\{0,1,2,3\}\), an evaluation point \(x\in\mathbb R\), and pilot-clipping parameters \(c_f,C_f\in\mathbb R\).
\end{definitionv}

\begin{definitionv}{synth_6}[Cubic correction resolution]
The cubic resolution \(K_3(n)\), for \(n\in\mathbb N\), is defined by
\[
K_3(n)=
\begin{cases}
1, & n<n_0,\\[2pt]
2^{\left\lfloor \dfrac{\log m_0}{\log 2}\right\rfloor}, & n\ge n_0,
\end{cases}
\]
where \(n_0\) is the threshold for using the cubic estimator and \(m_0\) is the size of block \(0\) at sample size \(n\). The floor is taken in the nonnegative integers, with value \(0\) when its argument is negative; the logarithm at \(0\) is assigned value \(0\).
\end{definitionv}

\begin{definitionv}{synth_5}[Quadratic correction resolution]
The quadratic resolution \(K_2(n)\), for \(n\in\mathbb N\), is defined by
\[
K_2(n)=
\begin{cases}
1, & n<n_0,\\[2pt]
2^{\left\lfloor \frac{4}{3}\,\frac{\log m_0}{\log 2}\right\rfloor},
& n\ge n_0,
\end{cases}
\]
where \(n_0\) is the threshold and \(m_0\) is the size of block \(0\) at sample size \(n\). The floor is taken in the nonnegative integers, with value \(0\) when its argument is negative.
\end{definitionv}

\begin{algorithmv}{def:cubic-estimator}[Cubic estimator \(\widehat T_A\)]
For \(A\in\{Y,D\}\), the inputs are the sample size \(n\), the threshold \(n_0\), the deterministic sample-channel blocks \(\mathcal B_b^G\) of sizes \(m_b\), the dyadic resolution maps \(K_0(n),K_2(n),K_3(n)\), the normalized marked histograms \(H_{i,K}^{(b)}\), the clipped pilot vector \(\widehat F\), the histogram residuals \(R_{i,K}^{(b)}\), the seven-coordinate index set \(\mathcal I\), the coordinate channels \(g(i)\), the observable marks \(W_{i,r}\), and the observed covariates \(X_r^{g(i)}\).
\begin{enumerate}
\item For \(n\ge n_0\), set the resolutions and cell midpoints to
\[
K_0=K_0(n),\qquad K_2=K_2(n),\qquad K_3=K_3(n),
\qquad
x_{\ell,K}=\frac{\ell-\tfrac12}{K},
\quad \ell=1,\ldots,K.
\]
The resolutions satisfy \(K_0\mid K_3\) and \(K_3\mid K_2\).

\item For \(n\ge n_0\), form the observable linear correction
\[
\mathsf L_A=
\sum_{i\in\mathcal I}
\left\{
\frac{1}{m_1}
\sum_{r\in\mathcal B_1^{g(i)}}
W_{i,r}\,\partial_i\Phi_A
\bigl(\widehat F(X_r^{g(i)})\bigr)
-
\int_0^1
\partial_i\Phi_A(\widehat F(x))\widehat F_i(x)\,dx
\right\}.
\]

\item For \(n\ge n_0\), form the observable quadratic correction
\[
\mathsf Q_A=
\frac{1}{2K_2}
\sum_{i,j\in\mathcal I}\sum_{\ell=1}^{K_2}
\partial_{ij}\Phi_A(\widehat F(x_{\ell,K_2}))
R_{i,K_2}^{(1)}(x_{\ell,K_2})
R_{j,K_2}^{(2)}(x_{\ell,K_2}).
\]

\item For \(n\ge n_0\), form the observable cubic correction
\[
\mathsf C_A=
\frac{1}{6K_3}
\sum_{i,j,k\in\mathcal I}\sum_{\ell=1}^{K_3}
\partial_{ijk}\Phi_A(\widehat F(x_{\ell,K_3}))
R_{i,K_3}^{(1)}(x_{\ell,K_3})
R_{j,K_3}^{(2)}(x_{\ell,K_3})
R_{k,K_3}^{(3)}(x_{\ell,K_3}).
\]
Repeated coordinate indices use distinct evaluation blocks.

\item Output
\[
\widehat T_A=
\begin{cases}
\displaystyle
\int_0^1\Phi_A(\widehat F(x))\,dx+
\mathsf L_A+\mathsf Q_A+\mathsf C_A,
& n\ge n_0,\\
0,& n<n_0.
\end{cases}
\]
Every integral is an exact finite sum, with \(K_0\mid K_3\) and \(K_3\mid K_2\). For \(n<n_0\), the score interval takes the finite-size fallback
\[
I_n=\Theta.
\]
\end{enumerate}
\end{algorithmv}

The estimator is an exact finite-sum cubic expansion around the clipped pilot. Distinct evaluation blocks supply repeated residual factors, while the nested dyadic partitions make the pilot-based derivative weights constant on correction cells.

\subsection{Observable reduction and mean-square control}

The following reduction identifies each marked coordinate with an observable measure and supplies the derivative envelope used in the Taylor remainder.

\begin{lemmav}{lem:observable-marked-reduction}[Observable marked density reduction]
Let \(n\in\mathbb N\) and let \(P\) be a transport law. Suppose:
\begin{itemize}
\item \textbf{(Lower density envelope.)} \(0<c_f<1\).
\item \textbf{(Upper density envelope.)} \(1<C_f\).
\item \textbf{(Hölder radius.)} \(1<L\).
\item \textbf{(Model membership.)} \(P\in\mathcal M_n\) for these constants, as specified in \cref{obj:def:model-class}.
\end{itemize}
Let \(F\), the seven-coordinate marked density vector, be as in \cref{obj:def:marked-density-vector}, with coordinates indexed by \(0,\ldots,6\). Write \(P_{\mathrm T}\) for the target covariate law and \(P_{\mathrm S}\) for the observed source-record law, whose record consists of covariate \(X\), encouragement \(Z\), receipt \(D\), and outcome \(Y\).

For every measurable \(B\subseteq[0,1]\) and every \(i\in\{0,\ldots,6\}\),
\[
\int_B F_i(x)\,dx
=
\begin{cases}
P_{\mathrm T}(B), & i=0,\\
\displaystyle\int_{\{X\in B\}}(1-Z)\,dP_{\mathrm S}, & i=1,\\
\displaystyle\int_{\{X\in B\}}Z\,dP_{\mathrm S}, & i=2,\\
\displaystyle\int_{\{X\in B\}}(1-Z)Y\,dP_{\mathrm S}, & i=3,\\
\displaystyle\int_{\{X\in B\}}ZY\,dP_{\mathrm S}, & i=4,\\
\displaystyle\int_{\{X\in B\}}(1-Z)D\,dP_{\mathrm S}, & i=5,\\
\displaystyle\int_{\{X\in B\}}ZD\,dP_{\mathrm S}, & i=6.
\end{cases}
\]
Moreover, the source arm densities \(q_0\) and \(q_1\) satisfy
\[
q_0(x)=F_1(x)\ge \frac{c_f}{4},
\qquad
q_1(x)=F_2(x)\ge \frac{c_f}{4},
\qquad x\in[0,1].
\]
For each \(A\in\{Y,D\}\), the transported contrast \(T_A\) and the integrand \(\Phi_A\) of \cref{obj:def:smooth-functional} satisfy
\[
T_A=\int_0^1\Phi_A(F(x))\,dx.
\]

For every \(A\in\{Y,D\}\) and every \(v\in\mathbb R^7\) in the deterministic clipping rectangle
\[
c_f\le v_0\le C_f,\qquad
\frac{c_f}{4}\le v_1,v_2\le\frac{3C_f}{4},\qquad
0\le v_i\le\frac{3C_f}{4}\quad(i=3,\ldots,6),
\]
every order \(k\in\{0,\ldots,4\}\) and every choice of coordinate indices \(i_1,\ldots,i_k\in\{0,\ldots,6\}\) obey
\[
\left|
\frac{\partial^k\Phi_A}{\partial v_{i_1}\cdots\partial v_{i_k}}(v)
\right|
\le
48(1+C_f)^2\left(\frac{4}{c_f}\right)^5.
\]
At order zero, the expression on the left is \(\lvert\Phi_A(v)\rvert\).
\end{lemmav}

\begin{proof}[Proof of \cref{obj:lem:observable-marked-reduction}]
We establish the observable identities, the lower bounds, the functional representation, and then the derivative bounds.

\begin{enumerate}
\item For the target coordinate, \cref{obj:def:marked-density-vector} gives \(F_0=f_{\mathrm T}\). By \cref{obj:ass:target-density-holder}, this density is continuous on \([0,1]\), hence measurable there. For every measurable \(B\subseteq[0,1]\), \cref{obj:ass:target-density-bounds} therefore yields
\[
P_{\mathrm T}(B)
=\int_B\max\{f_{\mathrm T}(x),0\}\,dx
=\int_B f_{\mathrm T}(x)\,dx
=\int_B F_0(x)\,dx.
\]
The second equality uses \(f_{\mathrm T}(x)\ge c_f>0\) on \(B\).
% lean: CausalSmith.Stat.TransportCaceRoughnuisanceLength.observable_marked_reduction

\item For the source coordinates, define the encouragement-arm probabilities by
\[
\pi_z(x)=
\begin{cases}
1-e(x),&z=0,\\
e(x),&z=1,
\end{cases}
\qquad x\in[0,1],\quad z\in\{0,1\}.
\]
By the definition of the source propensity,
\[
P_{\mathrm S}(Z=z\mid X)=\pi_z(X)
\qquad P_{\mathrm S}\text{-almost surely}.
\]
Conditioning on \(X\) under the observed source law therefore gives
\[
P_{\mathrm S}(X\in B,Z=z)
=
\mathbb E_{P_{\mathrm S}}\!\left[
\mathbf 1_{\{X\in B\}}\pi_z(X)
\right]
=
\int_B f_{\mathrm S}(x)\pi_z(x)\,dx.
\]
Here the last equality uses the source covariate density identity in
\cref{obj:ass:source-density-bounds}; measurability follows from the stated regularity assumptions. Integrating an arm indicator also gives
\[
\int_{\{X\in B\}}\mathbf 1_{\{Z=z\}}\,dP_{\mathrm S}
=P_{\mathrm S}(X\in B,Z=z).
\]
Since \(Z\) is binary, \(\mathbf 1_{\{Z=0\}}=1-Z\) and
\(\mathbf 1_{\{Z=1\}}=Z\). Thus the two assignment coordinates satisfy
\[
\int_B F_1(x)\,dx
=\int_B f_{\mathrm S}(x)(1-e(x))\,dx
=\int_{\{X\in B\}}(1-Z)\,dP_{\mathrm S},
\]
\[
\int_B F_2(x)\,dx
=\int_B f_{\mathrm S}(x)e(x)\,dx
=\int_{\{X\in B\}}Z\,dP_{\mathrm S}.
\]

For either response \(A\in\{Y,D\}\), restriction by the same indicator and the source arm-moment identity in \cref{obj:ass:arm-means-holder} give
\[
\begin{aligned}
\int_{\{X\in B\}}\mathbf 1_{\{Z=z\}}A\,dP_{\mathrm S}
&=\int_{\{X\in B,\;Z=z\}}A\,dP_{\mathrm S}\\
&=\int_B f_{\mathrm S}(x)\pi_z(x)m_{Az}(x)\,dx
=\int_B r_{Az}(x)\,dx.
\end{aligned}
\]
Applying this identity first to \(Y\) in arms \(0,1\), and then to \(D\) in arms \(0,1\), yields respectively
\[
\begin{aligned}
\int_B F_3(x)\,dx&=\int_{\{X\in B\}}(1-Z)Y\,dP_{\mathrm S},&
\int_B F_4(x)\,dx&=\int_{\{X\in B\}}ZY\,dP_{\mathrm S},\\
\int_B F_5(x)\,dx&=\int_{\{X\in B\}}(1-Z)D\,dP_{\mathrm S},&
\int_B F_6(x)\,dx&=\int_{\{X\in B\}}ZD\,dP_{\mathrm S}.
\end{aligned}
\]
These arguments apply also when \(B\) is empty.
% lean: CausalSmith.Stat.TransportCaceRoughnuisanceLength.source_assignment_density
% lean: CausalSmith.Stat.TransportCaceRoughnuisanceLength.source_assignment_mark_integral
% lean: CausalSmith.Stat.TransportCaceRoughnuisanceLength.source_marked_arm_density

\item Fix \(x\in[0,1]\). By
\cref{obj:ass:source-density-bounds,obj:ass:propensity-overlap},
\(f_{\mathrm S}(x)\ge c_f>0\) and both \(e(x)\) and \(1-e(x)\) are at least \(1/4\). Consequently,
\[
\frac{c_f}{4}
=c_f\frac14
\le f_{\mathrm S}(x)\frac14
\le f_{\mathrm S}(x)(1-e(x))
=q_0(x)=F_1(x),
\]
\[
\frac{c_f}{4}
=c_f\frac14
\le f_{\mathrm S}(x)\frac14
\le f_{\mathrm S}(x)e(x)
=q_1(x)=F_2(x).
\]
% lean: CausalSmith.Stat.TransportCaceRoughnuisanceLength.markedDensityVector_assignment_lower_bounds

\item The preceding bounds make both arm densities strictly positive. The identities in \cref{obj:def:marked-density-vector} therefore permit cancellation:
\[
\frac{r_{Az}(x)}{q_z(x)}=m_{Az}(x),
\qquad A\in\{Y,D\},\quad z\in\{0,1\},\quad x\in[0,1].
\]
Substituting into the integrand of \cref{obj:def:smooth-functional} gives, for each response,
\[
\Phi_A(F(x))
=f_{\mathrm T}(x)\bigl(m_{A1}(x)-m_{A0}(x)\bigr).
\]
Integrating this pointwise equality proves
\[
T_A
=\int_{[0,1]}f_{\mathrm T}(x)
       \bigl(m_{A1}(x)-m_{A0}(x)\bigr)\,dx
=\int_0^1\Phi_A(F(x))\,dx.
\]
The closed-interval set integral and the displayed interval integral agree because their endpoint difference has Lebesgue measure zero.
% lean: CausalSmith.Stat.TransportCaceRoughnuisanceLength.Phi_markedDensityVector_eq_transport_integrand
% lean: CausalSmith.Stat.TransportCaceRoughnuisanceLength.transportedForm_eq_integral_Phi_markedDensityVector

\item Fix \(A\in\{Y,D\}\) and a vector \(v\) in the clipping rectangle. Define the reciprocal scale and the desired envelope by
\[
R_{\mathrm{inv}}=\frac4{c_f},
\qquad
M_{\mathrm{der}}=48(1+C_f)^2R_{\mathrm{inv}}^5.
\]
The hypotheses imply \(R_{\mathrm{inv}}>1\) and \(1+C_f>0\). To identify the response coordinates, set
\[
j_{Az}=
\begin{cases}
3+z,&A=Y,\\
5+z,&A=D,
\end{cases}
\qquad A\in\{Y,D\},\quad z\in\{0,1\}.
\]
On the open set \(\{v\in\mathbb R^7:v_1>0,\ v_2>0\}\), define the two arm contributions by
\[
H_{Az}(v)=\frac{v_0v_{j_{Az}}}{v_{1+z}},
\qquad
\Phi_A(v)=H_{A1}(v)-H_{A0}(v).
\]
The clipping rectangle lies in this open set.

For order zero, the rectangle supplies
\[
0\le v_0\le C_f,\qquad
0\le v_{j_{Az}}\le\frac{3C_f}{4},\qquad
v_{1+z}\ge\frac{c_f}{4}>0.
\]
In particular,
\[
v_0v_{j_{Az}}\le(1+C_f)^2,
\qquad
1=\frac{c_f}{4}R_{\mathrm{inv}}
\le v_{1+z}R_{\mathrm{inv}},
\]
and hence
\[
0\le H_{Az}(v)\le(1+C_f)^2R_{\mathrm{inv}}.
\]
Since
\[
R_{\mathrm{inv}}
=R_{\mathrm{inv}}\cdot1
\le R_{\mathrm{inv}}R_{\mathrm{inv}}^4
=R_{\mathrm{inv}}^5,
\]
we obtain
\[
|\Phi_A(v)|
\le 2(1+C_f)^2R_{\mathrm{inv}}
\le48(1+C_f)^2R_{\mathrm{inv}}^5
=M_{\mathrm{der}}.
\]
This treats the empty list of coordinate indices at order zero.
% lean: CausalSmith.Stat.TransportCaceRoughnuisanceLength.Phi_value_abs_le_fourthDerivativeEnvelope

\item For the positive orders, fix arbitrary coordinate indices
\(i_1,\ldots,i_4\in\{0,\ldots,6\}\), allowing repetitions. Every shorter list can be extended to such a list. Define its coordinate directions by
\[
\eta^{(\ell)}_j=
\begin{cases}
1,&j=i_\ell,\\
0,&j\ne i_\ell,
\end{cases}
\qquad 1\le\ell\le4,\quad 0\le j\le6.
\]
Thus \(|\eta^{(\ell)}_j|\le1\). For a fixed arm \(A,z\), define
\[
\mathsf U_\ell
=\eta^{(\ell)}_0v_{j_{Az}}
 +v_0\eta^{(\ell)}_{j_{Az}},
\qquad 1\le\ell\le4.
\]
The rectangle gives
\[
|v_0|\le C_f,\qquad |v_{j_{Az}}|\le C_f,\qquad
0<v_{1+z},\qquad
|v_{1+z}^{-1}|
=v_{1+z}^{-1}
\le\left(\frac{c_f}{4}\right)^{-1}
=R_{\mathrm{inv}}.
\]
The product and reciprocal rules yield
\[
\partial_{i_1}H_{Az}(v)
=\mathsf U_1v_{1+z}^{-1}
-v_0v_{j_{Az}}\eta^{(1)}_{1+z}v_{1+z}^{-2}.
\]
Moreover,
\[
|\mathsf U_1|
\le|\eta^{(1)}_0|\,|v_{j_{Az}}|
   +|v_0|\,|\eta^{(1)}_{j_{Az}}|
\le2C_f.
\]
The two terms therefore obey
\[
|\mathsf U_1v_{1+z}^{-1}|\le2C_fR_{\mathrm{inv}},
\qquad
|v_0v_{j_{Az}}\eta^{(1)}_{1+z}v_{1+z}^{-2}|
\le C_f^2R_{\mathrm{inv}}^2.
\]
Using
\[
R_{\mathrm{inv}}\le R_{\mathrm{inv}}^2,\qquad
C_f\le(1+C_f)^2,\qquad
C_f^2\le(1+C_f)^2,
\]
we obtain the one-arm bound
\[
|\partial_{i_1}H_{Az}(v)|
\le2C_fR_{\mathrm{inv}}+C_f^2R_{\mathrm{inv}}^2
\le3(1+C_f)^2R_{\mathrm{inv}}^2.
\]
Both arm functions are smooth on the stated open set, so differentiation of their difference and \(R_{\mathrm{inv}}^2\le R_{\mathrm{inv}}^5\) give
\[
|\partial_{i_1}\Phi_A(v)|
\le|\partial_{i_1}H_{A1}(v)|
  +|\partial_{i_1}H_{A0}(v)|
\le6(1+C_f)^2R_{\mathrm{inv}}^2
\le M_{\mathrm{der}}.
\]
% lean: CausalSmith.Stat.TransportCaceRoughnuisanceLength.iteratedFDeriv_one_Phi_apply
% lean: CausalSmith.Stat.TransportCaceRoughnuisanceLength.marked_arm_first_derivative_abs_le
% lean: CausalSmith.Stat.TransportCaceRoughnuisanceLength.Phi_first_derivative_abs_le_fourthDerivativeEnvelope

\item For order two, define the paired direction coefficients by
\[
\mathsf V_{\ell,\ell'}
=\eta^{(\ell)}_0\eta^{(\ell')}_{j_{Az}}
 +\eta^{(\ell')}_0\eta^{(\ell)}_{j_{Az}},
\qquad 1\le\ell,\ell'\le4.
\]
Their bounds, and those of the linear coefficients, are
\[
|\mathsf V_{\ell,\ell'}|
\le1\cdot1+1\cdot1=2,
\qquad
|\mathsf U_\ell|\le2C_f.
\]
As functions of \(v\), these coefficients satisfy
\[
\partial_{i_{\ell'}}\mathsf U_\ell
=\mathsf V_{\ell,\ell'},
\qquad
\partial_{i_{\ell'}}\mathsf V_{\ell,\nu}=0,
\qquad
\partial_{i_\ell}v_{1+z}^{-p}
=-p\,\eta^{(\ell)}_{1+z}v_{1+z}^{-(p+1)}
\quad(p\ge1).
\]
Differentiating the first-order expression thus gives
\[
\begin{aligned}
\partial_{i_1}\partial_{i_2}H_{Az}(v)
={}&\mathsf V_{1,2}v_{1+z}^{-1}
-\mathsf U_1\eta^{(2)}_{1+z}v_{1+z}^{-2}
-\mathsf U_2\eta^{(1)}_{1+z}v_{1+z}^{-2}\\
&+2v_0v_{j_{Az}}\eta^{(1)}_{1+z}
                      \eta^{(2)}_{1+z}v_{1+z}^{-3}.
\end{aligned}
\]
The four terms have absolute values bounded respectively by
\[
2R_{\mathrm{inv}},\qquad
2C_fR_{\mathrm{inv}}^2,\qquad
2C_fR_{\mathrm{inv}}^2,\qquad
2C_f^2R_{\mathrm{inv}}^3.
\]
Consequently,
\[
\begin{aligned}
|\partial_{i_1}\partial_{i_2}H_{Az}(v)|
&\le2R_{\mathrm{inv}}
    +4C_fR_{\mathrm{inv}}^2
    +2C_f^2R_{\mathrm{inv}}^3\\
&\le(2+4+2)(1+C_f)^2R_{\mathrm{inv}}^5\\
&\le24(1+C_f)^2R_{\mathrm{inv}}^5.
\end{aligned}
\]
Here each reciprocal power is at most \(R_{\mathrm{inv}}^5\), and each of
\(1,C_f,C_f^2\) is at most \((1+C_f)^2\). Subtracting the two arm derivatives gives
\[
|\partial_{i_1}\partial_{i_2}\Phi_A(v)|
\le|\partial_{i_1}\partial_{i_2}H_{A1}(v)|
  +|\partial_{i_1}\partial_{i_2}H_{A0}(v)|
\le M_{\mathrm{der}}.
\]
% lean: CausalSmith.Stat.TransportCaceRoughnuisanceLength.cubicRatioD2_basis_abs_le
% lean: CausalSmith.Stat.TransportCaceRoughnuisanceLength.dPhi2_abs_le

\item For order three, the rectangle also supplies the nonnegative numerator bounds
\[
0\le v_0\le C_f,\qquad
0\le v_{j_{Az}}\le C_f,
\qquad
0<v_{1+z},\qquad
0<v_{1+z}^{-1}\le R_{\mathrm{inv}}.
\]
Differentiating the second-order expression by the same product and reciprocal rules yields
\[
\begin{aligned}
\partial_{i_1}\partial_{i_2}\partial_{i_3}H_{Az}(v)
={}&-\mathsf V_{1,2}\eta^{(3)}_{1+z}v_{1+z}^{-2}\\
&-\bigl(\mathsf V_{1,3}\eta^{(2)}_{1+z}
       +\mathsf V_{2,3}\eta^{(1)}_{1+z}\bigr)v_{1+z}^{-2}\\
&+2\mathsf U_1\eta^{(2)}_{1+z}\eta^{(3)}_{1+z}v_{1+z}^{-3}\\
&+2\mathsf U_2\eta^{(1)}_{1+z}\eta^{(3)}_{1+z}v_{1+z}^{-3}\\
&+2\mathsf U_3\eta^{(1)}_{1+z}\eta^{(2)}_{1+z}v_{1+z}^{-3}\\
&-6v_0v_{j_{Az}}\eta^{(1)}_{1+z}
                    \eta^{(2)}_{1+z}\eta^{(3)}_{1+z}v_{1+z}^{-4}.
\end{aligned}
\]
The first paired coefficient has absolute value at most \(2\), and the grouped coefficient satisfies
\[
\left|\mathsf V_{1,3}\eta^{(2)}_{1+z}
       +\mathsf V_{2,3}\eta^{(1)}_{1+z}\right|
\le2\cdot1+2\cdot1=4.
\]
Thus the six displayed terms have absolute values bounded respectively by
\[
2R_{\mathrm{inv}}^2,\qquad
4R_{\mathrm{inv}}^2,\qquad
4C_fR_{\mathrm{inv}}^3,\qquad
4C_fR_{\mathrm{inv}}^3,\qquad
4C_fR_{\mathrm{inv}}^3,\qquad
6C_f^2R_{\mathrm{inv}}^4.
\]
Adding these bounds gives
\[
\begin{aligned}
|\partial_{i_1}\partial_{i_2}\partial_{i_3}H_{Az}(v)|
&\le6R_{\mathrm{inv}}^2
    +12C_fR_{\mathrm{inv}}^3
    +6C_f^2R_{\mathrm{inv}}^4\\
&\le(6+12+6)(1+C_f)^2R_{\mathrm{inv}}^5\\
&=24(1+C_f)^2R_{\mathrm{inv}}^5.
\end{aligned}
\]
The second inequality uses \(R_{\mathrm{inv}}^2,R_{\mathrm{inv}}^3,
R_{\mathrm{inv}}^4\le R_{\mathrm{inv}}^5\) and the same numerator bounds as before. For \(A=D\), the two arm terms use the coordinate triples \((0,6,2)\) and \((0,5,1)\); for \(A=Y\), they use \((0,4,2)\) and \((0,3,1)\). In either case,
\[
|\partial_{i_1}\partial_{i_2}\partial_{i_3}\Phi_A(v)|
\le|\partial_{i_1}\partial_{i_2}\partial_{i_3}H_{A1}(v)|
  +|\partial_{i_1}\partial_{i_2}\partial_{i_3}H_{A0}(v)|
\le M_{\mathrm{der}}.
\]
% lean: CausalSmith.Stat.TransportCaceRoughnuisanceLength.cubicRatioD3_basis_abs_le
% lean: CausalSmith.Stat.TransportCaceRoughnuisanceLength.dPhi3_abs_le

\item The remaining order is four. Differentiating the third-order expression gives
\[
\begin{aligned}
\partial_{i_1}\partial_{i_2}\partial_{i_3}\partial_{i_4}H_{Az}(v)
={}&
2v_{1+z}^{-3}
\sum_{1\le\ell<\ell'\le4}
\mathsf V_{\ell,\ell'}
\prod_{\substack{1\le\nu\le4\\\nu\ne\ell,\ell'}}
\eta^{(\nu)}_{1+z}\\
&-6v_{1+z}^{-4}
\sum_{\ell=1}^4
\mathsf U_\ell
\prod_{\substack{1\le\nu\le4\\\nu\ne\ell}}
\eta^{(\nu)}_{1+z}\\
&+24v_0v_{j_{Az}}v_{1+z}^{-5}
\prod_{\nu=1}^4\eta^{(\nu)}_{1+z}.
\end{aligned}
\]
There are six pairs in the first sum. Each corresponding term satisfies
\[
\left|
2\mathsf V_{\ell,\ell'}
\prod_{\substack{1\le\nu\le4\\\nu\ne\ell,\ell'}}
\eta^{(\nu)}_{1+z}v_{1+z}^{-3}
\right|
\le2\cdot2\cdot R_{\mathrm{inv}}^3
=4R_{\mathrm{inv}}^3.
\]
Each of the four terms in the second sum satisfies
\[
\left|
6\mathsf U_\ell
\prod_{\substack{1\le\nu\le4\\\nu\ne\ell}}
\eta^{(\nu)}_{1+z}v_{1+z}^{-4}
\right|
\le6\cdot2C_fR_{\mathrm{inv}}^4
=12C_fR_{\mathrm{inv}}^4.
\]
The final term is bounded by
\[
\left|
24v_0v_{j_{Az}}v_{1+z}^{-5}
\prod_{\nu=1}^4\eta^{(\nu)}_{1+z}
\right|
\le24C_f^2R_{\mathrm{inv}}^5.
\]
Therefore,
\[
\begin{aligned}
|\partial_{i_1}\partial_{i_2}\partial_{i_3}\partial_{i_4}H_{Az}(v)|
&\le24R_{\mathrm{inv}}^3
    +48C_fR_{\mathrm{inv}}^4
    +24C_f^2R_{\mathrm{inv}}^5\\
&\le24R_{\mathrm{inv}}^5
    +48C_fR_{\mathrm{inv}}^5
    +24C_f^2R_{\mathrm{inv}}^5\\
&=24(1+C_f)^2R_{\mathrm{inv}}^5.
\end{aligned}
\]
The middle inequality uses \(R_{\mathrm{inv}}>1\). Subtracting the two arm derivatives finally gives
\[
\begin{aligned}
|\partial_{i_1}\partial_{i_2}\partial_{i_3}\partial_{i_4}\Phi_A(v)|
&\le
|\partial_{i_1}\partial_{i_2}\partial_{i_3}\partial_{i_4}H_{A1}(v)|
+
|\partial_{i_1}\partial_{i_2}\partial_{i_3}\partial_{i_4}H_{A0}(v)|\\
&\le48(1+C_f)^2R_{\mathrm{inv}}^5
=48(1+C_f)^2\left(\frac4{c_f}\right)^5.
\end{aligned}
\]
Together with the preceding cases, this proves the asserted bound for every order \(k\in\{0,\ldots,4\}\) and every choice of coordinate indices.
% lean: CausalSmith.Stat.TransportCaceRoughnuisanceLength.cubicRatioD4_basis_abs_le
% lean: CausalSmith.Stat.TransportCaceRoughnuisanceLength.dPhi4_abs_le
\end{enumerate}
\end{proof}

The measure identities justify the normalized marked counts, and the lower arm-density bounds keep every functional denominator uniformly positive on the clipping rectangle. The next result combines pilot approximation, projection cancellation, and block-conditional variance calculations.

\begin{lemmav}{lem:marked-cubic-mse}[Uniform cubic mean-square bound]
Let \(c_f\) and \(C_f\) be the lower and upper density envelopes, and let \(L\) be the common Hölder radius. Suppose that:
\begin{itemize}
\item \textbf{(Envelope constants.)} \(0<c_f<1\), \(1<C_f\), and \(1<L\).
\item \textbf{(Sample size.)} The common source and target sample size \(n\) satisfies
\[
n\ge n_0=256.
\]
\item \textbf{(Model.)} \(P\in\mathcal M_n\), with \(\mathcal M_n\) defined in \cref{obj:def:model-class}.
\end{itemize}
For each \(A\in\{Y,D\}\), let \(\widehat T_A\) be the statistic in \cref{obj:def:cubic-estimator}, and let \(T_A\) be the target-averaged source contrast,
\[
T_A
=
\int_{\text{covariate space}}
f_{\mathrm T}(x)\bigl(m_{A1}(x)-m_{A0}(x)\bigr)\,dx.
\]
Then, under the equal-size source–target sampling law specified by \(P\),
\[
\mathbb E_P\!\left[(\widehat T_A-T_A)^2\right]
\le C_{\mathrm{mse}}\,n^{-2/3},
\]
where the finite, explicitly computable constant \(C_{\mathrm{mse}}=C_{\mathrm{mse}}(c_f,C_f,L)\) is
\[
\begin{aligned}
C_{\mathrm{mse}}
={}&3\Bigg\{
\left[2(1+C_f)^2(4/c_f)^5\right]^2 7^8
\Big[
2^7\,8!\,
\big\{(1+C_f+C_f^2)^4 5^{4/5}
 +(1+C_f+C_f^2)5^{7/5}\big\}\\
&\hspace{48mm}
+2^8[3(1+C_f)L]^8 5^{4/5}
\Big]\\
&\quad
+7^4[48(1+C_f)^2(4/c_f)^5]^2
 [3(1+C_f)L]^4 2^{1/2}5^{2/3}\\
&\quad
+2\Big[
\big\{24(1+C_f)^2(4/c_f)^5
       7^2[3(1+C_f)L]^2\big\}^2
 7^2\\
&\hspace{24mm}\cdot
\big\{2(1+C_f+C_f^2)5^{1/5}
 +2^{5/4}[3(1+C_f)L]^2 5^{1/5}\big\}
 2^{1/2}5^{1/2}\\
&\hspace{16mm}
+\big\{8(1+C_f)^2(4/c_f)^5
       7^3[3(1+C_f)L]^3\big\}^2
 2^{3/4}5^{3/4}
\Big]
\Bigg\}\\
&+3\Bigg\{
10\cdot7^2[48(1+C_f)^2(4/c_f)^5]^2\\
&\quad
+60\cdot7^4(2^2-1)^2
 [48(1+C_f)^2(4/c_f)^5]^2
 [1+2(1+C_f)]^4[2+C_f+C_f^2]^2\\
&\quad
+310\cdot7^6(2^3-1)^2
 [48(1+C_f)^2(4/c_f)^5]^2
 [1+2(1+C_f)]^6[2+C_f+C_f^2]^3
\Bigg\}.
\end{aligned}
\]
\end{lemmav}

\begin{proof}[Proof of \cref{obj:lem:marked-cubic-mse}]
Fix \(n\ge256\), \(P\in\mathcal M_n\), and \(A\in\{Y,D\}\). All expectations below are under the equal-size source--target sampling law of \(P\). Use the marked-density vector \(F\), the clipped pilot \(\widehat F\), and the corrections \(\mathsf L_A,\mathsf Q_A,\mathsf C_A\) from \cref{obj:def:marked-density-vector,obj:def:cubic-estimator}. Define the training sigma-field and four envelope constants by
\[
\begin{aligned}
\Gamma
&:=\sigma\!\left(
O_r^{\mathrm S},X_r^{\mathrm T}:
r\in\mathcal B_0^{\mathrm S}
=\mathcal B_0^{\mathrm T}
\right),\\
U_{\mathrm{env}}&:=2(1+C_f),&
H_{\mathrm{env}}&:=3(1+C_f)L,\\
V_{\mathrm{bin}}&:=1+C_f+C_f^2,&
M_{\Phi}&:=48(1+C_f)^2(4/c_f)^5.
\end{aligned}
\]
The sampling conditions in
\cref{obj:ass:source-iid,obj:ass:target-iid,obj:ass:sample-independence}
make the three evaluation blocks independent of the training records and mutually independent.

\begin{enumerate}
\item \textit{Envelope and resolution bounds.}
The density, overlap, and arm-mean bounds imply that \(F(x)\) belongs to the clipping rectangle of \cref{obj:lem:observable-marked-reduction} for every \(x\in[0,1]\). The pilot belongs to the same rectangle by \cref{obj:synth_8}. Consequently,
\[
0\le F_i(x),\widehat F_i(x)\le C_f,
\qquad
|F_i(x)-\widehat F_i(x)|\le2C_f
\quad(i\in\mathcal I,\ x\in[0,1]).
\]
Moreover, \cref{obj:lem:observable-marked-reduction} supplies the functional representation and derivative bounds
\[
T_A=\int_0^1\Phi_A(F(x))\,dx,
\qquad
\bigl|\partial_{i_1\cdots i_s}\Phi_A(\widehat F(x))\bigr|
\le M_{\Phi},
\quad 0\le s\le4.
\]
The product seminorm inequality, applied first to the arm densities and then to their marked versions, gives
\[
[q_z]_\beta\le(1+C_f)L,\qquad
[r_{Az}]_\beta
\le C_fL+(1+C_f)L
\le H_{\mathrm{env}},
\]
while \([f_{\mathrm T}]_\beta\le L\le H_{\mathrm{env}}\). Thus
\[
|F_i(x)-F_i(y)|
\le H_{\mathrm{env}}|x-y|^{1/8}
\quad(i\in\mathcal I,\ x,y\in[0,1]).
\]
% lean: CausalSmith.Stat.TransportCaceRoughnuisanceLength.markedDensityVector_holder_modulus

The exact block sizes and the dyadic rules of
\cref{obj:synth_3,obj:synth_4,obj:synth_5,obj:synth_6}
give
\[
m_b\ge\lfloor n/4\rfloor\ge n/4-1\ge n/5>0
\quad(b=0,1,2,3),
\]
and
\[
\frac12m_0^{4/5}\le K_0\le m_0^{4/5},\qquad
\frac12m_0^{4/3}\le K_2\le m_0^{4/3},\qquad
\frac12m_0\le K_3\le m_0.
\]
Indeed, rounding a nonnegative logarithmic exponent down changes its power of two by a factor between \(1/2\) and \(1\). In particular,
\[
K_2\le n^{4/3},\qquad K_3\le n.
\]
The nesting property in \cref{obj:def:cubic-estimator} ensures that the pilot is constant on every correction cell, including its assigned endpoints.
% lean: CausalSmith.Stat.TransportCaceRoughnuisanceLength.block_size_lower
% lean: CausalSmith.Stat.TransportCaceRoughnuisanceLength.pilotResolution_power_bounds
% lean: CausalSmith.Stat.TransportCaceRoughnuisanceLength.quadraticResolution_power_bounds
% lean: CausalSmith.Stat.TransportCaceRoughnuisanceLength.cubicResolution_le_blockSize

\item \textit{Exact projected means.}
For every positive integer \(K\), use the cells \(E_{\ell,K}\) of \cref{obj:synth_7}, with \(0\le\ell<K\), and define the cell-averaging operator on integrable functions by
\[
(\Pi_K\varphi)(x)
:=
K\sum_{\ell=0}^{K-1}\mathbf1\{x\in E_{\ell,K}\}
\int_{E_{\ell,K}}\varphi(y)\,dy,
\qquad x\in[0,1].
\]
For \(K_0\mid K\), define
\[
\begin{aligned}
\delta_i(x)&:=F_i(x)-\widehat F_i(x),\\
\delta^{\mathrm{av}}_{i,K}(x)
&:=(\Pi_KF_i)(x)-\widehat F_i(x),\\
\delta^{\mathrm{av}}_{i,K,\ell}
&:=\delta^{\mathrm{av}}_{i,K}(x_{\ell+1,K}),\\
\eta^{(b)}_{i,K,\ell}
&:=H_{i,K}^{(b)}(x_{\ell+1,K})
-(\Pi_KF_i)(x_{\ell+1,K}),
\qquad b=1,2,3.
\end{aligned}
\]
Here and below the spatial error functions have domain \([0,1]\). The cell-average shifts satisfy
\[
|\delta^{\mathrm{av}}_{i,K,\ell}|
\le2C_f\le U_{\mathrm{env}}.
\]

For \(0\le r<n\), define the observable bin scores and their means by
\[
\mathsf B_{i,K,\ell,r}
:=
W_{i,r}\mathbf1\{X_r^{g(i)}\in E_{\ell,K}\},
\qquad
\pi_{i,K,\ell}:=\int_{E_{\ell,K}}F_i(x)\,dx.
\]
The marked moment identities in \cref{obj:lem:observable-marked-reduction}, together with the boundedness in \cref{obj:ass:outcome-consistency}, give
\[
0\le\mathsf B_{i,K,\ell,r}\le1,\qquad
\mathbb E_P[\mathsf B_{i,K,\ell,r}]=\pi_{i,K,\ell},
\qquad
0\le\pi_{i,K,\ell}\le C_f/K.
\]
Thus
\[
\mathbb E_P[H_{i,K}^{(b)}(x_{\ell+1,K})\mid\Gamma]
=(\Pi_KF_i)(x_{\ell+1,K}),
\qquad
\mathbb E_P[\eta^{(b)}_{i,K,\ell}\mid\Gamma]=0.
\]

Define the three projected corrections and their total by
\[
\begin{aligned}
\overline{\mathsf L}_A
&:=\sum_{i\in\mathcal I}\int_0^1
\partial_i\Phi_A(\widehat F(x))\delta_i(x)\,dx,\\
\overline{\mathsf Q}_A
&:=\frac1{2K_2}
\sum_{i,j\in\mathcal I}\sum_{\ell=0}^{K_2-1}
\partial_{ij}\Phi_A(\widehat F(x_{\ell+1,K_2}))
\delta^{\mathrm{av}}_{i,K_2,\ell}
\delta^{\mathrm{av}}_{j,K_2,\ell},\\
\overline{\mathsf C}_A
&:=\frac1{6K_3}
\sum_{i,j,k\in\mathcal I}\sum_{\ell=0}^{K_3-1}
\partial_{ijk}\Phi_A(\widehat F(x_{\ell+1,K_3}))
\delta^{\mathrm{av}}_{i,K_3,\ell}
\delta^{\mathrm{av}}_{j,K_3,\ell}
\delta^{\mathrm{av}}_{k,K_3,\ell},\\
\overline T_A
&:=\int_0^1\Phi_A(\widehat F(x))\,dx
+\overline{\mathsf L}_A+\overline{\mathsf Q}_A
+\overline{\mathsf C}_A.
\end{aligned}
\]
These are training-measurable quantities. Since the coefficient and pilot are constant on each pilot cell, the linear spatial term is exactly
\[
\overline{\mathsf L}_A
=
\frac1{K_0}\sum_{i\in\mathcal I}\sum_{\ell=0}^{K_0-1}
\partial_i\Phi_A(\widehat F(x_{\ell+1,K_0}))
\delta^{\mathrm{av}}_{i,K_0,\ell}.
\]
The sample covariates belong to \([0,1]\) almost surely, so partitioning the empirical linear term into these cells yields
\[
\mathsf L_A-\overline{\mathsf L}_A
=
\frac1{K_0}\sum_{i\in\mathcal I}\sum_{\ell=0}^{K_0-1}
\partial_i\Phi_A(\widehat F(x_{\ell+1,K_0}))
\eta^{(1)}_{i,K_0,\ell}.
\]

At a correction-cell midpoint,
\[
R^{(b)}_{i,K}
=\delta^{\mathrm{av}}_{i,K,\ell}
+\eta^{(b)}_{i,K,\ell}.
\]
Expanding the quadratic product gives its constant product and three terms containing centered factors; expanding the cubic product gives its constant product and seven such terms. Every nonconstant term has conditional mean zero, because its centered factors belong to distinct evaluation blocks. This remains true when coordinate indices coincide. Therefore, almost surely,
\[
\mathbb E_P[\mathsf L_A\mid\Gamma]=\overline{\mathsf L}_A,\qquad
\mathbb E_P[\mathsf Q_A\mid\Gamma]=\overline{\mathsf Q}_A,\qquad
\mathbb E_P[\mathsf C_A\mid\Gamma]=\overline{\mathsf C}_A,
\]
and hence
\[
\widehat T_A-\overline T_A
=
(\mathsf L_A-\overline{\mathsf L}_A)
+(\mathsf Q_A-\overline{\mathsf Q}_A)
+(\mathsf C_A-\overline{\mathsf C}_A).
\]
The bounded coefficients and finite cell sums supply the integrability needed for these conditional identities.
% lean: CausalSmith.Stat.TransportCaceRoughnuisanceLength.linearProjectionMean_eq_cell_sum
% lean: CausalSmith.Stat.TransportCaceRoughnuisanceLength.linearTerm_sub_projectionMean_eq_centered_cell_sum
% lean: CausalSmith.Stat.TransportCaceRoughnuisanceLength.condExp_linearTerm_eq_linearProjectionMean
% lean: CausalSmith.Stat.TransportCaceRoughnuisanceLength.condExp_quadraticTerm_eq_quadraticProjectionMean
% lean: CausalSmith.Stat.TransportCaceRoughnuisanceLength.condExp_cubicTerm_eq_cubicProjectionMean
% lean: CausalSmith.Stat.TransportCaceRoughnuisanceLength.cubicEstimator_sub_projectedEstimatorMean_eq_centered_corrections

\item \textit{Linear conditional second moment.}
For each coordinate define
\[
\gamma^{\mathrm{lin}}_{i,\ell}
:=\partial_i\Phi_A(\widehat F(x_{\ell+1,K_0})),
\qquad
\mathsf Z^{\mathrm{lin}}_i
:=\frac1{K_0}\sum_{\ell=0}^{K_0-1}
\gamma^{\mathrm{lin}}_{i,\ell}\eta^{(1)}_{i,K_0,\ell}.
\]
Conditional on training, the signed covariance sum for this cell average is the variance of a single weighted record, divided by \(m_1\):
\[
\begin{aligned}
\mathbb E_P[(\mathsf Z^{\mathrm{lin}}_i)^2\mid\Gamma]
&=
\frac1{K_0^2}\sum_{\ell,r=0}^{K_0-1}
\gamma^{\mathrm{lin}}_{i,\ell}\gamma^{\mathrm{lin}}_{i,r}
\operatorname{Cov}_P\!\left(
H_{i,K_0}^{(1)}(x_{\ell+1,K_0}),
H_{i,K_0}^{(1)}(x_{r+1,K_0})
\right)\\
&=
\frac1{m_1}
\operatorname{Var}_P\!\left(
\sum_{\ell=0}^{K_0-1}
\gamma^{\mathrm{lin}}_{i,\ell}
\mathsf B_{i,K_0,\ell,r_0}
\,\middle|\,\Gamma
\right)
\le\frac{M_\Phi^2}{m_1}
\le5M_\Phi^2n^{-1},
\end{aligned}
\]
where \(r_0\) is any record index in evaluation block \(1\), in channel \(g(i)\). The variance inequality follows because the cells are disjoint, each mark lies in \([0,1]\), and therefore the weighted single-record score has absolute value at most \(M_\Phi\). This calculation applies separately to the source and target channels.

Since \(\mathsf L_A-\overline{\mathsf L}_A=\sum_i\mathsf Z_i^{\mathrm{lin}}\), the finite-sum square inequality gives
\[
\begin{aligned}
\mathbb E_P[(\mathsf L_A-\overline{\mathsf L}_A)^2\mid\Gamma]
&\le7\sum_{i\in\mathcal I}
\mathbb E_P[(\mathsf Z_i^{\mathrm{lin}})^2\mid\Gamma]\\
&\le5\cdot7^2M_\Phi^2n^{-1}
\le10\cdot7^2M_\Phi^2n^{-1}.
\end{aligned}
\]
The same finite-sum argument proves square integrability of the centered linear correction.
% lean: CausalSmith.Stat.TransportCaceRoughnuisanceLength.normalized_flatBlockHeight_covariance_sum_le
% lean: CausalSmith.Stat.TransportCaceRoughnuisanceLength.condExp_sq_bounded_linear_cellAverage_le
% lean: CausalSmith.Stat.TransportCaceRoughnuisanceLength.linear_conditional_variance_with_memLp

\item \textit{Quadratic and cubic conditional second moments.}
For two coordinates using the same sample channel, direct expansion of the empirical bin sums gives
\[
\begin{aligned}
&\operatorname{Cov}_P\!\left(
H_{i,K}^{(b)}(x_{\ell+1,K}),
H_{j,K}^{(b)}(x_{r+1,K})
\right)\\
&\qquad=
\frac{K^2}{m_b}
\left\{
\mathbb E_P[
\mathsf B_{i,K,\ell,r_0}\mathsf B_{j,K,r,r_0}]
-\pi_{i,K,\ell}\pi_{j,K,r}
\right\}.
\end{aligned}
\]
For different channels this covariance is zero. On the same cell, bounded marks bound the product moment by \(C_f/K\); on distinct cells the product is zero. It follows that
\[
\left|
\operatorname{Cov}_P\!\left(
H_{i,K}^{(b)}(x_{\ell+1,K}),
H_{j,K}^{(b)}(x_{r+1,K})
\right)
\right|
\le
\begin{cases}
(C_fK+C_f^2)/m_b\le V_{\mathrm{bin}}K/m_b,
&\ell=r,\\
C_f^2/m_b\le V_{\mathrm{bin}}/m_b,
&\ell\ne r.
\end{cases}
\]

To apply these bounds, let \(s\in\{2,3\}\), let
\(\boldsymbol i=(i_1,\ldots,i_s)\in\mathcal I^s\), and let
\(\varnothing\ne\mathcal J\subseteq\{1,\ldots,s\}\). Define
\[
\begin{aligned}
\partial_{\boldsymbol i}\Phi_A
&:=\partial_{i_1\cdots i_s}\Phi_A,\\
\gamma^{(s)}_{\boldsymbol i,\mathcal J,\ell}
&:=
\frac1{s!}
\partial_{\boldsymbol i}\Phi_A(\widehat F(x_{\ell+1,K_s}))
\prod_{j\notin\mathcal J}
\delta^{\mathrm{av}}_{i_j,K_s,\ell},\\
\mathsf Z^{(s)}_{\boldsymbol i,\mathcal J}
&:=
\frac1{K_s}\sum_{\ell=0}^{K_s-1}
\gamma^{(s)}_{\boldsymbol i,\mathcal J,\ell}
\prod_{j\in\mathcal J}\eta^{(j)}_{i_j,K_s,\ell},\\
C_{\mathrm{coef},s}
&:=\frac{M_\Phi(1+U_{\mathrm{env}})^{s-1}}{s!}.
\end{aligned}
\]
Here \(K_s\) means \(K_2\) when \(s=2\) and \(K_3\) when \(s=3\), and an empty product equals \(1\). Since at most \(s-1\) shifts occur,
\[
|\gamma^{(s)}_{\boldsymbol i,\mathcal J,\ell}|
\le C_{\mathrm{coef},s}.
\]
The product expansion from the preceding step now reads
\[
\mathsf Q_A-\overline{\mathsf Q}_A
=\sum_{\boldsymbol i\in\mathcal I^2}
\sum_{\varnothing\ne\mathcal J\subseteq\{1,2\}}
\mathsf Z^{(2)}_{\boldsymbol i,\mathcal J},
\qquad
\mathsf C_A-\overline{\mathsf C}_A
=\sum_{\boldsymbol i\in\mathcal I^3}
\sum_{\varnothing\ne\mathcal J\subseteq\{1,2,3\}}
\mathsf Z^{(3)}_{\boldsymbol i,\mathcal J}.
\]

Put \(\nu:=|\mathcal J|\). Independence of the distinct evaluation blocks gives
\[
\mathbb E_P\!\left[
\prod_{j\in\mathcal J}
\eta^{(j)}_{i_j,K_s,\ell}\eta^{(j)}_{i_j,K_s,r}
\,\middle|\,\Gamma
\right]
=
\prod_{j\in\mathcal J}
\operatorname{Cov}_P\!\left(
H_{i_j,K_s}^{(j)}(x_{\ell+1,K_s}),
H_{i_j,K_s}^{(j)}(x_{r+1,K_s})
\right).
\]
There are \(K_s\) diagonal cell pairs and \(K_s(K_s-1)\) off-diagonal pairs. Expanding the square and applying the covariance bounds therefore yields
\[
\mathbb E_P\!\left[
(\mathsf Z^{(s)}_{\boldsymbol i,\mathcal J})^2
\,\middle|\,\Gamma
\right]
\le
C_{\mathrm{coef},s}^2V_{\mathrm{bin}}^\nu
\frac{5^\nu(K_s^{\nu-1}+1)}{n^\nu}.
\]
This also gives square integrability of every summand.

The constants satisfy \(U_{\mathrm{env}}>4\) and \(V_{\mathrm{bin}}>0\), so
\[
\begin{aligned}
\max\{10V_{\mathrm{bin}},25V_{\mathrm{bin}}^2\}
&\le4(1+U_{\mathrm{env}})^2(1+V_{\mathrm{bin}})^2,\\
\max\{10V_{\mathrm{bin}},25V_{\mathrm{bin}}^2,
125V_{\mathrm{bin}}^3\}
&\le36(1+U_{\mathrm{env}})^2(1+V_{\mathrm{bin}})^3.
\end{aligned}
\]
For \(\nu=1\), the preceding bound has factor \(10V_{\mathrm{bin}}/n\); for \(\nu=2\), it has factor \(25V_{\mathrm{bin}}^2(K_s+1)/n^2\); and for \(\nu=3\), it has factor \(125V_{\mathrm{bin}}^3(K_s^2+1)/n^3\).
Using \(n\ge1\), these inequalities give, for each summand,
\[
\begin{aligned}
\mathbb E_P[(\mathsf Z^{(2)}_{\boldsymbol i,\mathcal J})^2\mid\Gamma]
&\le
M_\Phi^2(1+U_{\mathrm{env}})^4(1+V_{\mathrm{bin}})^2
(n^{-1}+K_2n^{-2}),\\
\mathbb E_P[(\mathsf Z^{(3)}_{\boldsymbol i,\mathcal J})^2\mid\Gamma]
&\le
M_\Phi^2(1+U_{\mathrm{env}})^6(1+V_{\mathrm{bin}})^3
(n^{-1}+K_3n^{-2}+K_3^2n^{-3}).
\end{aligned}
\]
For example, \(K_s+1\le K_s+n\) absorbs the pair contribution, and \(K_3^2+1\le K_3^2+n^2\) absorbs the triple contribution.

There are exactly \(7^s(2^s-1)\) summands: \(7^s\) coordinate tuples, each with \(2^s-1\) nonempty subsets of evaluation positions. Define
\[
\mathcal V_s
:=
7^{2s}(2^s-1)^2M_\Phi^2
(1+U_{\mathrm{env}})^{2s}(1+V_{\mathrm{bin}})^s,
\qquad s=2,3.
\]
Applying the finite-sum square inequality to the displayed expansions gives
\[
\begin{aligned}
\mathbb E_P[(\mathsf Q_A-\overline{\mathsf Q}_A)^2\mid\Gamma]
&\le\mathcal V_2(n^{-1}+K_2n^{-2}),\\
\mathbb E_P[(\mathsf C_A-\overline{\mathsf C}_A)^2\mid\Gamma]
&\le\mathcal V_3(n^{-1}+K_3n^{-2}+K_3^2n^{-3}).
\end{aligned}
\]
% lean: CausalSmith.Stat.TransportCaceRoughnuisanceLength.conditional_variance_centeredSingle_cellAverage
% lean: CausalSmith.Stat.TransportCaceRoughnuisanceLength.conditional_variance_centeredPair_cellAverage
% lean: CausalSmith.Stat.TransportCaceRoughnuisanceLength.conditional_variance_centeredCubic_cellAverage
% lean: CausalSmith.Stat.TransportCaceRoughnuisanceLength.quadratic_conditional_variance_with_memLp
% lean: CausalSmith.Stat.TransportCaceRoughnuisanceLength.cubic_conditional_variance_with_memLp

\item \textit{Fluctuation bound at the common rate.}
The resolution bounds imply
\[
n^{-1}\le n^{-2/3},\qquad
n^{-1}+K_2n^{-2}\le2n^{-2/3},
\]
and
\[
K_3n^{-2}\le n^{-1},\qquad
K_3^2n^{-3}\le n^{-1},\qquad
n^{-1}+K_3n^{-2}+K_3^2n^{-3}\le3n^{-2/3}.
\]
Enlarging the nonnegative quadratic and cubic constants gives
\[
\begin{aligned}
\mathbb E_P[(\mathsf L_A-\overline{\mathsf L}_A)^2\mid\Gamma]
&\le10\cdot7^2M_\Phi^2n^{-2/3},\\
\mathbb E_P[(\mathsf Q_A-\overline{\mathsf Q}_A)^2\mid\Gamma]
&\le60\mathcal V_2n^{-2/3},\\
\mathbb E_P[(\mathsf C_A-\overline{\mathsf C}_A)^2\mid\Gamma]
&\le310\mathcal V_3n^{-2/3}.
\end{aligned}
\]
Define
\[
\mathcal V_{\mathrm{tot}}
:=3\{10\cdot7^2M_\Phi^2+60\mathcal V_2+310\mathcal V_3\}.
\]
Applying the three-term square inequality to the exact fluctuation identity, and then taking expectations, gives
\[
\mathbb E_P[(\widehat T_A-\overline T_A)^2\mid\Gamma]
\le\mathcal V_{\mathrm{tot}}n^{-2/3},
\qquad
\mathbb E_P[(\widehat T_A-\overline T_A)^2]
\le\mathcal V_{\mathrm{tot}}n^{-2/3}.
\]
The component square-integrability statements also give
\[
\widehat T_A-\overline T_A\in L^2(P).
\]
% lean: CausalSmith.Stat.TransportCaceRoughnuisanceLength.quadratic_conditional_variance_sample_rate
% lean: CausalSmith.Stat.TransportCaceRoughnuisanceLength.cubic_conditional_variance_sample_rate
% lean: CausalSmith.Stat.TransportCaceRoughnuisanceLength.linear_conditional_variance_sample_rate
% lean: CausalSmith.Stat.TransportCaceRoughnuisanceLength.centered_corrections_conditional_second_moment_sample_rate
% lean: CausalSmith.Stat.TransportCaceRoughnuisanceLength.centered_corrections_expected_conditional_second_moment_sample_rate
% lean: CausalSmith.Stat.TransportCaceRoughnuisanceLength.cubicEstimator_fluctuation_second_moment_sample_rate

\item \textit{Pilot moments.}
For a fixed coordinate and cell, define the centered training-record scores and their variance by
\[
\xi_{i,K,\ell,r}
:=\mathsf B_{i,K,\ell,r}-\pi_{i,K,\ell},
\qquad
\sigma_{i,K,\ell}^2
:=\mathbb E_P[\xi_{i,K,\ell,r}^2].
\]
Bounded bin scores give
\[
|\xi_{i,K,\ell,r}|\le1,\qquad
\mathbb E_P[\xi_{i,K,\ell,r}]=0,\qquad
\sigma_{i,K,\ell}^2
\le C_f/K+(C_f/K)^2
\le V_{\mathrm{bin}}/K.
\]
The centered training histogram is exactly
\[
H_{i,K}^{(0)}(x_{\ell+1,K})
-(\Pi_KF_i)(x_{\ell+1,K})
=
\frac K{m_0}\sum_{r\in\mathcal B_0^{g(i)}}
\xi_{i,K,\ell,r}.
\]
Independence therefore gives its second-moment bound
\[
\mathbb E_P\!\left[
\left\{H_{i,K}^{(0)}(x_{\ell+1,K})
-(\Pi_KF_i)(x_{\ell+1,K})\right\}^2
\right]
\le V_{\mathrm{bin}}K/m_0.
\]

For the eighth moment, expand the eighth power of the centered record sum and group terms by their equality pattern among the eight record positions. A pattern containing a singleton has expectation zero. Every surviving pattern has between one and four blocks. For a pattern with \(j\) blocks, boundedness of the centered scores bounds each block moment of order at least two by \(\sigma_{i,K,\ell}^2\), and there are at most \(m_0^j\) choices of its record indices. Its absolute contribution is consequently at most
\[
(m_0\sigma_{i,K,\ell}^2)^j
\le m_0\sigma_{i,K,\ell}^2
+(m_0\sigma_{i,K,\ell}^2)^4,
\qquad 1\le j\le4.
\]
There are at most \(8!\) equality patterns: encoding each block as a cycle in increasing order embeds the set of partitions into the permutations of eight positions. Thus
\[
\mathbb E_P\!\left[
\left(\sum_{r\in\mathcal B_0^{g(i)}}\xi_{i,K,\ell,r}\right)^8
\right]
\le8!\left\{
(m_0\sigma_{i,K,\ell}^2)^4
+m_0\sigma_{i,K,\ell}^2
\right\}.
\]
Scaling by \((K/m_0)^8\) yields
\[
\mathbb E_P\!\left[
\left\{H_{i,K}^{(0)}(x_{\ell+1,K})
-(\Pi_KF_i)(x_{\ell+1,K})\right\}^8
\right]
\le8!\left\{
V_{\mathrm{bin}}^4(K/m_0)^4
+V_{\mathrm{bin}}(K/m_0)^7
\right\}.
\]
% lean: CausalSmith.Stat.TransportCaceRoughnuisanceLength.iidBlockHeight_eighth_moment_le

Averaging the Hölder modulus over a cell gives
\[
|(\Pi_KF_i)(x)-F_i(x)|
\le H_{\mathrm{env}}K^{-1/8},
\qquad x\in E_{\ell,K}.
\]
Each clipping interval contains \(F_i(x)\), so clipping contracts the distance to that value. Splitting the histogram error at its cell mean, using
\((a+b)^2\le2(a^2+b^2)\) and
\((a+b)^8\le2^7(a^8+b^8)\) for nonnegative \(a,b\), gives
\[
\begin{aligned}
\mathbb E_P[\delta_i(x)^2]
&\le2V_{\mathrm{bin}}K_0/m_0
+2H_{\mathrm{env}}^2K_0^{-1/4},\\
\mathbb E_P[|\delta_i(x)|^8]
&\le2^7\,8!\left\{
V_{\mathrm{bin}}^4(K_0/m_0)^4
+V_{\mathrm{bin}}(K_0/m_0)^7
\right\}
+2^7H_{\mathrm{env}}^8K_0^{-1}.
\end{aligned}
\]
The dyadic bounds imply
\[
\begin{aligned}
K_0/m_0&\le5^{1/5}n^{-1/5},&
K_0^{-1/4}&\le2^{1/4}5^{1/5}n^{-1/5},\\
(K_0/m_0)^4&\le5^{4/5}n^{-4/5},&
(K_0/m_0)^7&\le5^{7/5}n^{-7/5}
\le5^{7/5}n^{-4/5},\\
K_0^{-1}&\le2\cdot5^{4/5}n^{-4/5}.
\end{aligned}
\]
Define
\[
\begin{aligned}
C_{\mathrm{pilot},2}
&:=2V_{\mathrm{bin}}5^{1/5}
+2^{5/4}H_{\mathrm{env}}^25^{1/5},\\
C_{\mathrm{pilot},8}
&:=2^7\,8!\left\{
V_{\mathrm{bin}}^45^{4/5}
+V_{\mathrm{bin}}5^{7/5}\right\}
+2^8H_{\mathrm{env}}^85^{4/5}.
\end{aligned}
\]
We have obtained, for every \(x\in[0,1]\) and \(i\in\mathcal I\),
\[
\mathbb E_P[\delta_i(x)^2]
\le C_{\mathrm{pilot},2}n^{-1/5},
\qquad
\mathbb E_P[|\delta_i(x)|^8]
\le C_{\mathrm{pilot},8}n^{-4/5}.
\]
Finally, define the spatial pilot errors by
\[
\mathsf E_{i,\mathrm{pilot}}:=\int_0^1|\delta_i(x)|\,dx.
\]
Cauchy--Schwarz on the unit interval and integration of the pointwise bounds give
\[
\begin{aligned}
\mathbb E_P[\mathsf E_{i,\mathrm{pilot}}^2]
&\le\mathbb E_P\int_0^1\delta_i(x)^2\,dx
\le C_{\mathrm{pilot},2}n^{-1/5},\\
\mathbb E_P\int_0^1|\delta_i(x)|^8\,dx
&\le C_{\mathrm{pilot},8}n^{-4/5}.
\end{aligned}
\]
Joint measurability and the bound \(|\delta_i|\le2C_f\) justify these integrations.
% lean: CausalSmith.Stat.TransportCaceRoughnuisanceLength.pilot_second_moment
% lean: CausalSmith.Stat.TransportCaceRoughnuisanceLength.pilot_eighth_moment
% lean: CausalSmith.Stat.TransportCaceRoughnuisanceLength.pilot_error_L1_second_moment

\item \textit{Spatial Taylor remainder.}
For \(s=1,2,3\), define the spatial Taylor terms by
\[
\mathsf T^{\mathrm{Taylor}}_{A,s}
:=
\frac1{s!}
\sum_{(i_1,\ldots,i_s)\in\mathcal I^s}
\int_0^1
\partial_{i_1\cdots i_s}\Phi_A(\widehat F(x))
\prod_{j=1}^s\delta_{i_j}(x)\,dx.
\]
In particular,
\(\mathsf T^{\mathrm{Taylor}}_{A,1}
=\overline{\mathsf L}_A\).

The remainder estimate follows from an exact ratio identity. For \(z=0,1\), define
\[
\Psi_{Az}(\mathbf w)
:=\frac{w_t w_{r_{Az}}}{w_{q_z}},
\qquad
\mathbf w\in\mathbb R^7,\quad w_{q_z}>0,
\]
using the named-coordinate convention of \cref{obj:synth_14}. For each \(x\in[0,1]\), direct expansion and reduction to a common denominator give
\[
\begin{aligned}
&\Psi_{Az}(F(x))-\Psi_{Az}(\widehat F(x))
-\sum_{s=1}^3\frac1{s!}
D^s\Psi_{Az}(\widehat F(x))
[\delta(x),\ldots,\delta(x)]\\
&\quad=
\frac{
\delta_{q_z}(x)^2
\{\widehat F_{q_z}(x)\delta_t(x)
-\widehat F_t(x)\delta_{q_z}(x)\}
\{\widehat F_{q_z}(x)\delta_{r_{Az}}(x)
-\widehat F_{r_{Az}}(x)\delta_{q_z}(x)\}
}{
\widehat F_{q_z}(x)^4 F_{q_z}(x)
}.
\end{aligned}
\]
Here
\[
\delta(x):=(\delta_0(x),\ldots,\delta_6(x)).
\]
Both denominator coordinates are at least \(c_f/4\). Each numerator factor in braces is bounded by \(C_f\) times the sum of the absolute increments in its two coordinates. Bounding these sums and \(|\delta_{q_z}|\) by the full coordinate sum therefore gives
\[
\left|
\Psi_{Az}(F)-\Psi_{Az}(\widehat F)
-\sum_{s=1}^3\frac1{s!}
D^s\Psi_{Az}(\widehat F)[\delta,\ldots,\delta]
\right|
\le C_f^2(4/c_f)^5
\left(\sum_{i\in\mathcal I}|\delta_i|\right)^4.
\]
Since \(\Phi_A=\Psi_{A1}-\Psi_{A0}\), define its pointwise remainder by
\[
\begin{aligned}
\mathcal R_{A,4}(x)
:=\Phi_A(F(x))-\Phi_A(\widehat F(x))
-\sum_{s=1}^3\frac1{s!}
\sum_{(i_1,\ldots,i_s)\in\mathcal I^s}
\partial_{i_1\cdots i_s}\Phi_A(\widehat F(x))
\prod_{j=1}^s\delta_{i_j}(x).
\end{aligned}
\]
The two ratio bounds and \(C_f^2\le(1+C_f)^2\) imply
\[
|\mathcal R_{A,4}(x)|
\le2(1+C_f)^2(4/c_f)^5
\left(\sum_i|\delta_i(x)|\right)^4
=
\frac{M_\Phi}{24}
\left(\sum_i|\delta_i(x)|\right)^4.
\]
% lean: CausalSmith.Stat.TransportCaceRoughnuisanceLength.cubicRatio_cubic_taylor_remainder_eq
% lean: CausalSmith.Stat.TransportCaceRoughnuisanceLength.Phi_cubic_taylor_remainder_abs_le

Define the integrated remainder and its majorant by
\[
\mathsf R_{\mathrm{rem},A}
:=\int_0^1\mathcal R_{A,4}(x)\,dx,
\qquad
\mathsf G_{\mathrm{rem}}
:=\frac{M_\Phi}{24}
\int_0^1\left(\sum_i|\delta_i(x)|\right)^4dx.
\]
Then
\[
|\mathsf R_{\mathrm{rem},A}|\le\mathsf G_{\mathrm{rem}},
\qquad
0\le\mathsf G_{\mathrm{rem}}
\le\frac{M_\Phi}{24}(14C_f)^4.
\]
Cauchy--Schwarz on \([0,1]\), followed by the finite-sum eighth-power inequality, gives
\[
\begin{aligned}
\mathbb E_P[\mathsf G_{\mathrm{rem}}^2]
&\le
(M_\Phi/24)^2
\mathbb E_P\int_0^1
\left(\sum_i|\delta_i(x)|\right)^8dx\\
&\le
(M_\Phi/24)^2\,7^7
\sum_{i\in\mathcal I}\mathbb E_P\int_0^1|\delta_i(x)|^8dx\\
&\le
(M_\Phi/24)^2\,7^8C_{\mathrm{pilot},8}n^{-4/5}.
\end{aligned}
\]
Define
\[
\mathcal B_{\mathrm{rem}}
:=(M_\Phi/24)^2\,7^8C_{\mathrm{pilot},8}.
\]
Since \(n^{-4/5}\le n^{-2/3}\),
\[
\mathbb E_P[\mathsf G_{\mathrm{rem}}^2]
\le\mathcal B_{\mathrm{rem}}n^{-2/3}.
\]
Integrating the defining Taylor identity also yields
\[
T_A=
\int_0^1\Phi_A(\widehat F(x))\,dx
+\sum_{s=1}^3\mathsf T^{\mathrm{Taylor}}_{A,s}
+\mathsf R_{\mathrm{rem},A}.
\]
% lean: CausalSmith.Stat.TransportCaceRoughnuisanceLength.pilot_remainder_integral_abs_le_majorant
% lean: CausalSmith.Stat.TransportCaceRoughnuisanceLength.pilotRemainderMajorant_integrable_sq_and_moment
% lean: CausalSmith.Stat.TransportCaceRoughnuisanceLength.pilotRemainderMajorant_sample_rate
% lean: CausalSmith.Stat.TransportCaceRoughnuisanceLength.spatialTaylorRemainder_eq_transport_sub_coordinate_terms

\item \textit{Exact spatial projection errors.}
For \(K_0\mid K\), define
\[
\varepsilon_{i,K}(x):=F_i(x)-(\Pi_KF_i)(x),
\qquad x\in[0,1].
\]
On each cell,
\[
\delta_i=\delta^{\mathrm{av}}_{i,K}+\varepsilon_{i,K},
\qquad
\int_{E_{\ell,K}}\varepsilon_{i,K}(x)\,dx=0,
\qquad
|\varepsilon_{i,K}(x)|\le H_{\mathrm{env}}K^{-1/8}.
\]
The pilot and all averaged shifts are constant on that cell. Hence any cell product containing exactly one averaging residual has integral zero. More explicitly, for any function \(\vartheta:\mathbb R^7\to\mathbb R\) and any finite list of coordinates \(i_1,\ldots,i_s\),
\[
\int_{E_{\ell,K}}
\vartheta(\widehat F(x))
\left\{\prod_{j=1}^s\delta^{\mathrm{av}}_{i_j,K}(x)\right\}
\varepsilon_{i,K}(x)\,dx=0.
\]
This includes \(s=0\), with the empty product equal to \(1\).
% lean: CausalSmith.Stat.TransportCaceRoughnuisanceLength.spatial_pilot_single_residual_cancellation

Define the two projection errors by
\[
\mathsf R_{\mathrm{quad},A}
:=\mathsf T^{\mathrm{Taylor}}_{A,2}
-\overline{\mathsf Q}_A,
\qquad
\mathsf R_{\mathrm{cub},A}
:=\mathsf T^{\mathrm{Taylor}}_{A,3}
-\overline{\mathsf C}_A.
\]
The nested-cell identities identify these as differences between the true spatial Taylor terms and the exact midpoint projection means. For the quadratic error, expanding the two factors and cancelling the single-residual terms gives
\[
\mathsf R_{\mathrm{quad},A}
=
\frac12\sum_{i,j\in\mathcal I}\sum_{\ell=0}^{K_2-1}
\partial_{ij}\Phi_A(\widehat F(x_{\ell+1,K_2}))
\int_{E_{\ell,K_2}}
\varepsilon_{i,K_2}(x)\varepsilon_{j,K_2}(x)\,dx.
\]
Thus
\[
|\mathsf R_{\mathrm{quad},A}|
\le\frac{7^2}{2}M_\Phi H_{\mathrm{env}}^2K_2^{-1/4}.
\]
Using the lower dyadic bound,
\[
K_2^{-1/2}
\le2^{1/2}5^{2/3}n^{-2/3}.
\]
Define
\[
\mathcal B_{\mathrm{quad}}
:=7^4M_\Phi^2H_{\mathrm{env}}^4
2^{1/2}5^{2/3}.
\]
Squaring the absolute bound and enlarging its factor \(1/4\) gives the pointwise inequality
\[
\mathsf R_{\mathrm{quad},A}^2
\le\frac{7^4M_\Phi^2H_{\mathrm{env}}^4}{4}K_2^{-1/2}
\le\mathcal B_{\mathrm{quad}}n^{-2/3}.
\]
% lean: CausalSmith.Stat.TransportCaceRoughnuisanceLength.quadraticCellProjectionBias_eq_spatial_sub_projectionMean
% lean: CausalSmith.Stat.TransportCaceRoughnuisanceLength.quadraticCellProjectionBias_abs_le
% lean: CausalSmith.Stat.TransportCaceRoughnuisanceLength.quadraticCellProjectionBias_sq_sample_rate

For the cubic error, work on a \(K_3\)-cell. Expanding the three factors leaves, after single-residual cancellation,
\[
\begin{aligned}
&\int_{E_{\ell,K_3}}
\left\{
\delta_i\delta_j\delta_k
-\delta^{\mathrm{av}}_{i,K_3,\ell}
 \delta^{\mathrm{av}}_{j,K_3,\ell}
 \delta^{\mathrm{av}}_{k,K_3,\ell}
\right\}dx\\
&\quad=
\delta^{\mathrm{av}}_{i,K_3,\ell}
\int_{E_{\ell,K_3}}\varepsilon_{j,K_3}\varepsilon_{k,K_3}\,dx
+\delta^{\mathrm{av}}_{j,K_3,\ell}
\int_{E_{\ell,K_3}}\varepsilon_{i,K_3}\varepsilon_{k,K_3}\,dx\\
&\qquad+
\delta^{\mathrm{av}}_{k,K_3,\ell}
\int_{E_{\ell,K_3}}\varepsilon_{i,K_3}\varepsilon_{j,K_3}\,dx
+\int_{E_{\ell,K_3}}
\varepsilon_{i,K_3}\varepsilon_{j,K_3}\varepsilon_{k,K_3}\,dx.
\end{aligned}
\]
To retain the pilot error, use the exact cell-average identity
\[
\frac{|\delta^{\mathrm{av}}_{i,K_3,\ell}|}{K_3}
=
\left|\int_{E_{\ell,K_3}}\delta_i(x)\,dx\right|
\le\int_{E_{\ell,K_3}}|\delta_i(x)|\,dx.
\]
Consequently, the absolute value of the preceding cell integral is at most
\[
H_{\mathrm{env}}^2K_3^{-1/4}
\int_{E_{\ell,K_3}}(|\delta_i|+|\delta_j|+|\delta_k|)\,dx
+\frac{H_{\mathrm{env}}^3K_3^{-3/8}}{K_3}.
\]
Multiplying by the cubic coefficient, bounded by \(M_\Phi/6\), and summing over cells and coordinate triples gives
\[
|\mathsf R_{\mathrm{cub},A}|
\le
C_{\mathrm{bias},2}K_3^{-1/4}
\sum_{i\in\mathcal I}\mathsf E_{i,\mathrm{pilot}}
+C_{\mathrm{bias},3}K_3^{-3/8},
\]
where
\[
C_{\mathrm{bias},2}
:=\frac{M_\Phi}{2}7^2H_{\mathrm{env}}^2,
\qquad
C_{\mathrm{bias},3}
:=\frac{M_\Phi}{6}7^3H_{\mathrm{env}}^3.
\]
Define the square majorant by
\[
\mathsf S_{\mathrm{cub}}
:=
2C_{\mathrm{bias},2}^2K_3^{-1/2}
\,7\sum_{i\in\mathcal I}\mathsf E_{i,\mathrm{pilot}}^2
+2C_{\mathrm{bias},3}^2K_3^{-3/4}.
\]
The two-term square inequality and the coordinate-sum square inequality imply
\[
\mathsf R_{\mathrm{cub},A}^2\le\mathsf S_{\mathrm{cub}}.
\]
The pilot \(L^1\) moment bound therefore gives
\[
\mathbb E_P[\mathsf S_{\mathrm{cub}}]
\le
2C_{\mathrm{bias},2}^2\,7^2C_{\mathrm{pilot},2}
K_3^{-1/2}n^{-1/5}
+2C_{\mathrm{bias},3}^2K_3^{-3/4}.
\]
The lower bound \(K_3\ge n/10\) yields
\[
\begin{aligned}
K_3^{-1/2}n^{-1/5}
&\le2^{1/2}5^{1/2}n^{-7/10}
\le2^{1/2}5^{1/2}n^{-2/3},\\
K_3^{-3/4}
&\le2^{3/4}5^{3/4}n^{-3/4}
\le2^{3/4}5^{3/4}n^{-2/3}.
\end{aligned}
\]
Define
\[
\mathcal B_{\mathrm{cub}}
:=
2\left\{
C_{\mathrm{bias},2}^2\,7^2C_{\mathrm{pilot},2}
2^{1/2}5^{1/2}
+C_{\mathrm{bias},3}^2\,2^{3/4}5^{3/4}
\right\}.
\]
We conclude that
\[
\mathbb E_P[\mathsf S_{\mathrm{cub}}]
\le\mathcal B_{\mathrm{cub}}n^{-2/3}.
\]
% lean: CausalSmith.Stat.TransportCaceRoughnuisanceLength.cubic_midpoint_cell_projection_L1_bound
% lean: CausalSmith.Stat.TransportCaceRoughnuisanceLength.cubic_spatial_projection_bias_abs_le_L1
% lean: CausalSmith.Stat.TransportCaceRoughnuisanceLength.cubic_bias_resolution_terms_sample_rate

\item \textit{Squared bias of the projected mean.}
The integrated Taylor identity and the definitions of the projection errors give the exact signed identity
\[
\overline T_A-T_A
=
-\left(
\mathsf R_{\mathrm{rem},A}
+\mathsf R_{\mathrm{quad},A}
+\mathsf R_{\mathrm{cub},A}
\right).
\]
Using the three-term square inequality and the preceding majorants gives
\[
(\overline T_A-T_A)^2
\le
3\left\{
\mathsf G_{\mathrm{rem}}^2
+\mathcal B_{\mathrm{quad}}n^{-2/3}
+\mathsf S_{\mathrm{cub}}
\right\}.
\]
The right-hand side is integrable, by the remainder bound and the spatial pilot moment bounds. Define
\[
\mathcal B_{\mathrm{tot}}
:=3\left(
\mathcal B_{\mathrm{rem}}
+\mathcal B_{\mathrm{quad}}
+\mathcal B_{\mathrm{cub}}
\right).
\]
Integration yields
\[
\mathbb E_P[(\overline T_A-T_A)^2]
\le\mathcal B_{\mathrm{tot}}n^{-2/3},
\qquad
\overline T_A-T_A\in L^2(P).
\]
% lean: CausalSmith.Stat.TransportCaceRoughnuisanceLength.combined_spatial_projection_bias_second_moment
% lean: CausalSmith.Stat.TransportCaceRoughnuisanceLength.projectedEstimatorMean_sub_transport_eq_spatial_bias
% lean: CausalSmith.Stat.TransportCaceRoughnuisanceLength.projectedEstimatorMean_squared_bias_sample_rate

\item \textit{Orthogonality and final assembly.}
Define the fluctuation and projected bias by
\[
\mathsf R_{\mathrm{fluc},A}:=\widehat T_A-\overline T_A,
\qquad
\mathsf R_{\mathrm{bias},A}:=\overline T_A-T_A.
\]
The exact conditional means imply
\[
\mathbb E_P[\mathsf R_{\mathrm{fluc},A}\mid\Gamma]=0.
\]
The projected bias is training-measurable, and both quantities are square-integrable. Their product is therefore integrable, and pulling out the training-measurable factor gives
\[
\begin{aligned}
\mathbb E_P[
\mathsf R_{\mathrm{bias},A}\mathsf R_{\mathrm{fluc},A}]
&=
\mathbb E_P\!\left[
\mathsf R_{\mathrm{bias},A}
\mathbb E_P[\mathsf R_{\mathrm{fluc},A}\mid\Gamma]
\right]
=0.
\end{aligned}
\]
Expanding the square consequently gives the exact risk decomposition
\[
\begin{aligned}
\mathbb E_P[(\widehat T_A-T_A)^2]
&=
\mathbb E_P[\mathsf R_{\mathrm{fluc},A}^2]
+\mathbb E_P[\mathsf R_{\mathrm{bias},A}^2]\\
&\le
(\mathcal V_{\mathrm{tot}}+\mathcal B_{\mathrm{tot}})n^{-2/3}.
\end{aligned}
\]
Finally, substitution of the defined constants gives
\[
\begin{aligned}
\mathcal V_{\mathrm{tot}}+\mathcal B_{\mathrm{tot}}
&=
3\left(
\mathcal B_{\mathrm{rem}}
+\mathcal B_{\mathrm{quad}}
+\mathcal B_{\mathrm{cub}}
\right)
+3\left\{
10\cdot7^2M_\Phi^2
+60\mathcal V_2+310\mathcal V_3
\right\}\\
&=C_{\mathrm{mse}},
\end{aligned}
\]
with exactly the expanded expression in the statement. Each term is finite and explicitly computable from \((c_f,C_f,L)\), since \(c_f>0\). All bounds apply to both response choices and every \(P\in\mathcal M_n\), proving the asserted uniform inequality.
% lean: CausalSmith.Stat.TransportCaceRoughnuisanceLength.condExp_cubicEstimator_sub_projectedEstimatorMean_eq_zero
% lean: CausalSmith.Stat.TransportCaceRoughnuisanceLength.integral_projected_bias_mul_fluctuation_eq_zero
% lean: CausalSmith.Stat.TransportCaceRoughnuisanceLength.cubicEstimator_mse_eq_fluctuation_add_projected_bias
% lean: CausalSmith.Stat.TransportCaceRoughnuisanceLength.cubicEstimator_mse_sample_rate
% lean: CausalSmith.Stat.TransportCaceRoughnuisanceLength.marked_cubic_mse
\end{enumerate}
\end{proof}

The eighth-moment pilot bound controls the squared fourth-order remainder. The quadratic and cubic resolutions control projection errors, and conditioning on the training block preserves within-block coordinate dependence while retaining independence across the evaluation blocks. The displayed constant includes all approximation and fluctuation terms.

\subsection{Supporting length bounds}

The mean-square result calibrates both transported contrasts. Identification supplies the population affine-score identity at the target effect, and the bounded effect domain supplies the constant branch.

\begin{theoremv}{thm:honest-upper-full-elbow}[Honest upper length envelope]
Fix the noncoverage probability \(\alpha\), the density envelopes \(c_f,C_f\), and the Hölder radius \(L\), satisfying:
\begin{itemize}
\item \textbf{(Noncoverage.)} \(0<\alpha<1\).
\item \textbf{(Lower density envelope.)} \(0<c_f<1\).
\item \textbf{(Upper density envelope.)} \(1<C_f\).
\item \textbf{(Hölder radius.)} \(1<L\).
\end{itemize}
The score interval \(I_n\) of \cref{obj:def:score-interval} belongs to the class \(\mathfrak H_n\) of honest intervals in \cref{obj:def:honest-intervals} for every positive integer \(n\).

Set the sample-size threshold to
\[
n_0=256.
\]
Write \(\mu(P)\) for the transported first stage, namely the target-averaged source receipt contrast:
\[
\mu(P)=T_D(P).
\]
There exists a finite constant \(C_0=C_0(\alpha,c_f,C_f,L)>0\) such that, for every integer \(n\ge n_0\) and every \(P\in\mathcal M_n\), the model class of \cref{obj:def:model-class},
\[
\mathbb E_P\!\left[\lambda_{\Theta}(I_n)\right]
\le C_0\min\left\{1,\frac{n^{-1/3}}{\mu(P)}\right\},
\]
where \(\lambda_{\Theta}\) denotes Lebesgue length restricted to the effect domain \(\Theta\), and the expectation is under the joint law of the source and target samples, each of size \(n\).

With the same constant \(C_0\), for every integer \(n\ge n_0\) and every real \(a\) satisfying \(0<a\le 1/4\), the minimax expected-length frontier \(L_n(a)\) of \cref{obj:def:length-frontier} satisfies
\[
L_n(a)\le C_0\min\left\{1,\frac{n^{-1/3}}{a}\right\}.
\]
\end{theoremv}

The upper result gives all-positive-sample-size honesty and, from \(n_0\) onward, both a law-specific length bound and a uniform bound on each strength slice. The converse is stated next; its causally legal experiment and testing argument are developed in \cref{sec:lower-appendix}.

\begin{theoremv}{thm:honest-lower-full-elbow}[Honest length lower bound]
Fix constants satisfying:
\begin{itemize}
\item \textbf{(Coverage.)} The noncoverage probability satisfies \(0<\alpha<1\).
\item \textbf{(Density envelopes.)} The lower and upper density envelopes satisfy \(0<c_f<1<C_f\).
\item \textbf{(Smoothness.)} The common Hölder radius satisfies \(L>1\).
\end{itemize}
Let \(\mathcal M_n\) be the model class in \cref{obj:def:model-class}, and let \(\mathcal M_n(a)\) be its strength slice in \cref{obj:def:strength-slice}. Write \(\theta(P)\) for the target-population average treatment effect among compliers under \(P\), and let the effect domain and sample-size threshold be
\[
\Theta=[-1,1],
\qquad n_0=256.
\]
There exists a constant \(c_0=c_0(\alpha,c_f,C_f,L)>0\) such that, for every integer \(n\ge n_0\) and every \(0<a\le 1/4\), the following holds. Every measurable confidence-set procedure \(C_n\), based on \(n\) source records and \(n\) target covariate records, with \(C_n(\omega)\subseteq\Theta\) for every sample realization \(\omega\), and satisfying
\[
P\bigl(\theta(P)\in C_n\bigr)\ge 1-\alpha
\qquad\text{for every }P\in\mathcal M_n,
\]
obeys
\[
\sup_{P\in\mathcal M_n(a)}
\mathbb E_P\!\left[\lambda_{\Theta}(C_n)\right]
\ge c_0\min\left\{1,\frac{n^{-1/3}}{a}\right\},
\qquad
\lambda_{\Theta}(C_n)
=\operatorname{Leb}(C_n\cap\Theta).
\]
Here probabilities and expectations are taken under the observed two-sample law. Moreover, the minimax expected-length frontier \(L_n(a)\) of \cref{obj:def:length-frontier} satisfies
\[
L_n(a)\ge c_0\min\left\{1,\frac{n^{-1/3}}{a}\right\}.
\]
\end{theoremv}

Together these bounds yield the exact frontier stated in \cref{obj:thm:sharp-length-frontier-full-elbow}. Full proof scripts for the anchored statements are collected in \cref{sec:deferred-proofs}.



\section{Lower-bound construction and proof}
\label{sec:lower-appendix}

The converse in \cref{obj:thm:honest-lower-full-elbow} rests on a \leanref{S-6}{lower experiment}, specified in \cref{obj:def:legal-iv-mixture}, whose components satisfy the causal and transport restrictions while spanning a continuum of complier effects. The construction couples perturbations of the target covariate distribution, encouragement probabilities, and receipt--outcome cells. Throughout, the observed experiment consists of independent source and target samples, each of size \(n\).

We first specify the interval cells on which the perturbations are placed.

\begin{definitionv}{synth_18}[Resolution-\(K\) interval cells]
For a positive integer \(K\) and an index \(\ell\in\{0,\ldots,K-1\}\), define
\[
\operatorname{cell}(K,\ell)=
\begin{cases}
[\ell/K,(\ell+1)/K),&0\le\ell<K-1,\\
[(K-1)/K,1],&\ell=K-1.
\end{cases}
\]
Thus each cell includes its left endpoint, and the final cell also includes the endpoint \(1\).
\end{definitionv}

The endpoint convention assigns every covariate value in the unit interval to one cell. A centered sinusoidal bump then supplies the within-cell perturbation.

\begin{definitionv}{synth_12}[Tiled and coupled perturbations]
The tiled perturbation \(u_{\lambda}\), for \(c_{\star}\in\mathbb R\), \(n\in\mathbb N\), and a cellwise sign vector \(\lambda\in\{-1,1\}^{K_{\mathrm L}}\), is defined for every \(x\in\mathbb R\) by
\[
K_{\mathrm L}=\left\lceil n^{4/3}\right\rceil,
\qquad
h_n=c_{\star}K_{\mathrm L}^{-1/8},
\qquad
u_{\lambda}(x)=
\sum_{\ell=0}^{K_{\mathrm L}-1}
\begin{cases}
h_n\lambda_\ell\sin\!\bigl(2\pi(K_{\mathrm L}x-\ell)\bigr),
& x\in\operatorname{cell}(K_{\mathrm L},\ell),\\
0,&\text{otherwise}.
\end{cases}
\]
Each sign \(\lambda_\ell\) equals \(1\) when the supplied Boolean entry for cell \(\ell\) is true and \(-1\) when it is false. For \(0\leq\ell<K_{\mathrm L}\), the cells are
\[
\operatorname{cell}(K_{\mathrm L},\ell)=
\begin{cases}
[\ell/K_{\mathrm L},1],
& \ell+1=K_{\mathrm L},\\
[\ell/K_{\mathrm L},(\ell+1)/K_{\mathrm L}),
& \ell+1\ne K_{\mathrm L}.
\end{cases}
\]
At \(n=0\), the sum is empty and equals zero, with the convention \(0^{-1/8}=0\).
The coupled perturbation \(v_{\lambda,\tau}\), for any \(\tau\in\mathbb R\), is defined by
\[
v_{\lambda,\tau}(x)=\tau u_{\lambda}(x),
\qquad x\in\mathbb R.
\]
\end{definitionv}

The bump \leanref{sym:B}{\ensuremath{B}} in \cref{obj:synth_16} has zero integral and positive squared integral. Its endpoint values allow adjacent rescaled copies to meet continuously, while its antisymmetry will make the effect separation common across cellwise sign choices. The following definition specifies the resolution and height of the tiled perturbation.

\begin{definitionv}{synth_15}[Baseline receipt–outcome cell laws]
For a baseline compliance share \(b\in(0,1/4]\) and \(z\in\{0,1\}\), set
\[
r_z^{(b)}=\frac{1+(2z-1)b}{2}.
\]
Define the baseline conditional receipt–outcome probability mass function on \(\{0,1\}^2\) by
\[
R_z^{(b)}(d,y)=
\begin{cases}
r_z^{(b)}/2,&d=1,\\
(1-r_z^{(b)})/2,&d=0,
\end{cases}
\qquad z,d,y\in\{0,1\}.
\]
Equivalently, for each \(y\in\{0,1\}\),
\[
R_0^{(b)}(0,y)=R_1^{(b)}(1,y)=\frac{1+b}{4},
\qquad
R_0^{(b)}(1,y)=R_1^{(b)}(0,y)=\frac{1-b}{4}.
\]
Under this baseline law, receipt has probability \(r_z^{(b)}\) conditional on instrument value \(z\), and the binary outcome is independent of receipt with probability \(1/2\) for each outcome value. In the lower-mixture construction, the baseline share is \(b=\max\{a,n^{-1/3}\}\) for \(n\ge n_0=256\) and \(0<a\le1/4\).
\end{definitionv}

In \cref{obj:synth_12}, \leanref{sym:K_{\mathrm L}}{\ensuremath{K_{\mathrm L}}} is the number of perturbation cells, \leanref{sym:h_n}{\ensuremath{h_n}} is the perturbation height, and \leanref{sym:\lambda}{\ensuremath{\lambda}} is the cellwise sign vector. The amplitude \leanref{sym:c_{\star}}{\ensuremath{c_{\star}}} controls the height relative to the resolution. The tiled perturbation \leanref{sym:u_{\lambda}}{\ensuremath{u_{\lambda}}} is used in the probability-law construction of \cref{obj:def:legal-iv-mixture}.

Before forming those laws, we specify the baseline receipt--outcome probabilities. They retain a positive compliance share while assigning equal probability to the two outcome values within each receipt cell.

\begin{algorithmv}{def:legal-iv-mixture}[Legal IV mixture construction]
The legal IV mixture is the triple \((Q_{\lambda,\tau},M_{\tau},P_{\star})\), consisting of a component law, its uniform sign-mixture experiment, and a center law, defined for the following parameters:
\begin{itemize}
\item \textbf{(Sample size.)} \(n\in\mathbb N\) satisfies
\[
n\ge n_0=256,
\qquad
K_{\mathrm L}=\left\lceil n^{4/3}\right\rceil.
\]
Here \(K_{\mathrm L}\) is the number of perturbation cells.
\item \textbf{(Baseline strength.)} The strength parameter \(a\) satisfies \(0<a\le1/4\), and the baseline compliance share is
\[
b=\bar a=\max\{a,n^{-1/3}\}.
\]
\item \textbf{(Amplitude and admissibility.)} The amplitude \(c_{\star}\) satisfies \(c_{\star}<1\). With perturbation height
\[
h=h_n=c_{\star}K_{\mathrm L}^{-1/8},
\]
require \(n>0\) and \(\operatorname{Adm}_{n,a}(c_{\star})\), where
\[
\operatorname{Adm}_{n,a}(c_{\star})
\Longleftrightarrow
\left(
0<c_{\star}
\ \text{and}\
h\le\frac{3}{16}
\ \text{and}\
h^2\le b\left(\frac14-h^2\right)
\right).
\]
\item \textbf{(Sign and continuum indices.)} The cellwise sign vector and continuum index satisfy
\[
\lambda\in\{-1,1\}^{K_{\mathrm L}},
\qquad
\tau\in[-1,1].
\]
For the tiled covariate perturbation \(u_{\lambda}\) indexed by these signs, set
\[
v_{\lambda,\tau}=\tau u_{\lambda}.
\]
\end{itemize}

Define the source and target covariate densities and encouragement probabilities by
\[
f_{\mathrm S}=1,
\qquad
f_{\mathrm T}=1+u_{\lambda},
\qquad
e_1^{(\lambda)}=\frac12+u_{\lambda},
\qquad
e_0^{(\lambda)}=\frac12-u_{\lambda}.
\]
For \(z,d,y\in\{0,1\}\), the source joint cells conditional on covariates are
\[
Q_{\lambda,\tau}(Z=z,D=d,Y=y\mid X=x)
=
e_z^{(\lambda)}(x)R_z^{(b)}(d,y)
+\frac{v_{\lambda,\tau}(x)(2y-1)}4,
\]
where \(R_z^{(b)}(d,y)\) is the baseline receipt–outcome cell probability at compliance share \(b\). The cells conditional on encouragement and covariates are
\[
Q_{\lambda,\tau}(D=d,Y=y\mid Z=z,X=x)
=
\frac{
e_z^{(\lambda)}(x)R_z^{(b)}(d,y)
+v_{\lambda,\tau}(x)(2y-1)/4
}{
e_z^{(\lambda)}(x)
}.
\]

Let \(C\) be the complier indicator, \(D(z)\) potential receipt under encouragement \(z\), and \(Y(d)\) the potential outcome under receipt state \(d\). Define the complier outcome margins by
\[
M_{1y}(x)
=
P(C=1,Y(1)=y\mid X=x)
=
\frac b2
-
\frac{(2y-1)u_{\lambda}(x)v_{\lambda,\tau}(x)}
{2\{1/4-u_{\lambda}(x)^2\}},
\]
\[
M_{0y}(x)
=
P(C=1,Y(0)=y\mid X=x)
=
\frac b2
+
\frac{(2y-1)u_{\lambda}(x)v_{\lambda,\tau}(x)}
{2\{1/4-u_{\lambda}(x)^2\}},
\qquad y\in\{0,1\}.
\]
Assign always-takers and never-takers the conditional masses
\[
P(D(0)=D(1)=1,Y(1)=y\mid X=x)
=
Q_{\lambda,\tau}(D=1,Y=y\mid Z=0,X=x),
\]
\[
P(D(0)=D(1)=0,Y(0)=y\mid X=x)
=
Q_{\lambda,\tau}(D=0,Y=y\mid Z=1,X=x).
\]
For compliers, set \(D(0)=0\), \(D(1)=1\), and
\[
P(Y(d)=y\mid X=x,C=1)
=
\frac{M_{dy}(x)}b,
\qquad d,y\in\{0,1\},
\]
coupling \(Y(0)\) and \(Y(1)\) independently under these binary marginals. Set
\[
Y(0)=0\quad\text{for always-takers},
\qquad
Y(1)=0\quad\text{for never-takers}.
\]
Draw encouragement according to
\[
Z\ \perp\ (D(0),D(1),Y(0),Y(1))\mid X,
\qquad
P(Z=1\mid X=x)=e_1^{(\lambda)}(x),
\]
and observe
\[
D=D(Z),
\qquad
Y=Y(D).
\]
Use the same conditional primitive full-data law given \(X\) in source and target. Represent the resulting combined law \(Q_{\lambda,\tau}\) by
\[
Q_{\lambda,\tau}(S=\mathrm{source})
=
Q_{\lambda,\tau}(S=\mathrm{target})
=
\frac12.
\]
This allocation is a convention for representing the constructed witness; admissibility specifies the parameter domain of the construction.

Define the center law by
\[
P_{\star}
=
\left.Q_{\lambda,\tau}\right|_{u_{\lambda}=v_{\lambda,\tau}=0},
\]
retaining the baseline compliance share \(b\) and the representation convention
\[
P_{\star}(S=\mathrm{source})
=
P_{\star}(S=\mathrm{target})
=
\frac12.
\]
For each component and for the center, take independent source and target product samples, each of size \(n\). Writing \(P_{\mathrm S}\) and \(P_{\mathrm T}\) for the one-record observed source and target laws induced by the law in question, the center experiment and uniform sign mixture are
\[
P_{\star}^{(n,n)}
=
\left.
\left(P_{\mathrm S}^{\otimes n}\otimes P_{\mathrm T}^{\otimes n}\right)
\right|_{P_{\star}},
\]
\[
M_{\tau}
=
2^{-K_{\mathrm L}}
\sum_{\lambda\in\{-1,1\}^{K_{\mathrm L}}}
\left.
\left(P_{\mathrm S}^{\otimes n}\otimes P_{\mathrm T}^{\otimes n}\right)
\right|_{Q_{\lambda,\tau}}.
\]
\end{algorithmv}

The baseline receipt probability \leanref{sym:r_z^{(a)}}{\ensuremath{r_z^{(b)}}} and receipt--outcome mass function \leanref{sym:R_z^{(a)}}{\ensuremath{R_z^{(b)}(d,y)}} in \cref{obj:synth_15} provide the cells around which the lower family is perturbed. Their receipt contrast equals the baseline compliance share \(b\). The next construction specifies both the observed cells and a compatible potential-outcome law, with an explicit admissibility predicate governing the probability assignments.

\begin{definitionv}{synth_17}[Distances between observed experiments]
For probability measures \(P\) and \(Q\) on the same measurable observation space \((\Omega,\mathcal F)\), define the chi-square divergence by
\[
\chi^2(P,Q)=
\begin{cases}
\displaystyle\int_\Omega
\left(\frac{dP}{dQ}-1\right)^2\,dQ,&P\ll Q,\\
+\infty,&P\not\ll Q.
\end{cases}
\]
The integral is understood in the extended nonnegative reals. Define total-variation distance with the probability-distance normalization by
\[
\operatorname{TV}(P,Q)
=\sup_{A\in\mathcal F}|P(A)-Q(A)|
=\frac12\int_\Omega
\left|\frac{dP}{d\nu}-\frac{dQ}{d\nu}\right|\,d\nu,
\qquad \nu=P+Q.
\]
Thus \(\operatorname{TV}(P,Q)\in[0,1]\). For the lower mixture, take \(P=M_\tau\) and \(Q=P_\star^{(n,n)}\) on the observation space of the independent source and target samples, each of size \(n\).
\end{definitionv}

The admissibility predicate \leanref{sym:\operatorname{Adm}_{n,a}(c_{\star})}{\ensuremath{\operatorname{Adm}_{n,a}(c_{\star})}} in \cref{obj:def:legal-iv-mixture} controls two distinct probability requirements. The height restriction keeps the perturbed observed cells nonnegative and encouragement probabilities positive. The squared-height restriction ensures compatibility of the complier outcome margins: each margin is nonnegative, and the two outcome masses for either receipt state sum to \(b\). Dividing those masses by \(b\) therefore supplies binary conditional outcome distributions for compliers.

The component law \leanref{sym:Q_{\lambda,\tau}(component)}{\ensuremath{Q_{\lambda,\tau}}} in \cref{obj:def:legal-iv-mixture} assigns receipt types to always-takers, never-takers, and compliers, with conditional randomization of encouragement and observed variables obtained by consistency. Sharing the conditional potential-outcome law across populations supplies the transport restrictions. The source--target allocation in that construction represents the full-data law; the sampling experiment remains the product of two independent samples with fixed counts \(n\) and \(n\).

The continuum index \leanref{sym:\tau}{\ensuremath{\tau}} and the scaled perturbation \leanref{sym:v_{\lambda,\tau}}{\ensuremath{v_{\lambda,\tau}}} in \cref{obj:def:legal-iv-mixture} vary the outcome component of the family. Its uniform sign mixture \leanref{sym:M_{\tau}}{\ensuremath{M_{\tau}}} is compared with the equal-size observed product experiment at the center law \leanref{sym:P_{\star}}{\ensuremath{P_{\star}}}. We use the following distance conventions for that comparison.

\begin{lemmav}{lem:legal-iv-mixture-full-elbow}[Legal mixture and effect separation]
Fix the prescribed noncoverage probability \(\alpha\), density envelopes \(c_f,C_f\), and Hölder radius \(L\), satisfying:
\begin{itemize}
\item \textbf{(Noncoverage.)} \(0<\alpha<1\).
\item \textbf{(Lower envelope.)} \(0<c_f<1\).
\item \textbf{(Upper envelope.)} \(1<C_f\).
\item \textbf{(Radius.)} \(1<L\).
\end{itemize}
There exists an amplitude \(c_{\star}\in(0,1)\), chosen uniformly in the sample size and strength floor, with the following properties. For every integer \(n\ge n_0=256\) and every \(0<a\le1/4\), set
\[
K_{\mathrm L}=\lceil n^{4/3}\rceil,
\qquad
\bar a=\max\{a,n^{-1/3}\},
\qquad
h_n=c_{\star}K_{\mathrm L}^{-1/8}.
\]
The amplitude satisfies the construction's admissibility conditions
\[
c_{\star}>0,
\qquad
h_n\le\frac{3}{16},
\qquad
h_n^2\le\bar a\left(\frac14-h_n^2\right).
\]

For every \(\tau\in[-1,1]\) and every cellwise sign vector
\(\lambda\in\{-1,1\}^{K_{\mathrm L}}\), the constructed component
\(Q_{\lambda,\tau}\) belongs to the strength slice
\(\mathcal M_n(a)\) of \cref{obj:def:strength-slice}.
For each \(s\in\{0,1\}\) and every \(x\in[0,1]\), its pointwise conditional-mass averages satisfy
\[
\mathbb E_{Q_{\lambda,\tau}}[C\mid X=x,S=s]=\bar a,
\qquad
\mathbb E_{Q_{\lambda,\tau}}
 [(Y(1)-Y(0))C\mid X=x,S=s]
 =
 -\frac{u_{\lambda}(x)v_{\lambda,\tau}(x)}
 {1/4-u_{\lambda}(x)^2},
\]
where \(C=\mathbf 1\{D(1)>D(0)\}\), \(u_{\lambda}\) is the construction's tiled covariate perturbation, and \(v_{\lambda,\tau}=\tau u_{\lambda}\). These pointwise averages also define versions of the corresponding conditional means under \(Q_{\lambda,\tau}\). Its transported first stage \(\mu\) and target complier effect \(\theta\) satisfy
\[
\mu(Q_{\lambda,\tau})=\bar a,
\qquad
\theta(Q_{\lambda,\tau})=-\frac{\tau J_n}{\bar a},
\]
where the common separation functional \(J_n\) is characterized by
\[
J_n=-\bar a\,\theta(Q_{\lambda,1}).
\]

The unperturbed center \(P_{\star}\) belongs to
\(\mathcal M_n(a)\), and
\[
\mu(P_{\star})=\bar a,
\qquad
\theta(P_{\star})=0,
\qquad
2^{3/4}c_{\star}^2n^{-1/3}
\le J_n\le4c_{\star}^2n^{-1/3}.
\]
For the uniform sign mixture and the center's equal-size observed product experiment,
\[
M_{\tau}
=
2^{-K_{\mathrm L}}
\sum_{\lambda\in\{-1,1\}^{K_{\mathrm L}}}
\left(Q_{\lambda,\tau,\mathrm S}^{\otimes n}
\otimes Q_{\lambda,\tau,\mathrm T}^{\otimes n}\right),
\qquad
P_{\star}^{(n,n)}
=
P_{\star,\mathrm S}^{\otimes n}
\otimes P_{\star,\mathrm T}^{\otimes n},
\]
the subscripts \(\mathrm S,\mathrm T\) denote the observed source-record and target-covariate laws. Uniformly over \(\tau\in[-1,1]\),
\[
1+\chi^2(M_{\tau},P_{\star}^{(n,n)})
\le
\exp\{C_{\mathrm{mix}}c_{\star}^4\},
\qquad
C_{\mathrm{mix}}=\frac{19321}{1800},
\qquad
\operatorname{TV}(M_{\tau},P_{\star}^{(n,n)})
\le\frac{1-\alpha}{2}.
\]
Moreover, \(M_{\tau}\ll P_{\star}^{(n,n)}\), and its likelihood ratio minus one is square-integrable under \(P_{\star}^{(n,n)}\).
\end{lemmav}

\begin{proof}[Proof of \cref{obj:lem:legal-iv-mixture-full-elbow}]
\begin{enumerate}
\item
We first construct a ceiling for amplitudes that gives uniform model membership. Define
\[
c_{\max}
=
\min\left\{
\frac1{100},\,
1-c_f,\,
C_f-1,\,
\frac{L-1}{1+2\pi},\,
\frac{L-3/4}{6\pi}
\right\}.
\]
The hypotheses imply that every entry in this minimum is positive, so
\[
0<c_{\max}\le\frac1{100}<1.
\]
For the membership verification, let
\[
0<c_{\star}\le c_{\max},\qquad
n\ge256,\qquad
0<a\le\frac14,\qquad
\tau\in[-1,1],\qquad
\lambda\in\{-1,1\}^{K_{\mathrm L}}.
\]
Use the construction of \cref{obj:def:legal-iv-mixture}, and write
\[
K_{\mathrm L}=\lceil n^{4/3}\rceil,\qquad
b=\bar a=\max\{a,n^{-1/3}\},\qquad
h=h_n=c_{\star}K_{\mathrm L}^{-1/8},\qquad
u=u_{\lambda},\qquad v=v_{\lambda,\tau}=\tau u.
\]
Since \(n\ge64\),
\[
n^{-1/3}\le64^{-1/3}=\frac14,
\qquad
0<a\le b\le\frac14.
\]
Moreover,
\[
1\le n^{4/3}\le K_{\mathrm L},
\qquad
0<h\le c_{\star}\le\frac1{100},
\]
and monotonicity of the negative power gives
\[
h^2=c_{\star}^2K_{\mathrm L}^{-1/4}
\le c_{\star}^2n^{-1/3}
\le\frac{b}{10000}.
\]
Thus
\[
h^2\le\frac1{10000},
\qquad
\frac14-h^2\ge\frac14-\frac1{10000}\ge\frac1{10000},
\]
and consequently
\[
c_{\star}>0,\qquad
h\le\frac3{16},\qquad
h^2\le\frac{b}{10000}
\le b\left(\frac14-h^2\right).
\]
This proves admissibility throughout the prescribed ranges.

For clarity, the tiled perturbation has the following definition on all of \(\mathbb R\):
\[
u(x)=
\begin{cases}
h\lambda_{\ell}B(K_{\mathrm L}x-\ell),
&x\in\operatorname{cell}(K_{\mathrm L},\ell),
\quad 0\le\ell<K_{\mathrm L},\\
0,&x\notin[0,1],
\end{cases}
\qquad
B(t)=\sin(2\pi t),\quad t\in[0,1].
\]
Here the cells are those of \cref{obj:synth_18}, and \(B\) is the restriction of the bump in \cref{obj:synth_16}. Their disjointness, the endpoint zeros, and the bump normalization give
\[
|v(x)|\le |u(x)|\le h
\quad(x\in\mathbb R),
\qquad
\int_0^1u(x)\,dx
=
\frac{h}{K_{\mathrm L}}
\sum_{\ell=0}^{K_{\mathrm L}-1}
\lambda_{\ell}\int_0^1B(t)\,dt
=0.
\]
In particular, the construction's covariate densities and encouragement probabilities satisfy
\[
f_{\mathrm S}=1,\qquad f_{\mathrm T}=1+u,\qquad
e_z^{(\lambda)}=\frac12+(2z-1)u,\quad z\in\{0,1\},
\]
\[
\int_0^1f_{\mathrm T}(x)\,dx=1,
\qquad
c_f\le1-c_{\star}\le f_{\mathrm T}(x)
\le1+c_{\star}\le C_f,
\]
and
\[
\frac14\le\frac12-h
\le e_z^{(\lambda)}(x)
\le\frac12+h\le\frac34.
\]
The constant source density satisfies the same density envelopes. These are the density and overlap bounds required in
\cref{obj:ass:source-density-bounds,obj:ass:target-density-bounds,obj:ass:propensity-overlap}.
% lean: CausalSmith.Stat.TransportCaceRoughnuisanceLength.exists_membership_amplitude_ceiling
% lean: CausalSmith.Stat.TransportCaceRoughnuisanceLength.lower_admissible_of_small_amplitude
% lean: CausalSmith.Stat.TransportCaceRoughnuisanceLength.legalIVComponent_density_overlap_of_ceiling

\item
We verify the Hölder conditions with the modulus used by the construction. Set
\[
\beta=\frac18.
\]
On a single cell, the sine Lipschitz bound gives
\[
|u(x)-u(x')|
\le 2\pi hK_{\mathrm L}|x-x'|.
\]
The same bound holds across different cells. Indeed, after ordering \(x\le x'\), suppose they belong to cells \(\ell<\ell'\), and define the intervening endpoints by
\[
x_{\ell,+}=\frac{\ell+1}{K_{\mathrm L}},
\qquad
x_{\ell',-}=\frac{\ell'}{K_{\mathrm L}}.
\]
The bump vanishes at these endpoints, and
\[
x\le x_{\ell,+}\le x_{\ell',-}\le x'.
\]
Therefore
\[
\begin{aligned}
|u(x)-u(x')|
&\le |u(x)|+|u(x')|\\
&\le 2\pi hK_{\mathrm L}
\left\{(x_{\ell,+}-x)+(x'-x_{\ell',-})\right\}\\
&\le 2\pi hK_{\mathrm L}|x-x'|.
\end{aligned}
\]
The opposite ordering follows by symmetry.

For \(x,x'\in[0,1]\), define
\[
d_{x,x'}=|x-x'|.
\]
When \(d_{x,x'}=0\), the Hölder bound is immediate. If
\(0<K_{\mathrm L}d_{x,x'}\le1\), then
\[
\begin{aligned}
|u(x)-u(x')|
&\le2\pi c_{\star}K_{\mathrm L}^{-1/8}
(K_{\mathrm L}d_{x,x'})\\
&\le2\pi c_{\star}K_{\mathrm L}^{-1/8}
(K_{\mathrm L}d_{x,x'})^{1/8}
=2\pi c_{\star}d_{x,x'}^{1/8}.
\end{aligned}
\]
If \(K_{\mathrm L}d_{x,x'}>1\), the supremum bound instead gives
\[
|u(x)-u(x')|
\le2h
\le2c_{\star}d_{x,x'}^{1/8}.
\]
Combining the cases yields the uniform modulus
\[
|u(x)-u(x')|
\le c_{\star}(1+2\pi)|x-x'|^{1/8}.
\]
It also proves continuity, including at cell boundaries. Hence
\[
[f_{\mathrm T}]_{\beta},\,[e_1^{(\lambda)}]_{\beta}
\le c_{\star}(1+2\pi)\le L-1<L,
\]
while
\[
\|f_{\mathrm T}\|_{\infty}\le1+c_{\star}\le L,
\qquad
\|e_1^{(\lambda)}\|_{\infty}\le\frac12+c_{\star}\le L.
\]
The source density is constant with supremum norm one and seminorm zero.

The baseline cells of \cref{obj:synth_15} and finite summation in
\cref{obj:def:legal-iv-mixture} give the exact arm means
\[
m_{Dz}(x)=r_z^{(b)}=\frac{1+(2z-1)b}{2},
\qquad
m_{Yz}(x)=
\frac12+\frac{\tau u(x)}{1+2(2z-1)u(x)},
\quad z\in\{0,1\}.
\]
The receipt means are constant and belong to \([3/8,5/8]\). For the outcome means, the denominator satisfies
\[
1+2(2z-1)u(x)\ge1-\frac2{100}=\frac{49}{50}.
\]
Subtracting the two rational expressions gives
\[
m_{Yz}(x)-m_{Yz}(x')
=
\frac{\tau\{u(x)-u(x')\}}
{\{1+2(2z-1)u(x)\}\{1+2(2z-1)u(x')\}}.
\]
Since \(|\tau|\le1\) and \((49/50)^2\ge1/2\),
\[
|m_{Yz}(x)-m_{Yz}(x')|
\le2|u(x)-u(x')|.
\]
Taking \(u(x')=0\) in the same rational comparison gives
\[
\left|m_{Yz}(x)-\frac12\right|
\le2|u(x)|\le\frac1{50},
\]
so the outcome means belong to \([12/25,13/25]\). Their Hölder seminorms satisfy
\[
[m_{Yz}]_{\beta}
\le2c_{\star}(1+2\pi)
\le6\pi c_{\star}
\le L-\frac34<L.
\]
Their supremum norms are at most \(13/25<L\). Thus all density, propensity, and arm-mean Hölder conditions hold, with the range conditions in
\cref{obj:ass:arm-means-holder}. The corresponding center functions are the constants
\[
f_{\mathrm S}=f_{\mathrm T}=1,\qquad
e=\frac12,\qquad
m_{Dz}=r_z^{(b)},\qquad
m_{Yz}=\frac12,
\]
and satisfy the same conditions.
% lean: CausalSmith.Stat.TransportCaceRoughnuisanceLength.tiledPerturbation_holder_modulus
% lean: CausalSmith.Stat.TransportCaceRoughnuisanceLength.rational_arm_mean_abs_sub_le
% lean: CausalSmith.Stat.TransportCaceRoughnuisanceLength.legalIVComponent_arm_means_holder

\item
We next verify the primitive probability masses and calculate their complier margins. By admissibility and \(|v|\le|u|\),
\[
|u(x)v(x)|\le u(x)^2\le h^2
\le b\left(\frac14-h^2\right)
\le b\left(\frac14-u(x)^2\right),
\]
and
\[
\frac14-u(x)^2\ge\frac14-h^2>0.
\]
It follows that
\[
\left|
\frac{u(x)v(x)}
{2\{1/4-u(x)^2\}}
\right|\le\frac b2.
\]
Consequently the complier margins \(M_{dy}(x)\) defined in
\cref{obj:def:legal-iv-mixture} satisfy
\[
M_{dy}(x)\ge0,
\qquad
M_{d0}(x)+M_{d1}(x)=b,
\quad d,y\in\{0,1\}.
\]

For always-takers and never-takers, the relevant masses in that construction are
\[
\frac{1-b}{4}
+\frac{v(x)(2y-1)}{4e_0^{(\lambda)}(x)},
\qquad
\frac{1-b}{4}
+\frac{v(x)(2y-1)}{4e_1^{(\lambda)}(x)},
\quad y\in\{0,1\}.
\]
Using \(|u|,|v|\le3/16\), we obtain
\[
e_z^{(\lambda)}(x)\ge\frac5{16},
\qquad
\left|\frac{v(x)}{4e_z^{(\lambda)}(x)}\right|
\le\frac{3/16}{4(5/16)}=\frac3{20},
\qquad
\frac{1-b}{4}\ge\frac3{16}.
\]
Thus these masses are nonnegative. Summing over \(y\), the perturbations cancel, and each of the two strata has total mass \((1-b)/2\). The independent complier coupling has joint masses
\[
P(C=1,Y(0)=y_0,Y(1)=y_1\mid X=x)
=
\frac{M_{0y_0}(x)M_{1y_1}(x)}{b},
\quad y_0,y_1\in\{0,1\},
\]
whose sum is
\[
\frac{
\left(\sum_{y_0=0}^1M_{0y_0}(x)\right)
\left(\sum_{y_1=0}^1M_{1y_1}(x)\right)
}{b}
=b.
\]
The three principal strata therefore have total mass
\[
\frac{1-b}{2}+b+\frac{1-b}{2}=1.
\]
The remaining primitive cells have mass zero, as specified by the construction.

The receipt assignments \((1,1)\), \((0,0)\), and \((0,1)\) give monotonicity. Conditional independent encouragement, followed by
\[
D=D(Z),\qquad Y=Y(D),
\]
gives receipt consistency, outcome consistency, instrument randomization, and exclusion. All potential outcomes are binary. These verify
\cref{obj:ass:receipt-consistency,obj:ass:outcome-consistency,obj:ass:instrument-randomization,obj:ass:exclusion,obj:ass:monotonicity}.

Finite summation of the same complier masses gives, for either population and every \(x\in[0,1]\),
\[
\mathbb E[C\mid X=x,S=s]=b,
\qquad s\in\{0,1\},
\]
and
\[
\begin{aligned}
\mathbb E[(Y(1)-Y(0))C\mid X=x,S=s]
&=
\sum_{y_0,y_1\in\{0,1\}}
(y_1-y_0)\frac{M_{0y_0}(x)M_{1y_1}(x)}b\\
&=M_{11}(x)-M_{01}(x)
=-\frac{u(x)v(x)}{1/4-u(x)^2}.
\end{aligned}
\]
These are also versions of the conditional means. Indeed, the construction uses the same primitive conditional mass function in both populations. Define their covariate densities by
\[
f_{S=0}(x)=1+u(x),\qquad f_{S=1}(x)=1,
\quad x\in[0,1].
\]
Integrating the finite sums over any Borel set \(E_X\subseteq\mathbb R\) yields
\[
\int_{\{X\in E_X\}} C\,dQ_{\lambda,\tau}(\,\cdot\mid S=s)
=
\int_{E_X\cap[0,1]}b f_{S=s}(x)\,dx
\]
and
\[
\int_{\{X\in E_X\}}(Y(1)-Y(0))C\,
dQ_{\lambda,\tau}(\,\cdot\mid S=s)
=
\int_{E_X\cap[0,1]}
-\frac{u(x)v(x)}{1/4-u(x)^2}f_{S=s}(x)\,dx.
\]
The displayed conditional-mean functions are measurable and bounded. These identities verify their conditional-mean property and both transport requirements.

The arm-mean formulas from the preceding step give
\[
\Delta_D(x)=b,\qquad
\Delta_Y(x)=-\frac{u(x)v(x)}{1/4-u(x)^2}.
\]
Hence
\[
\mu(Q_{\lambda,\tau})
=
\int_0^1(1+u(x))b\,dx=b.
\]
The target complier share is also \(b\), by the conditional-mean identity with \(E_X=\mathbb R\).

Define the separation functional by
\[
J_n
=
h^2\int_0^1
\frac{B(t)^2}{1/4-h^2B(t)^2}\,dt.
\]
The reflection identity \(B(1-t)=-B(t)\) from \cref{obj:synth_16} implies
\[
\int_0^1
\frac{B(t)^3}{1/4-h^2B(t)^2}\,dt
=
-\int_0^1
\frac{B(t)^3}{1/4-h^2B(t)^2}\,dt
=0.
\]
Changing variables separately on each cell therefore gives
\[
\int_0^1\frac{u(x)^2}{1/4-u(x)^2}\,dx=J_n,
\qquad
\int_0^1\frac{u(x)^3}{1/4-u(x)^2}\,dx=0.
\]
Since \(v=\tau u\), the target complier-weighted outcome effect equals
\[
\begin{aligned}
\mathbb E_{Q_{\lambda,\tau}}[(Y(1)-Y(0))C\mid S=0]
&=-\tau\int_0^1
\left\{
\frac{u(x)^2}{1/4-u(x)^2}
+\frac{u(x)^3}{1/4-u(x)^2}
\right\}\,dx\\
&=-\tau J_n.
\end{aligned}
\]
Dividing by the positive target complier share gives
\[
\theta(Q_{\lambda,\tau})=-\frac{\tau J_n}{b}.
\]
In particular,
\[
J_n=-b\,\theta(Q_{\lambda,1}).
\]
For the center, the same calculations with zero perturbations give
\[
\mu(P_{\star})=b,\qquad \theta(P_{\star})=0.
\]
% lean: CausalSmith.Stat.TransportCaceRoughnuisanceLength.primitiveMass_nonneg_of_bounds
% lean: CausalSmith.Stat.TransportCaceRoughnuisanceLength.lowerPointwiseMean_complier_margins
% lean: CausalSmith.Stat.TransportCaceRoughnuisanceLength.legal_mixture_margins

\item
Independent source and target product sampling supplies
\cref{obj:ass:source-iid,obj:ass:target-iid,obj:ass:sample-independence}.
The preceding calculations establish the density identities, support, regularity, causal restrictions, conditional-mean transport, and positive strength required by
\cref{obj:def:model-class}. Thus
\[
Q_{\lambda,\tau}\in\mathcal M_n,\qquad
P_{\star}\in\mathcal M_n.
\]
Together with
\[
a\le b=\mu(Q_{\lambda,\tau})=\mu(P_{\star})\le\frac14,
\]
this proves, by \cref{obj:def:strength-slice},
\[
Q_{\lambda,\tau}\in\mathcal M_n(a),
\qquad
P_{\star}\in\mathcal M_n(a).
\]
We have established these memberships uniformly for every
\(c_{\star}\in(0,c_{\max}]\), every admissible \(n,a\), and every \(\lambda,\tau\).
% lean: CausalSmith.Stat.TransportCaceRoughnuisanceLength.legal_mixture_membership

\item
We now choose one amplitude within this membership ceiling. Since \(\alpha<1\),
\[
\exp\{C_{\mathrm{mix}}\,0^4\}=1<1+(1-\alpha)^2,
\qquad
C_{\mathrm{mix}}=\frac{19321}{1800}.
\]
Continuity at zero supplies \(\delta>0\) such that
\[
0\le c_{\mathrm{test}}<\delta
\quad\Longrightarrow\quad
\exp\{C_{\mathrm{mix}}c_{\mathrm{test}}^4\}
<1+(1-\alpha)^2.
\]
Choose
\[
c_{\star}
=
\min\left\{c_{\max},\,\frac1{100},\,\frac{\delta}{2}\right\}.
\]
Then
\[
0<c_{\star}\le c_{\max}<1,\qquad
c_{\star}\le\frac1{100},\qquad
\exp\{C_{\mathrm{mix}}c_{\star}^4\}\le1+(1-\alpha)^2.
\]
This choice depends on \(\alpha,c_f,C_f,L\) and is uniform in \(n,a,\lambda,\tau\). The admissibility calculation in the first step applies to this amplitude for every prescribed \(n,a\).
% lean: CausalSmith.Stat.TransportCaceRoughnuisanceLength.calibrated_mixture_amplitude

\item
Fix \(n\ge256\) and \(0<a\le1/4\). We bound the common separation functional. The ceiling bounds and \(n^{4/3}\ge1\) give
\[
n^{4/3}\le K_{\mathrm L}
\le n^{4/3}+1\le2n^{4/3}.
\]
Since \(h^2\le1/10000\le1/8\) and \(B(t)^2\le1\),
\[
\frac18\le\frac14-h^2B(t)^2\le\frac14,
\]
so
\[
4B(t)^2
\le\frac{B(t)^2}{1/4-h^2B(t)^2}
\le8B(t)^2.
\]
Using \(\int_0^1B(t)^2\,dt=1/2\), supplied by \cref{obj:synth_16}, we obtain
\[
2
\le\int_0^1\frac{B(t)^2}{1/4-h^2B(t)^2}\,dt
\le4,
\qquad
2h^2\le J_n\le4h^2.
\]
Finally,
\[
(2n^{4/3})^{-1/4}
\le K_{\mathrm L}^{-1/4}\le n^{-1/3},
\]
and therefore
\[
2^{3/4}c_{\star}^2n^{-1/3}
=
2c_{\star}^2(2n^{4/3})^{-1/4}
\le J_n
\le4c_{\star}^2n^{-1/3}.
\]
% lean: CausalSmith.Stat.TransportCaceRoughnuisanceLength.legal_mixture_separation

\item
Fix \(\tau\in[-1,1]\). We calculate the distance of the uniform sign mixture from the center. With respect to Lebesgue measure on \([0,1]\) and counting measure on the three binary source coordinates, the center source density is
\[
p_{\star,\mathrm S}(x,z,d,y)
=\frac12R_z^{(b)}(d,y),
\qquad
x\in[0,1],\quad z,d,y\in\{0,1\}.
\]
The baseline formulas in \cref{obj:synth_15} give
\[
R_z^{(b)}(d,y)\ge\frac{1-b}{4}\ge\frac3{16}>0.
\]
The source cells of \cref{obj:def:legal-iv-mixture} consequently yield the one-record likelihood ratios
\[
\operatorname{LR}_{\mathrm S,\lambda,\tau}(x,z,d,y)
=
1+
u_{\lambda}(x)
\frac{(2z-1)R_z^{(b)}(d,y)+\tau(2y-1)/4}
{R_z^{(b)}(d,y)/2},
\]
\[
\operatorname{LR}_{\mathrm T,\lambda}(x)=1+u_{\lambda}(x).
\]
The first formula is used on the center source support
\([0,1]\times\{0,1\}^3\), and the second on the center target support \([0,1]\); take both ratios to be zero off their respective supports. They are nonnegative, measurable, and bounded, and satisfy
\[
Q_{\lambda,\tau,\mathrm S}
=
\operatorname{LR}_{\mathrm S,\lambda,\tau}\,P_{\star,\mathrm S},
\qquad
Q_{\lambda,\tau,\mathrm T}
=
\operatorname{LR}_{\mathrm T,\lambda}\,P_{\star,\mathrm T}.
\]

Define, for \(b\in(0,1/4]\) and \(\tau\in[-1,1]\),
\[
\gamma(b,\tau)
=
\sum_{z,d,y\in\{0,1\}}
\frac{\{(2z-1)R_z^{(b)}(d,y)+\tau(2y-1)/4\}^2}
{R_z^{(b)}(d,y)/2}.
\]
Expanding the square, its three finite sums are
\[
\sum_{z,d,y}
\frac{\{(2z-1)R_z^{(b)}(d,y)\}^2}{R_z^{(b)}(d,y)/2}
=4,
\]
\[
\sum_{z,d,y}
\frac{(2z-1)R_z^{(b)}(d,y)(2y-1)/4}
{R_z^{(b)}(d,y)/2}
=0,
\]
and
\[
\sum_{z,d,y}
\frac{\{(2y-1)/4\}^2}{R_z^{(b)}(d,y)/2}
=
\frac2{1+b}+\frac2{1-b}
=\frac4{1-b^2}.
\]
The middle sum cancels over \(y\); the last uses the four baseline cells of each mass \((1+b)/4\) and \((1-b)/4\). Thus
\[
\gamma(b,\tau)=4+\frac{4\tau^2}{1-b^2}.
\]
Since \(1-b^2\ge15/16\) and \(\tau^2\le1\),
\[
0\le\gamma(b,\tau),
\qquad
\gamma(b,\tau)+1
\le5+\frac{64}{15}=\frac{139}{15}.
\]

For two sign vectors, define their overlap by
\[
\mathscr O(\lambda,\lambda')
=
\sum_{\ell=0}^{K_{\mathrm L}-1}\lambda_{\ell}\lambda'_{\ell},
\qquad
\lambda,\lambda'\in\{-1,1\}^{K_{\mathrm L}}.
\]
Cellwise integration gives
\[
\begin{aligned}
\int_0^1u_{\lambda}(x)u_{\lambda'}(x)\,dx
&=
\frac{h^2}{K_{\mathrm L}}
\sum_{\ell=0}^{K_{\mathrm L}-1}
\lambda_{\ell}\lambda'_{\ell}\int_0^1B(t)^2\,dt\\
&=\frac{h^2}{2K_{\mathrm L}}\mathscr O(\lambda,\lambda').
\end{aligned}
\]
Expanding the likelihood products, the zero-integral linear terms cancel. Consequently,
\[
\int
\operatorname{LR}_{\mathrm S,\lambda,\tau}
\operatorname{LR}_{\mathrm S,\lambda',\tau}\,dP_{\star,\mathrm S}
=
1+\frac{\gamma(b,\tau)h^2}{2K_{\mathrm L}}
\mathscr O(\lambda,\lambda'),
\]
\[
\int
\operatorname{LR}_{\mathrm T,\lambda}
\operatorname{LR}_{\mathrm T,\lambda'}\,dP_{\star,\mathrm T}
=
1+\frac{h^2}{2K_{\mathrm L}}
\mathscr O(\lambda,\lambda').
\]
Both expressions are nonnegative, being integrals of products of nonnegative likelihoods.

For an observed sample
\[
\omega=
\bigl((x_i,z_i,d_i,y_i)_{i=1}^n,\,
(x_r^{\mathrm T})_{r=1}^n\bigr),
\]
define the product and mixture likelihoods by
\[
\operatorname{LR}_{\lambda,\tau}^{(n,n)}(\omega)
=
\prod_{i=1}^n
\operatorname{LR}_{\mathrm S,\lambda,\tau}(x_i,z_i,d_i,y_i)
\prod_{r=1}^n
\operatorname{LR}_{\mathrm T,\lambda}(x_r^{\mathrm T}),
\]
\[
\operatorname{LR}_{\tau}^{\mathrm{mix}}(\omega)
=
2^{-K_{\mathrm L}}
\sum_{\lambda\in\{-1,1\}^{K_{\mathrm L}}}
\operatorname{LR}_{\lambda,\tau}^{(n,n)}(\omega).
\]
Source--target independence and independent sampling within each channel give
\[
M_{\tau}
=
\operatorname{LR}_{\tau}^{\mathrm{mix}}\,P_{\star}^{(n,n)}.
\]
The bounded one-record likelihoods give integrable product pairs. The finite mixture then has an integrable square, since
\[
\bigl(\operatorname{LR}_{\tau}^{\mathrm{mix}}\bigr)^2
=
2^{-2K_{\mathrm L}}
\sum_{\lambda,\lambda'}
\operatorname{LR}_{\lambda,\tau}^{(n,n)}
\operatorname{LR}_{\lambda',\tau}^{(n,n)}.
\]
Its mean under the center is one.

Let \(\lambda,\lambda'\) be independent uniform sign vectors. Expanding this square and using the two channel inner products gives
\[
\begin{aligned}
1+\chi^2(M_{\tau},P_{\star}^{(n,n)})
&=
\mathbb E_{\lambda,\lambda'}
\left(1+
\frac{\gamma(b,\tau)h^2}{2K_{\mathrm L}}
\mathscr O(\lambda,\lambda')\right)^n
\left(1+
\frac{h^2}{2K_{\mathrm L}}
\mathscr O(\lambda,\lambda')\right)^n\\
&\le
\mathbb E_{\lambda,\lambda'}
\exp\left\{
\frac{nh^2\{\gamma(b,\tau)+1\}}
{2K_{\mathrm L}}\mathscr O(\lambda,\lambda')
\right\}.
\end{aligned}
\]
Here \(1+t\le e^t\) is applied to each nonnegative base, followed by multiplication of their \(n\)-th powers.

For each fixed \(\lambda\), the products \(\lambda_{\ell}\lambda'_{\ell}\) are independent uniform signs. Factoring their exponential moment and using the standard inequality
\[
\cosh(t)\le e^{t^2/2},\qquad t\in\mathbb R,
\]
gives
\[
\begin{aligned}
\mathbb E_{\lambda'}
\exp\left\{
\frac{nh^2\{\gamma(b,\tau)+1\}}
{2K_{\mathrm L}}\mathscr O(\lambda,\lambda')
\right\}
&=
\cosh\left(
\frac{nh^2\{\gamma(b,\tau)+1\}}{2K_{\mathrm L}}
\right)^{K_{\mathrm L}}\\
&\le
\exp\left\{
\frac{n^2h^4\{\gamma(b,\tau)+1\}^2}
{8K_{\mathrm L}}
\right\}.
\end{aligned}
\]
Averaging over \(\lambda\) preserves this bound. Finally,
\[
n^2\le K_{\mathrm L}^{3/2},
\qquad
h^4=c_{\star}^4K_{\mathrm L}^{-1/2},
\]
so
\[
\frac{n^2h^4\{\gamma(b,\tau)+1\}^2}{8K_{\mathrm L}}
=
\frac{n^2}{K_{\mathrm L}^{3/2}}
\frac{\{\gamma(b,\tau)+1\}^2}{8}c_{\star}^4
\le
\frac{(139/15)^2}{8}c_{\star}^4
=
\frac{19321}{1800}c_{\star}^4.
\]
This proves
\[
1+\chi^2(M_{\tau},P_{\star}^{(n,n)})
\le\exp\{C_{\mathrm{mix}}c_{\star}^4\}.
\]

With the distance normalization in \cref{obj:synth_17}, absolute continuity and Cauchy--Schwarz yield
\[
\begin{aligned}
\operatorname{TV}(M_{\tau},P_{\star}^{(n,n)})
&=
\frac12\int
\left|\operatorname{LR}_{\tau}^{\mathrm{mix}}-1\right|
\,dP_{\star}^{(n,n)}\\
&\le
\frac12\left\{
\int
\left(\operatorname{LR}_{\tau}^{\mathrm{mix}}-1\right)^2
\,dP_{\star}^{(n,n)}
\right\}^{1/2}\\
&=\frac12\sqrt{\chi^2(M_{\tau},P_{\star}^{(n,n)})}.
\end{aligned}
\]
The amplitude calibration therefore gives
\[
\chi^2(M_{\tau},P_{\star}^{(n,n)})
\le\exp\{C_{\mathrm{mix}}c_{\star}^4\}-1
\le(1-\alpha)^2,
\]
and, because \(1-\alpha>0\),
\[
\operatorname{TV}(M_{\tau},P_{\star}^{(n,n)})
\le\frac{1-\alpha}{2}.
\]
Every estimate is uniform over \(\tau\in[-1,1]\).
% lean: CausalSmith.Stat.TransportCaceRoughnuisanceLength.sourceLikelihood_pair_integral
% lean: Causalean.Stat.Minimax.Mixture.TwoChannel.one_add_chiSqDiv_uniformMixture_twoChannel_sign_le_exp
% lean: CausalSmith.Stat.TransportCaceRoughnuisanceLength.legal_mixture_distance

\item
The likelihood representation above also establishes
\[
M_{\tau}\ll P_{\star}^{(n,n)}.
\]
The finite sum of integrable product pairs proves integrability of
\((\operatorname{LR}_{\tau}^{\mathrm{mix}})^2\); its mean is one, and hence
\[
\int
\left(\operatorname{LR}_{\tau}^{\mathrm{mix}}-1\right)^2
\,dP_{\star}^{(n,n)}
<\infty.
\]
Thus the likelihood ratio minus one is square-integrable.

Finally, \(n\ge256\) gives \(n>0\), and the chosen amplitude satisfies both admissibility and \(c_{\star}<1\). The triple in
\cref{obj:def:legal-iv-mixture} is therefore defined throughout the stated parameter range. Uniform membership, the center margins, and the separation bounds were established above. For each component and each population, the finite-mass calculations give the asserted pointwise averages, while the integral identities establish their conditional-mean versions. Together with the distance and likelihood conclusions, these prove all the claims.
% lean: CausalSmith.Stat.TransportCaceRoughnuisanceLength.legal_mixture_finite_chiSquare
% lean: CausalSmith.Stat.TransportCaceRoughnuisanceLength.legal_iv_mixture_full_elbow
\end{enumerate}
\end{proof}

The total-variation normalization in \cref{obj:synth_17} measures the largest difference in event probabilities between the mixture and the center experiment. The next result combines this comparison with membership in the causal model and a common effect separation across sign choices.

\begin{definitionv}{synth_16}[Tiled perturbation bump]
Define \(\operatorname{bump}:\mathbb R\to[-1,1]\) by
\[
\operatorname{bump}(t)=
\begin{cases}
\sin(2\pi t),&t\in[0,1],\\
0,&t\notin[0,1].
\end{cases}
\]
Its restriction to \([0,1]\) is the function \(B(t)=\sin(2\pi t)\) used in the tiled perturbation. Its support is \([0,1]\), and its normalization is
\[
\|\operatorname{bump}\|_\infty=1,\qquad
\int_0^1\operatorname{bump}(t)\,dt=0,\qquad
\int_0^1\operatorname{bump}(t)^2\,dt=\frac12.
\]
It satisfies \(\operatorname{bump}(1-t)=-\operatorname{bump}(t)\) for \(t\in[0,1]\) and vanishes at both endpoints. The zero extension is continuous and globally Lipschitz with constant \(2\pi\); its restriction to \([0,1]\) is smooth. In particular, for \(\beta=1/8\),
\[
|\operatorname{bump}(s)-\operatorname{bump}(t)|
\le 2\pi |s-t|^\beta,\qquad s,t\in\mathbb R.
\]
\end{definitionv}

\Cref{obj:lem:legal-iv-mixture-full-elbow} places the entire lower family inside the strength slice while choosing one amplitude uniformly over sample sizes and strength floors. Its separation functional \leanref{sym:J_n}{\ensuremath{J_n}} has order \(n^{-1/3}\), and every sign component at a given continuum index has the same target effect. Uniform sign averaging keeps the observed mixture close to the center's equal-size product experiment. These properties connect effect separation to observational similarity within the causal model.

The actual compliance share is \(\bar a\), as specified in \cref{obj:def:legal-iv-mixture}, for both the components and the center. When \(a\ge n^{-1/3}\), this share equals \(a\), and the effect separation has order \(n^{-1/3}/a\). When \(0<a\le n^{-1/3}\), the share equals \(n^{-1/3}\), and the effect separation has constant order. In both regimes the family belongs to the slice indexed by \(a\). This distinction accounts for the bounded-length regime of the converse while retaining positive compliance and compatible complier outcome margins.

To translate a continuum of effect values into expected confidence-set length, the argument uses the following identity.

\begin{lemmav}{lem:random-set-length-identity}[Expected random-set length]
Let \(C_n\subseteq\Theta\) be a measurable random set under \(P\), where \(\Theta\) is the effect domain and \(\lambda_{\Theta}\) denotes Lebesgue length restricted to that domain. Then
\[
\mathbb E_P\!\left[\lambda_{\Theta}(C_n)\right]
=
\int_{\Theta}P(t\in C_n)\,dt.
\]
\end{lemmav}

The restricted Lebesgue length \leanref{sym:\lambda_{\Theta}}{\ensuremath{\lambda_{\Theta}}} in \cref{obj:lem:random-set-length-identity} aggregates pointwise inclusion probabilities over the effect domain. Its role in the converse is to convert coverage along the lower continuum into expected length at a single center experiment.

Specifically, the continuum supplied by \cref{obj:lem:legal-iv-mixture-full-elbow} spans the effect interval
\[
\left[-\frac{J_n}{\bar a},\frac{J_n}{\bar a}\right].
\]
Global honesty, as defined in \cref{obj:def:honest-intervals}, applies to every component of that continuum. The uniform total-variation comparison makes the center experiment a common reference for these coverage requirements, and \cref{obj:lem:random-set-length-identity} supplies their length interpretation. This is the continuum argument underlying \cref{obj:thm:honest-lower-full-elbow}; all probability comparisons concern the independent source and target product samples of size \(n\) each.

The two compliance regimes therefore give the lower expected-length scale \(\min\{1,n^{-1/3}/a\}\) stated in \cref{obj:thm:honest-lower-full-elbow}. Combining that converse with the attaining interval of \cref{obj:thm:honest-upper-full-elbow} yields the sharp two-sided comparison in \cref{obj:thm:sharp-length-frontier-full-elbow}. Under the prescribed density, smoothness, and coverage conditions, the minimax expected length on each strength slice is thus characterized by the rough estimation scale divided by the strength floor, capped at constant order.


\section{Comparator definitions and verification scope}
\label{sec:verification-appendix}

\subsection{Auxiliary comparators}

The \leanref{S-5}{comparison apparatus} supplements the model in \cref{obj:def:model-class} with supplied geometry, a finite partition, and an explicit continuous construction. These embedded objects provide geometry and calibration comparisons. The supplied-geometry score displays the familiar \(n^{-1/2}\) calibration when the weighting geometry is known, while the continuous witness confirms that the rough ambient geometry is nonempty. We begin by fixing the source and target covariate densities and the source encouragement propensity while retaining the ambient sampling and model restrictions.

\begin{definitionv}{def:fixed-geometry-subclass}[Fixed-geometry subclass $\mathcal M_n^{\mathrm{fix}}(\mathcal G)$]
The fixed-geometry subclass is
\[
\mathcal M_n^{\mathrm{fix}}(\mathcal G)
=
\bigl\{P\in\mathcal M_n:(f_{\mathrm S},f_{\mathrm T},e)=\mathcal G\bigr\},
\]
with equality at sample-size index \(n\). Here \(\mathcal G\) is the supplied triple of source covariate density \(f_{\mathrm S}\), target covariate density \(f_{\mathrm T}\), and source encouragement propensity \(e\). Membership in \(\mathcal M_n\) requires all of the following restrictions at that index:
\begin{itemize}
\item \textbf{(Sampling.)} \cref{obj:ass:source-iid,obj:ass:target-iid,obj:ass:sample-independence}.
\item \textbf{(Covariate densities.)} \cref{obj:ass:source-density-bounds,obj:ass:target-density-bounds,obj:ass:source-density-holder,obj:ass:target-density-holder}.
\item \textbf{(Propensity and arm means.)} \cref{obj:ass:propensity-overlap,obj:ass:propensity-holder,obj:ass:arm-means-holder}.
\item \textbf{(Causal restrictions.)} \cref{obj:ass:receipt-consistency,obj:ass:outcome-consistency,obj:ass:instrument-randomization,obj:ass:exclusion,obj:ass:monotonicity}.
\item \textbf{(Transport and strength.)} \cref{obj:ass:transport-complier-outcome,obj:ass:transport-complier-share,obj:ass:positive-strength}.
\end{itemize}
\end{definitionv}

The supplied geometry \leanref{sym:\mathcal G}{\ensuremath{\mathcal G}} in \cref{obj:def:fixed-geometry-subclass} fixes the three functions entering the transport and encouragement weights. The remaining components of each law continue to satisfy the restrictions defining the ambient model. With this triple available, the following score uses the observed source records directly.

\begin{definitionv}{def:known-geometry-score}[Supplied-geometry score $\widetilde T_A^{\mathcal G}$]
For \(A\in\{Y,D\}\) and every positive integer \(n\), the supplied-geometry score is
\[
\widetilde T_A^{\mathcal G}
=
\frac1n\sum_{i=1}^n
\frac{f_{\mathrm T}^0(X_i)}{f_{\mathrm S}^0(X_i)}
\left\{
\frac{Z_iA_i}{e^0(X_i)}
-
\frac{(1-Z_i)A_i}{1-e^0(X_i)}
\right\}.
\]
Here \((X_i,Z_i,D_i,Y_i)\) are the observed source records, comprising covariates, encouragement, receipt, and outcome, and the supplied geometry is
\[
\mathcal G=(f_{\mathrm S}^0,f_{\mathrm T}^0,e^0),
\]
where \(f_{\mathrm S}^0\) and \(f_{\mathrm T}^0\) are the supplied source and target covariate densities and \(e^0\) is the supplied source encouragement propensity. The statistic is a function of the observed source sample and \(\mathcal G\) alone.
\end{definitionv}

The score \leanref{sym:\widetilde T_A^{\mathcal G}}{\ensuremath{\widetilde T_A^{\mathcal G}}} in \cref{obj:def:known-geometry-score} combines the supplied density ratio with inverse encouragement probabilities. Its outcome and receipt versions enter an affine-score inversion, whose calibration radius is defined next.

\begin{definitionv}{synth_11}[Known geometry calibration radius]
The supplied-geometry calibration radius is defined by
\[
\rho_n^{\mathcal G}
=
\frac{4C_f}{c_f}
\sqrt{\frac{2\log(4/\alpha)}{n}},
\]
for real parameters \(\alpha\), the prescribed noncoverage probability, \(C_f\), the upper density envelope, and \(c_f\), the lower density envelope, and a natural-number sample size \(n\), satisfying:
\begin{itemize}
\item \textbf{(Noncoverage probability.)} \(0<\alpha<1\).
\item \textbf{(Upper envelope.)} \(1<C_f\).
\item \textbf{(Lower envelope.)} \(0<c_f<1\).
\item \textbf{(Sample size.)} \(0<n\).
\end{itemize}
\end{definitionv}

For fixed density envelopes and noncoverage probability, the radius \leanref{sym:\rho_n^{\mathcal G}}{\ensuremath{\rho_n^{\mathcal G}}} in \cref{obj:synth_11} is proportional to \(n^{-1/2}\). The interval definition uses twice this radius as the tolerance for the outcome score minus a candidate effect times the receipt score.

\begin{definitionv}{def:known-geometry-interval}[Supplied-geometry interval $I_n^{\mathcal G}$]
The supplied-geometry interval is
\[
I_n^{\mathcal G}
=
\begin{cases}
\mathcal A_n^{\mathcal G},&\mathcal A_n^{\mathcal G}\ne\varnothing,\\
\{0\},&\mathcal A_n^{\mathcal G}=\varnothing,
\end{cases}
\]
where the affine-score acceptance set is
\[
\mathcal A_n^{\mathcal G}
=
\left\{
v\in\Theta:
\left|\widetilde T_Y^{\mathcal G}
-v\widetilde T_D^{\mathcal G}\right|
\le 2\rho_n^{\mathcal G}
\right\}.
\]
Here \(\Theta\) is the bounded domain of the complier effect, \(\rho_n^{\mathcal G}\) is the calibration radius, and \(\widetilde T_Y^{\mathcal G}\) and \(\widetilde T_D^{\mathcal G}\) are the supplied-geometry outcome and receipt scores.
\end{definitionv}

The supplied-geometry interval \leanref{sym:I_n^{\mathcal G}}{\ensuremath{I_n^{\mathcal G}}} in \cref{obj:def:known-geometry-interval} specifies both score inversion and the singleton convention for an empty acceptance set. Together, \cref{obj:def:known-geometry-score,obj:synth_11,obj:def:known-geometry-interval} define a comparator statistic and its calibration.

A second comparison setting restricts the geometry to a fixed finite interval partition \leanref{sym:\Pi}{\ensuremath{\Pi}}. The following subclass imposes that restriction within the ambient model.

\begin{definitionv}{def:finite-cell-subclass}[Finite-cell subclass $\mathcal M_n^{\mathrm{cell}}(\Pi)$]
The finite-cell subclass \(\mathcal M_n^{\mathrm{cell}}(\Pi)\) consists of the laws \(P\in\mathcal M_n\) whose source covariate density \(f_{\mathrm S}\), target covariate density \(f_{\mathrm T}\), and source encouragement propensity \(e\) are constant on every cell of the fixed finite interval partition \(\Pi\), at sample-size index \(n\). Membership in \(\mathcal M_n\) requires all of the following restrictions at that index:
\begin{itemize}
\item \textbf{(Sampling.)} \cref{obj:ass:source-iid,obj:ass:target-iid,obj:ass:sample-independence}.
\item \textbf{(Covariate densities.)} \cref{obj:ass:source-density-bounds,obj:ass:target-density-bounds,obj:ass:source-density-holder,obj:ass:target-density-holder}.
\item \textbf{(Propensity and arm means.)} \cref{obj:ass:propensity-overlap,obj:ass:propensity-holder,obj:ass:arm-means-holder}.
\item \textbf{(Causal restrictions.)} \cref{obj:ass:receipt-consistency,obj:ass:outcome-consistency,obj:ass:instrument-randomization,obj:ass:exclusion,obj:ass:monotonicity}.
\item \textbf{(Transport and strength.)} \cref{obj:ass:transport-complier-outcome,obj:ass:transport-complier-share,obj:ass:positive-strength}.
\end{itemize}
\end{definitionv}

Membership in \cref{obj:def:finite-cell-subclass} requires cellwise constancy together with the density, smoothness, causal, and transport restrictions listed there. Because the ambient Hölder restrictions impose continuity on the connected interval \([0,1]\), constancy on the cells of an ordinary adjacent interval partition forces the constants to agree at every shared boundary. This continuous finite-cell subclass therefore reduces the three geometry functions to globally constant functions. It is therefore a globally constant geometry specialization embedded in the present continuous model. Its role here is to provide a syntax-level comparison object; the delivered rate conclusions are the frontier results stated in \cref{obj:thm:sharp-length-frontier-full-elbow}.

For a continuous comparison construction, we first define the centered direction that enters the source density, target density, and encouragement propensity.

\begin{definitionv}{synth_13}[Centered covariate witness direction]
The witness direction \(g_{\dagger}\) is defined, for every \(x\in\mathbb R\), by
\[
g_{\dagger}(x)
=
\left|x-\frac12\right|^{1/8}
-
\int_0^1 \left|u-\frac12\right|^{1/8}\,du.
\]
\end{definitionv}

The direction \leanref{sym:g_{\dagger}}{\ensuremath{g_{\dagger}}} in \cref{obj:synth_13} subtracts the average of the fractional-power profile on the covariate domain. The full-data construction below uses this same direction in all three geometry functions and specifies the principal strata, potential variables, and sampling scheme.

\begin{definitionv}{def:continuous-geometry-witness}[Continuous-geometry laws $P_{\sigma}^{\dagger}$]
For each \(\sigma\in\{-1,1\}\), the full-data law \(P_{\sigma}^{\dagger}\) is specified by the following covariate densities, principal-stratum distribution, potential variables, and sampling scheme. Define the perturbation amplitude by
\[
\varepsilon_{\dagger}
=
\frac14\min\{1-c_f,C_f-1,1/4,L-1\},
\]
where \(c_f\) and \(C_f\) are the fixed lower and upper density envelopes and \(L\) is the common Hölder radius. With \(g_{\dagger}\) denoting the covariate perturbation, set
\[
f_{\mathrm S}(x)=1+\sigma\varepsilon_{\dagger}g_{\dagger}(x),
\qquad
f_{\mathrm T}(x)=1-\sigma\varepsilon_{\dagger}g_{\dagger}(x),
\qquad
e(x)=\frac12+\sigma\varepsilon_{\dagger}g_{\dagger}(x).
\]
The law is completed as follows:
\begin{itemize}
\item \textbf{(Principal strata.)} Conditional on \(X\) in each population, the complier, always-taker, and never-taker probabilities are, respectively,
\[
\frac18,\qquad \frac7{16},\qquad \frac7{16}.
\]
\item \textbf{(Potential variables.)} The potential receipts \(D(0)\) and \(D(1)\) are the receipts corresponding to the assigned principal stratum, and the potential outcomes satisfy
\[
Y(0)=Y(1)=\frac12.
\]
\item \textbf{(Source observations.)} Draw encouragement \(Z\) conditionally independently of the potential variables given \(X\), with propensity \(e(X)\), and observe
\[
D=D(Z),\qquad Y=Y(D).
\]
\item \textbf{(Samples.)} Generate independent source and target samples with covariate densities \(f_{\mathrm S}\) and \(f_{\mathrm T}\), respectively.
\end{itemize}
For the combined full-data representation of these two conditional population laws, the population indicator \(S\) has probabilities
\[
P_{\sigma}^{\dagger}(S=\mathrm{source})
=
P_{\sigma}^{\dagger}(S=\mathrm{target})
=
\frac12.
\]
These probabilities specify the representation convention for \(P_{\sigma}^{\dagger}\); the ambient model class retains its stated assumptions.
\end{definitionv}

The continuous-geometry laws \leanref{sym:P_{\sigma}^{\dagger}}{\ensuremath{P_{\sigma}^{\dagger}}} in \cref{obj:def:continuous-geometry-witness} combine opposite source and target density perturbations with fixed principal-stratum probabilities and constant potential outcomes. Their amplitude \leanref{sym:\varepsilon_{\dagger}}{\ensuremath{\varepsilon_{\dagger}}} is specified from the density envelopes, the overlap margin, and the Hölder radius. The equal population probabilities serve the stated full-data representation convention. These definitions establish the availability of explicit comparison objects; the paper's proved conclusions remain the identification, estimation, and expected-length results stated elsewhere.

\subsection{Verification scope}

The supplied theorem statements and proofs are machine-checked in Lean 4 under their stated hypotheses. The verification metadata identifies the sampling and model restrictions referenced in \cref{obj:def:model-class} as assumed inputs and distinguishes the comparator definitions above from proved conclusions. Each checked result retains its own conditions; accompanying theorem-local cited-dependency disclosures, where applicable, identify published conclusions whose source proofs lie outside the formalized proof.


\section{Proofs of the main results}\label{sec:deferred-proofs}

\begin{proof}[Proof of \cref{obj:thm:sharp-length-frontier-full-elbow}]
\begin{enumerate}
\item By \cref{obj:thm:honest-upper-full-elbow}, there is a finite constant \(C_0>0\), depending only on \((\alpha,c_f,C_f,L)\), such that
\[
I_n\in\mathfrak H_n
\qquad\text{for every integer }n>0.
\]
For every integer \(n\ge256\), this same constant satisfies
\[
\mathbb E_P\!\left[\lambda_{\Theta}(I_n)\right]
\le C_0\min\left\{1,\frac{n^{-1/3}}{\mu(P)}\right\}
\qquad\text{for every }P\in\mathcal M_n,
\]
and
\[
L_n(a)\le C_0\min\left\{1,\frac{n^{-1/3}}{a}\right\}
\qquad\text{for every }0<a\le\frac14.
\]
% lean: CausalSmith.Stat.TransportCaceRoughnuisanceLength.honest_upper_full_elbow

\item By \cref{obj:thm:honest-lower-full-elbow}, choose a positive lower comparison constant \(c_{\mathrm{lower}}\), depending only on \((\alpha,c_f,C_f,L)\), for which
\[
c_{\mathrm{lower}}\min\left\{1,\frac{n^{-1/3}}{a}\right\}
\le L_n(a)
\qquad
\text{whenever }n\ge256,\quad 0<a\le\frac14.
\]
% lean: CausalSmith.Stat.TransportCaceRoughnuisanceLength.honest_lower_full_elbow

\item Define the final lower comparison constant by
\[
c_0:=\min\left\{c_{\mathrm{lower}},\frac{C_0}{2}\right\}.
\]
Both entries in this minimum are positive. Consequently,
\[
0<c_0,
\qquad
c_0\le\frac{C_0}{2}<C_0.
\]
Thus \(c_0\) and \(C_0\) have the required dependence and strict ordering.
% lean: CausalSmith.Stat.TransportCaceRoughnuisanceLength.sharp_length_frontier_full_elbow

\item Fix an integer \(n\ge256\) and a real \(a\) with \(0<a\le1/4\). Since \(n>0\) and \(a>0\),
\[
0\le\min\left\{1,\frac{n^{-1/3}}{a}\right\}.
\]
Multiplying \(c_0\le c_{\mathrm{lower}}\) by this nonnegative factor and applying the lower bound gives
\[
c_0\min\left\{1,\frac{n^{-1/3}}{a}\right\}
\le
c_{\mathrm{lower}}\min\left\{1,\frac{n^{-1/3}}{a}\right\}
\le L_n(a).
\]
Together with the upper bound already obtained, this yields
\[
c_0\min\left\{1,\frac{n^{-1/3}}{a}\right\}
\le L_n(a)
\le C_0\min\left\{1,\frac{n^{-1/3}}{a}\right\}.
\]
The honesty assertion and the law-specific expected-length bound follow from \cref{obj:thm:honest-upper-full-elbow} with the same interval sequence and the same \(C_0\).
% lean: CausalSmith.Stat.TransportCaceRoughnuisanceLength.sharp_length_frontier_full_elbow
\end{enumerate}
\end{proof}

\begin{proof}[Proof of \cref{obj:thm:honest-upper-full-elbow}]
\begin{enumerate}
\item
Let \(C_{\mathrm{mse}}\) be the mean-square envelope in
\cref{obj:lem:marked-cubic-mse}. Its displayed expression is finite and strictly positive under the stated conditions. For every positive integer \(n\), define the calibration radius by
\[
r_n:=\sqrt{\frac{2C_{\mathrm{mse}}}{\alpha}}\,n^{-1/3}\ge0.
\]
Thus the acceptance set in \cref{obj:synth_10} is
\[
\mathcal A_n(\omega)
=
\left\{\vartheta\in[-1,1]:
\left|\widehat T_Y(\omega)-\vartheta\widehat T_D(\omega)\right|
\le2r_n\right\},
\]
where the two statistics are those of \cref{obj:def:cubic-estimator}.

The statistics are measurable. Moreover,
\[
\mathcal A_n(\omega)\ne\varnothing
\quad\Longleftrightarrow\quad
|\widehat T_Y(\omega)|
\le|\widehat T_D(\omega)|+2r_n.
\]
Indeed, the values of \(\vartheta\widehat T_D(\omega)\) for
\(\vartheta\in[-1,1]\) form
\([ -|\widehat T_D(\omega)|,|\widehat T_D(\omega)|]\), whose distance from
\(\widehat T_Y(\omega)\) is at most \(2r_n\) precisely under the displayed inequality. Consequently, both the nonemptiness event and the graph of \(\mathcal A_n\) are measurable. For \(n<n_0\), the graph of \(I_n\) is the product of the sample space with \(\Theta=[-1,1]\). For \(n\ge n_0\), its graph combines the graph of \(\mathcal A_n\) on the nonemptiness event with the graph of \(\{0\}\) on its complement. Hence
\[
\{(\omega,\vartheta):\vartheta\in I_n(\omega)\}
\quad\text{is measurable}.
\]
Under every \(P\in\mathcal M_n\), the joint sampling law is a probability law, and the squared estimation errors are integrable for both responses. Below the threshold the errors are constant; above it their integrability follows from the square-integrable centered estimator and conditional-mean error.
% lean: CausalSmith.Stat.TransportCaceRoughnuisanceLength.cubicEstimator_error_sq_integrable
% lean: CausalSmith.Stat.TransportCaceRoughnuisanceLength.scoreInterval_measurableSet_graph

\item
We first check connectedness. Fix a realization \(\omega\), and take
\(\vartheta_1,\vartheta_2\in\mathcal A_n(\omega)\) with
\(\vartheta_1\le\vartheta_2\). For every
\(\vartheta\in[\vartheta_1,\vartheta_2]\), if
\(\widehat T_D(\omega)\ge0\), then
\[
-2r_n
\le \widehat T_Y(\omega)-\vartheta_2\widehat T_D(\omega)
\le \widehat T_Y(\omega)-\vartheta\widehat T_D(\omega)
\le \widehat T_Y(\omega)-\vartheta_1\widehat T_D(\omega)
\le2r_n.
\]
If \(\widehat T_D(\omega)<0\), the corresponding inequalities are
\[
-2r_n
\le \widehat T_Y(\omega)-\vartheta_1\widehat T_D(\omega)
\le \widehat T_Y(\omega)-\vartheta\widehat T_D(\omega)
\le \widehat T_Y(\omega)-\vartheta_2\widehat T_D(\omega)
\le2r_n.
\]
Also \(\vartheta\in[-1,1]\), so \(\vartheta\in\mathcal A_n(\omega)\).
The full-domain and singleton branches in \cref{obj:def:score-interval}
are connected subsets of \(\Theta\) as well. Thus, for every realization,
\[
I_n(\omega)\subseteq\Theta,
\qquad I_n(\omega)\ \text{is connected}.
\]
% lean: CausalSmith.Stat.TransportCaceRoughnuisanceLength.scoreInterval_subset_ordConnected

For coverage, fix \(P\in\mathcal M_n\).
\Cref{obj:lem:transported-cace-identification} supplies
\[
\mu(P)=T_D(P)>0,\qquad
\theta(P)=\frac{T_Y(P)}{\mu(P)}\in[-1,1],
\qquad
T_Y(P)=\theta(P)\mu(P).
\]
If \(n<n_0\), the interval equals \(\Theta\), and therefore
\[
P\{\theta(P)\in I_n\}=1\ge1-\alpha.
\]
Suppose now that \(n\ge n_0\). For each \(A\in\{Y,D\}\), define the error event
\[
\mathcal E_{A,n}
:=
\{\omega:|\widehat T_A(\omega)-T_A(P)|>r_n\}.
\]
The radius is strictly positive and satisfies
\[
r_n^2=\frac{2C_{\mathrm{mse}}}{\alpha}\,n^{-2/3}.
\]
Applying the second-moment tail inequality and
\cref{obj:lem:marked-cubic-mse} gives
\[
P(\mathcal E_{A,n})
\le
\frac{\mathbb E_P[(\widehat T_A-T_A(P))^2]}{r_n^2}
\le
\frac{C_{\mathrm{mse}}n^{-2/3}}
     {(2C_{\mathrm{mse}}/\alpha)n^{-2/3}}
=\frac{\alpha}{2}.
\]
% lean: CausalSmith.Stat.TransportCaceRoughnuisanceLength.cubicEstimator_error_tail_le

Define the simultaneous good event by
\[
\mathcal E_{\mathrm{cov},n}
:=(\mathcal E_{Y,n}\cup\mathcal E_{D,n})^c.
\]
On this event, identification and the triangle inequality yield
\[
\begin{aligned}
|\widehat T_Y-\theta(P)\widehat T_D|
&=
|(\widehat T_Y-T_Y(P))
 +\theta(P)(\mu(P)-\widehat T_D)|\\
&\le|\widehat T_Y-T_Y(P)|
  +|\theta(P)|\,|\widehat T_D-\mu(P)|\\
&\le2r_n.
\end{aligned}
\]
Hence \(\theta(P)\in\mathcal A_n\), so the acceptance set is nonempty and
\(\theta(P)\in I_n\). The union bound therefore gives
\[
P\{\theta(P)\in I_n\}
\ge P(\mathcal E_{\mathrm{cov},n})
\ge1-P(\mathcal E_{Y,n})-P(\mathcal E_{D,n})
\ge1-\alpha.
\]
Together with the measurable graph and connectedness established above,
this proves \(I_n\in\mathfrak H_n\) for every positive integer \(n\), by
\cref{obj:def:honest-intervals}.
% lean: CausalSmith.Stat.TransportCaceRoughnuisanceLength.scoreInterval_coverage_of_square_integrable

\item
Choose
\[
C_0:=
\max\left\{
2,\,
8\sqrt{\frac{2C_{\mathrm{mse}}}{\alpha}}
+8C_{\mathrm{mse}}
\right\}.
\]
This constant is finite, depends only on
\((\alpha,c_f,C_f,L)\), and satisfies
\[
C_0\ge2>0.
\]
% lean: CausalSmith.Stat.TransportCaceRoughnuisanceLength.honest_upper_full_elbow

\item
Fix \(n\ge n_0\) and \(P\in\mathcal M_n\). Define the bad-slope event by
\[
\mathcal E_{\mathrm{slope},n,P}
:=
\{\omega:|\widehat T_D(\omega)-\mu(P)|>\mu(P)/2\}.
\]
Since \(\mu(P)>0\), the second-moment tail inequality and
\cref{obj:lem:marked-cubic-mse} imply
\[
\begin{aligned}
P(\mathcal E_{\mathrm{slope},n,P})
&\le
\frac{\mathbb E_P[(\widehat T_D-\mu(P))^2]}
     {(\mu(P)/2)^2}\\
&\le
\frac{4C_{\mathrm{mse}}n^{-2/3}}{\mu(P)^2}.
\end{aligned}
\]
% lean: CausalSmith.Stat.TransportCaceRoughnuisanceLength.cubicEstimator_bad_slope_probability_le

On the complementary event, the reverse triangle inequality gives
\[
|\widehat T_D|
\ge \mu(P)-|\widehat T_D-\mu(P)|
\ge\frac{\mu(P)}2>0.
\]
For a nonzero slope, the unrestricted affine-score acceptance interval is
\[
\left\{\vartheta\in\mathbb R:
|\widehat T_Y-\vartheta\widehat T_D|\le2r_n\right\}
=
\left[
\frac{\widehat T_Y}{\widehat T_D}
-\frac{2r_n}{|\widehat T_D|},
\,
\frac{\widehat T_Y}{\widehat T_D}
+\frac{2r_n}{|\widehat T_D|}
\right].
\]
Its length is \(4r_n/|\widehat T_D|\). Intersecting with \(\Theta\) bounds
the length by both this quantity and \(2\). Moreover,
\[
\lambda_\Theta(I_n)=\lambda_\Theta(\mathcal A_n):
\]
the two sets agree when \(\mathcal A_n\) is nonempty, and otherwise both
the empty acceptance set and the singleton fallback have length zero.
Consequently, on the complementary event,
\[
\lambda_\Theta(I_n)
\le
\min\left\{2,\frac{4r_n}{|\widehat T_D|}\right\}
\le
\min\left\{2,\frac{8r_n}{\mu(P)}\right\}.
\]
% lean: CausalSmith.Stat.TransportCaceRoughnuisanceLength.restrictedLength_scoreInterval_le_on_good_slope

Define the nonnegative length majorant
\[
\ell_{\mathrm{good},n,P}:=\frac{8r_n}{\mu(P)}\ge0.
\]
Restricted length always lies between zero and the length of \(\Theta\).
Thus, on the bad-slope event it is at most \(2\), while on its complement
the preceding bound applies. For every realization,
\[
0\le\lambda_\Theta(I_n)
\le
\ell_{\mathrm{good},n,P}
+2\mathbf1_{\mathcal E_{\mathrm{slope},n,P}}.
\]
Integrating this inequality gives
\[
\begin{aligned}
\mathbb E_P[\lambda_\Theta(I_n)]
&\le
\ell_{\mathrm{good},n,P}
+2P(\mathcal E_{\mathrm{slope},n,P})\\
&\le
8\sqrt{\frac{2C_{\mathrm{mse}}}{\alpha}}\,
\frac{n^{-1/3}}{\mu(P)}
+
8C_{\mathrm{mse}}\frac{n^{-2/3}}{\mu(P)^2}.
\end{aligned}
\]
There is also the deterministic cap
\[
\mathbb E_P[\lambda_\Theta(I_n)]\le2.
\]
% lean: CausalSmith.Stat.TransportCaceRoughnuisanceLength.scoreInterval_expectedLength_le_of_square_integrable

Set
\[
\xi_{n,P}:=\frac{n^{-1/3}}{\mu(P)}>0,
\qquad
\xi_{n,P}^2=\frac{n^{-2/3}}{\mu(P)^2}.
\]
If \(\xi_{n,P}\le1\), then \(\xi_{n,P}^2\le\xi_{n,P}\), and hence
\[
\begin{aligned}
\mathbb E_P[\lambda_\Theta(I_n)]
&\le
\left(
8\sqrt{\frac{2C_{\mathrm{mse}}}{\alpha}}
+8C_{\mathrm{mse}}
\right)\xi_{n,P}\\
&\le C_0\xi_{n,P}
=C_0\min\{1,\xi_{n,P}\}.
\end{aligned}
\]
If \(\xi_{n,P}>1\), the cap gives
\[
\mathbb E_P[\lambda_\Theta(I_n)]
\le2\le C_0
=C_0\min\{1,\xi_{n,P}\}.
\]
This proves the asserted law-specific expected-length bound.
% lean: CausalSmith.Stat.TransportCaceRoughnuisanceLength.upper_length_two_regime_bound

\item
Fix \(n\ge n_0\) and \(0<a\le1/4\).
The proposed slice bound is nonnegative:
\[
0\le C_0\min\left\{1,\frac{n^{-1/3}}a\right\}.
\]
For every \(P\in\mathcal M_n(a)\),
\cref{obj:def:strength-slice} supplies \(P\in\mathcal M_n\) and
\(a\le\mu(P)\). Since the numerator is nonnegative and \(a>0\),
\[
\frac{n^{-1/3}}{\mu(P)}\le\frac{n^{-1/3}}a.
\]
The law-specific bound therefore yields
\[
\mathbb E_P[\lambda_\Theta(I_n)]
\le
C_0\min\left\{1,\frac{n^{-1/3}}{\mu(P)}\right\}
\le
C_0\min\left\{1,\frac{n^{-1/3}}a\right\}.
\]
If the slice is nonempty, taking its supremum preserves this bound.
For an empty slice, the supremum is taken to be zero, and the same
conclusion follows from the nonnegativity displayed above. In either case,
\[
\sup_{P\in\mathcal M_n(a)}
\mathbb E_P[\lambda_\Theta(I_n)]
\le
C_0\min\left\{1,\frac{n^{-1/3}}a\right\}.
\]
% lean: CausalSmith.Stat.TransportCaceRoughnuisanceLength.honest_upper_full_elbow

\item
To pass to the infimum, observe that every
\(C_n\in\mathfrak H_n\) and every \(P\in\mathcal M_n(a)\) satisfy
\[
0\le\lambda_\Theta(C_n(\omega))\le2,
\qquad
0\le\mathbb E_P[\lambda_\Theta(C_n)]\le2.
\]
These inequalities follow from nonnegativity of Lebesgue length,
restriction to \(\Theta=[-1,1]\), and the probability sampling law.
For a nonempty slice, the expected lengths consequently have a finite
supremum satisfying
\[
0\le
\sup_{P\in\mathcal M_n(a)}
\mathbb E_P[\lambda_\Theta(C_n)]
\le2.
\]
For an empty slice this supremum equals zero. Thus the collection of
suprema over honest procedures is bounded below by zero. It is also
nonempty, since \(I_n\in\mathfrak H_n\). Using \(I_n\) as a competitor in
\cref{obj:def:length-frontier} gives
\[
\begin{aligned}
L_n(a)
&=
\inf_{C_n\in\mathfrak H_n}
\sup_{P\in\mathcal M_n(a)}
\mathbb E_P[\lambda_\Theta(C_n)]\\
&\le
\sup_{P\in\mathcal M_n(a)}
\mathbb E_P[\lambda_\Theta(I_n)]\\
&\le
C_0\min\left\{1,\frac{n^{-1/3}}a\right\}.
\end{aligned}
\]
% lean: CausalSmith.Stat.TransportCaceRoughnuisanceLength.honest_upper_full_elbow
\end{enumerate}
\end{proof}

\begin{proof}[Proof of \cref{obj:thm:honest-lower-full-elbow}]
\begin{enumerate}
\item Choose the amplitude supplied by \cref{obj:lem:legal-iv-mixture-full-elbow}, and define
\[
0<c_{\star}<1,
\qquad
c_0=(1-\alpha)2^{3/4}c_{\star}^2>0.
\]
The amplitude is chosen uniformly in \(n\) and \(a\), so \(c_0\) depends on the prescribed constants \(\alpha,c_f,C_f,L\).
% lean: CausalSmith.Stat.TransportCaceRoughnuisanceLength.legal_iv_mixture_full_elbow

\item Fix \(n\ge256\) and \(0<a\le1/4\). Use the center, components, and observed mixtures of \cref{obj:lem:legal-iv-mixture-full-elbow}, with
\[
K_{\mathrm L}=\lceil n^{4/3}\rceil,
\qquad
\bar a=\max\{a,n^{-1/3}\},
\qquad
r_{\mathrm{sep}}=\frac{J_n}{\bar a}.
\]
Since the negative power is decreasing on the positive half-line,
\[
0<n^{-1/3}\le64^{-1/3}=\frac14,
\qquad
0<\bar a\le\frac14.
\]
The separation bound in that same result gives
\[
0<2^{3/4}c_{\star}^2n^{-1/3}\le J_n,
\qquad
r_{\mathrm{sep}}>0.
\]
For every sign vector \(\lambda\in\{-1,1\}^{K_{\mathrm L}}\), the component at \(\tau=-1\) belongs to \(\mathcal M_n(a)\subseteq\mathcal M_n\). Applying \cref{obj:lem:transported-cace-identification} to this component yields
\[
r_{\mathrm{sep}}
=\theta(Q_{\lambda,-1})
\in[-1,1].
\]
Consequently,
\[
0<r_{\mathrm{sep}}\le1,
\qquad
[-r_{\mathrm{sep}},r_{\mathrm{sep}}]\subseteq\Theta.
\]
The sampling conditions ensure that the center and component observed laws are probability laws.
% lean: CausalSmith.Stat.TransportCaceRoughnuisanceLength.actualStrength_bounds
% lean: CausalSmith.Stat.TransportCaceRoughnuisanceLength.transported_cace_identification

\item The strength floor implements the capped inverse-strength rate. Indeed, if \(a\le n^{-1/3}\), then
\[
\bar a=n^{-1/3},
\qquad
\frac{n^{-1/3}}{\bar a}=1
=\min\left\{1,\frac{n^{-1/3}}a\right\}.
\]
If \(a>n^{-1/3}\), then
\[
\bar a=a,
\qquad
\frac{n^{-1/3}}{\bar a}
=\frac{n^{-1/3}}a
=\min\left\{1,\frac{n^{-1/3}}a\right\}.
\]
Dividing the separation lower bound by \(\bar a>0\) and multiplying by \(1-\alpha>0\) therefore gives
\[
c_0\min\left\{1,\frac{n^{-1/3}}a\right\}
=(1-\alpha)2^{3/4}c_{\star}^2
  \frac{n^{-1/3}}{\bar a}
\le(1-\alpha)r_{\mathrm{sep}}.
\]
% lean: CausalSmith.Stat.TransportCaceRoughnuisanceLength.actualStrength_rate_identity

\item Let \(C_n\) be any procedure satisfying the theorem's hypotheses. For every \(t\in\mathbb R\), define the inclusion event
\[
E_t=\{\omega:t\in C_n(\omega)\}.
\]
These events are measurable by graph measurability. For \(t\in[-r_{\mathrm{sep}},r_{\mathrm{sep}}]\), define
\[
\tau_t=-\frac{t}{r_{\mathrm{sep}}}.
\]
The inequalities \(-r_{\mathrm{sep}}\le t\le r_{\mathrm{sep}}\), together with \(r_{\mathrm{sep}}>0\), give
\[
-1\le\tau_t\le1.
\]
By \cref{obj:lem:legal-iv-mixture-full-elbow}, every component indexed by \(\tau_t\) belongs to \(\mathcal M_n(a)\) and satisfies
\[
\theta(Q_{\lambda,\tau_t})
=-\frac{\tau_tJ_n}{\bar a}
=-\tau_t r_{\mathrm{sep}}
=t.
\]
Uniform honesty thus implies, for every sign vector,
\[
\left(
Q_{\lambda,\tau_t,\mathrm S}^{\otimes n}
\otimes Q_{\lambda,\tau_t,\mathrm T}^{\otimes n}
\right)(E_t)\ge1-\alpha.
\]
Averaging these inequalities over the finite sign family gives
\[
M_{\tau_t}(E_t)
=
2^{-K_{\mathrm L}}
\sum_{\lambda\in\{-1,1\}^{K_{\mathrm L}}}
\left(
Q_{\lambda,\tau_t,\mathrm S}^{\otimes n}
\otimes Q_{\lambda,\tau_t,\mathrm T}^{\otimes n}
\right)(E_t)
\ge1-\alpha.
\]
The normalized mixture is a probability law. The total-variation event inequality and the bound in \cref{obj:lem:legal-iv-mixture-full-elbow} yield
\[
M_{\tau_t}(E_t)-P_{\star}^{(n,n)}(E_t)
\le
\operatorname{TV}(M_{\tau_t},P_{\star}^{(n,n)})
\le\frac{1-\alpha}{2}.
\]
Hence, throughout the symmetric interval,
\[
P_{\star}^{(n,n)}(t\in C_n)
\ge\frac{1-\alpha}{2},
\qquad
t\in[-r_{\mathrm{sep}},r_{\mathrm{sep}}].
\]
% lean: CausalSmith.Stat.TransportCaceRoughnuisanceLength.uniform_data_mixture_coverage
% lean: Causalean.Stat.measureReal_sub_le_tvDist
% lean: CausalSmith.Stat.TransportCaceRoughnuisanceLength.honest_lower_full_elbow

\item Define the center's inclusion-probability function by
\[
p_{\star,n}(t)=P_{\star}^{(n,n)}(t\in C_n),
\qquad t\in\mathbb R.
\]
Graph measurability makes this function measurable, and
\[
0\le p_{\star,n}(t)\le1.
\]
It is therefore integrable over \(\Theta\). Applying \cref{obj:lem:random-set-length-identity} and using the interval inclusion established above gives
\[
\begin{aligned}
(1-\alpha)r_{\mathrm{sep}}
&=\int_{-r_{\mathrm{sep}}}^{r_{\mathrm{sep}}}
       \frac{1-\alpha}{2}\,dt\\
&\le\int_{-r_{\mathrm{sep}}}^{r_{\mathrm{sep}}}
       p_{\star,n}(t)\,dt\\
&\le\int_{\Theta}p_{\star,n}(t)\,dt\\
&=\mathbb E_{P_{\star}}
       \!\left[\lambda_{\Theta}(C_n)\right].
\end{aligned}
\]
The first inequality uses the pointwise inclusion bound; the second uses nonnegativity and
\([-r_{\mathrm{sep}},r_{\mathrm{sep}}]\subseteq\Theta\).
% lean: CausalSmith.Stat.TransportCaceRoughnuisanceLength.expectedLength_lower_of_inclusion

\item The range of expected lengths on the strength slice is bounded above. Pointwise,
\[
0\le\lambda_{\Theta}(C_n(\omega))
\le\operatorname{Leb}(\Theta)=2.
\]
For an integrable length, integration under an observed probability law gives
\[
\mathbb E_P\!\left[\lambda_{\Theta}(C_n)\right]
\le\int 2\,dP^{(n,n)}=2,
\qquad P\in\mathcal M_n(a).
\]
Under the convention assigning zero to a nonintegrable real integral, the complementary case also gives
\[
\mathbb E_P\!\left[\lambda_{\Theta}(C_n)\right]=0\le2.
\]
The measurable bounded lengths considered here are integrable. Since \(P_{\star}\in\mathcal M_n(a)\), the slice is nonempty, and comparison with its supremum now yields
\[
\begin{aligned}
c_0\min\left\{1,\frac{n^{-1/3}}a\right\}
&\le(1-\alpha)r_{\mathrm{sep}}\\
&\le\mathbb E_{P_{\star}}
       \!\left[\lambda_{\Theta}(C_n)\right]\\
&\le\sup_{P\in\mathcal M_n(a)}
       \mathbb E_P\!\left[\lambda_{\Theta}(C_n)\right]
\le2.
\end{aligned}
\]
This proves the asserted procedure-wise lower bound.
% lean: CausalSmith.Stat.TransportCaceRoughnuisanceLength.expectedLength_le_two
% lean: CausalSmith.Stat.TransportCaceRoughnuisanceLength.honest_lower_full_elbow

\item The honest interval class in \cref{obj:def:honest-intervals} is nonempty. Define the constant procedure by
\[
C_{n,\Theta}(\omega)=\Theta
\qquad\text{for every sample realization }\omega.
\]
It has a measurable graph and connected values. By \cref{obj:lem:transported-cace-identification},
\[
\theta(P)\in\Theta,
\qquad
P^{(n,n)}\bigl(\theta(P)\in C_{n,\Theta}\bigr)
=1\ge1-\alpha,
\qquad P\in\mathcal M_n.
\]
Thus
\[
C_{n,\Theta}\in\mathfrak H_n,
\qquad
\mathfrak H_n\ne\varnothing.
\]
Every member of this class satisfies the procedure-wise bound already established. Taking the infimum prescribed by \cref{obj:def:length-frontier} gives
\[
L_n(a)
=
\inf_{C_n\in\mathfrak H_n}
\sup_{P\in\mathcal M_n(a)}
\mathbb E_P\!\left[\lambda_{\Theta}(C_n)\right]
\ge
c_0\min\left\{1,\frac{n^{-1/3}}a\right\}.
\]
% lean: CausalSmith.Stat.TransportCaceRoughnuisanceLength.parameterSpace_honest
% lean: CausalSmith.Stat.TransportCaceRoughnuisanceLength.honest_lower_full_elbow
\end{enumerate}
\end{proof}

\bibliographystyle{plainnat}

\bibliography{references}

\end{document}
